% ======================================================================
%  PAPER B DRAFT v5 -- GKSL consistency and finite-level clarification
%  Consistency patch: fresh ancillas, pointer-neutral Hamiltonian,
%  quantitative level adiabaticity, qutrit dilation, active/passive sign, k/N.
%  Dynamical SU(2) Chern--Simons Boundary Theory on the Cosmological
%  Apparent Horizon: Screening, Horizon Memory, and Calibration Holonomy
% ======================================================================
\documentclass[aps,prd,twocolumn,10pt,superscriptaddress,nofootinbib]{revtex4-2}

\usepackage{amsmath,amssymb,amsthm}
\usepackage{bm}
\usepackage{booktabs}
\usepackage{microtype}
\usepackage[colorlinks=true,linkcolor=blue,citecolor=blue,urlcolor=blue]{hyperref}
\hypersetup{
  pdftitle={Dynamical SU(2) Chern--Simons Boundary Theory on the Cosmological Apparent Horizon: Screening, Horizon Memory, and Calibration Holonomy},
  pdfauthor={Bernard Sayegh}
}
\tolerance=2000
\emergencystretch=20pt

\newtheorem{proposition}{Proposition}
\newtheorem{corollary}[proposition]{Corollary}
\newtheorem{lemma}[proposition]{Lemma}

\begin{document}

\title{Dynamical SU(2) Chern--Simons Boundary Theory on the Cosmological Apparent Horizon:\\ Screening, Horizon Memory, and Calibration Holonomy}

\author{Bernard Sayegh}
\affiliation{Independent Researcher, Sydney, Australia}
\email[Contact author: ]{b.sayegh@bernardsayegh.com}

\date{\today}

\begin{abstract}
We develop a microscopic boundary theory for the two late-time responses of
the cosmological apparent horizon.  The processed branch is described by an
SU(2) Chern--Simons theory with slice-dependent level $k\propto H^{-2}$.
Its puncture channels separate into entropy-producing $j=1/2$ emergence
channels and flux-mediating $j=1$ interaction channels.  The boundary state
sum, its independent KMS formulation, and the source-reduced path integral
give the channel allocation $I_{\rm eff}=(1-q)/2$, the screening polynomial
$x(q)=(1-q)(2-q)/6$, and the Pendry-normalised coupling $\beta=1/12$, with
nonlinear corrections suppressed by $\mathcal O(1/k)\sim10^{-122}$.

The over-capacity branch is treated separately.  The rejected fraction
$\Delta=\Theta(\mathcal S-1)(1-\mathcal S^{-1})$ defines an additive
capacity-cycle history $X_H=\int\Delta\,d\lambda_H$, where
$d\lambda_H=Hdt=d\ln a$ on FLRW.  Full counting statistics of the
source-reduced processed/excess instrument gives
$d\langle\chi\rangle/d\lambda_H=\Delta$ with variance suppressed by $1/N$;
hence $\chi_0=X_H$ at the semiclassical saddle.  The shift-symmetric
three-form carrier is the covariant large-$N$ influence functional of this
counting record rather than an independent write ansatz.

We then construct an explicit local matter--boundary comparator EFT.  A
Ramsey-readable atomic phase is compared with a memory-neutral
processed-sector boundary phase, while a Glauber radiative amplitude is
compared with a geometric radiometric reference.  The two records attach to
the conjugate endpoints of one gauge-neutral current bilocal
$J^+W_\gamma J^-$.  Each affine endpoint has exact weight one, the normalized
conditional matrix element removes the overall portal overlap, and the
single defect therefore fixes $g_T=g_R=g_{\rm com}=1$ in the minimal
comparator sector.  A global hierarchical representation of the distance
ladder contains one serial defect crossing; multiple anchors and
supernovae are parallel measurements of the same two latent zero points and
do not reapply the gate.  Retarded passive orientation fixes the signed
native shifts
$\Delta a_B=X_H/\ln10$ and
$\Delta M_B=5X_H/\ln10$, yielding the parameter-independent relation
$\Delta M_B=5\Delta a_B$ and
$H_{0,\rm ladder}/H_{0,\rm phys}=e^{2X_H}$.

The convention-independent observable is nonclosure of a redundant
calibration loop.  We formulate it as a register-valued graph holonomy and
prove that node-zero-point redefinitions telescope from every closed loop.
We further give a source-reduced Hamiltonian dilation of the write
instrument, a concrete caesium-clock and SI-radiometric realization of the
matter endpoints, a symmetry-complete counterterm basis, an executable
native $(M_B,a_B)$ likelihood, estimator-level reductions for sirens,
chronometers, time-delay lenses, and megamasers, and a pre-registered
bright-siren--SN closure experiment.  These results complete the minimal
low-energy construction.  The remaining frontier is empirical execution
and the question whether the exact Standard Model plus non-perturbative
time-dependent CS/WZW theory uniquely selects this minimal portal.
\end{abstract}

\maketitle



\section{Introduction}
\label{sec:intro}

The cosmological apparent horizon is both a quasi-local gravitational
boundary and a dynamical thermodynamic screen.  A companion cosmology
paper~\cite{Sayegh2026} uses its processed throughput to define a fixed
late-time dark-sector interaction and studies the resulting suppression of
structure growth.  The present paper asks a more microscopic question:
what boundary theory generates the interaction, and what becomes of the
fraction of the horizon drive that exceeds the processed capacity?

The answer developed here has two branches on one boundary,
\begin{equation}
\dot Q_{\rm drive}
=\dot Q_{\rm processed}+\dot Q_{\rm excess}.
\label{eq:two_branch_split_intro}
\end{equation}
The processed branch is continuously assimilated into the coarse-grained
cosmological dynamics.  Its SU(2) Chern--Simons description produces three
fixed microscopic inputs for the companion cosmology:
\begin{equation}
I_{\rm eff}=\frac{1-q}{2},\qquad
x(q)=\frac{(1-q)(2-q)}{6},\qquad
\beta=\frac{1}{12}.
\label{eq:three_inputs_intro}
\end{equation}
The over-capacity branch is different.  It is not introduced as a new bulk
energy density, a photon optical depth, or a homogeneous lapse perturbation.
Instead, the rejected fraction writes an additive boundary history, which
is stored in a protected zero mode and read locally when two inequivalent
physical standards are interfaced.

The central conceptual distinction is therefore
\begin{align}
\dot Q_{\rm processed}
&\longrightarrow \{Q,I_{\rm eff},x(q),\beta\},\notag\\
\dot Q_{\rm excess}
&\longrightarrow \{X_H,\chi,\text{endpoint gate}\}.
\label{eq:two_branch_logic}
\end{align}
The first line controls physical growth.  The second line controls a
relative calibration-transfer observable.  Separating them prevents the
uncertainty of the operational recipient from contaminating the growth
calculation.

This paper makes seven contributions.
\begin{enumerate}
\item It extends the isolated-horizon SU(2) Chern--Simons construction to a
slice-wise dynamical cosmological level $k(t)$ and derives the processed
screening sector.
\item It absorbs the source-reduction calculation: the homogeneous drive
couples to the boundary monopole, the Gauss constraint freezes the
longitudinal adjoint current, and two transverse Ohmic channels remain.
\item It derives the over-capacity write law from a processed/excess quantum
instrument and its full counting statistics.  The geometric-capacity clock
sets the channel slots, the homogeneous monopole samples the intensive
rejected fraction, and the deterministic saddle gives $\chi_0=X_H$ with
$1/N$ fluctuations.
\item It derives the exact endpoint-response theorem and its common-gate
corollary, and evaluates the normalized matrix element of the gauge-neutral
SU(2)$_k$ current bilocal of total weight two.
\item It constructs explicit temporal and radiative matter portals using a
local phase ratio and a local radiative-amplitude ratio, and shows that they
are the two endpoints of one normalized physical defect rather than two
independent couplings.
\item It proves one engagement in the global hierarchical distance ladder:
parallel calibrators and Hubble-flow objects constrain shared zero points,
whereas only the single serial preparation of that state carries the gate.
A retarded passive portal fixes the endpoint orientation and the sign of the
native map.
\item It formulates calibration nonclosure as a register-valued graph
holonomy, removing arbitrary node zero points and forbidding response units
from being manufactured by subdividing a data-analysis pipeline.
\end{enumerate}

The epistemic status is deliberately graded.  The SU(2) representation
structure, screening polynomial, current weights, loop invariance,
capacity-cycle write law, and the normalized comparator matrix elements are
derivations within the stated source-reduced large-$N$ boundary model.  The
paper now supplies an explicit Stinespring/GKSL Hamiltonian dilation, a
concrete metrological realization of both endpoints, a counterterm
classification, a native likelihood algorithm, estimator-level register
reductions, and a closure-experiment protocol.  This is a complete minimal
open-system EFT and statistical implementation.  It is not a uniqueness
theorem for the non-perturbative Standard Model plus quantum-gravity
Hamiltonian, nor does it replace the empirical execution of the released
likelihood and proposed closure experiment.


\section{Macroscopic Framework}
\label{sec:macro}

This section provides a self-contained summary of the results from
the companion cosmology paper~\cite{Sayegh2026} that define the screening problem addressed
in this paper.

\subsection{The holographic interaction kernel}
\label{sec:kernel_recap}

In a flat FLRW universe with apparent horizon $r_A = c/H$, the
vacuum energy swept into the horizon per unit time is
\begin{equation}
\dot{M}_\Lambda
  = \rho_\Lambda \cdot 4\pi r_A^2 \cdot \dot{r}_A
  = \frac{9}{4}\,\Omega_\Lambda(1-\Omega_\Lambda)\,\frac{c^3}{G}\,.
\label{eq:sweep}
\end{equation}
The Hubble parameter cancels exactly, leaving an
$H_0$-independent kinematic relation.  The sweep rate approaches
the Schwarzschild processing limit $c^3/2G$ within $2\%$ for
$\Omega_\Lambda \simeq 0.68$; saturating the bound,
$\dot{M}_\Lambda = c^3/2G$, yields the quadratic
$9\Omega_\Lambda(1-\Omega_\Lambda)/2 = 1$ with roots
\begin{equation}
\Omega_\Lambda = \frac{1}{3}\,,\qquad \Omega_\Lambda = \frac{2}{3}\,,
\label{eq:roots}
\end{equation}
corresponding precisely to the onset of cosmic acceleration and
the near-present epoch.  Between these roots, the sweep rate
\emph{exceeds} the Schwarzschild benchmark ($\mathcal{S} > 1$);
this super-Schwarzschild window spans from $z \approx 0.63$
to $z \approx 0.02$.

Within this window, the excess fraction
$f_{\rm ss} = \max(0,\,1 - 1/\mathcal{S})$ is nonzero.  The holographic
interaction kernel itself is active at all epochs:
\begin{equation}
Q = \tfrac{9}{2}\,\beta\,H\rho_{\rm crit}\,\Omega_\Lambda\,\Omega_m\,,
\label{eq:Q}
\end{equation}
vanishing only at $\Omega_\Lambda = 0$ (matter domination)
and $\Omega_\Lambda = 1$ (de~Sitter).  (The complete interaction
of the companion cosmology paper carries the additional thermodynamic modulation,
$Q = \tfrac{9}{2}\beta\,I_{\rm eff}\,H\rho_{\rm crit}\Omega_\Lambda\Omega_m$
with $I_{\rm eff} = (1-q)/2$, whose microscopic origin is the
subject of this paper.)  The coupling is fixed by
\begin{equation}
\beta = \frac{c^5/24G}{c^5/2G} = \frac{1}{12}\,.
\label{eq:beta}
\end{equation}
The numerator is the Pendry-limited quantum channel throughput of
the Bekenstein--Hawking~\cite{Bekenstein1973,Hawking1975} pixelated horizon
($N = A/4\ell_P^2$ channels, each at the Pendry limit for graviton
modes at the de~Sitter temperature~\cite{GibbonsHawking1977}); the denominator is the
classical geometric capacity of the holographic
screen~\cite{Pendry1983}.  The value $\beta = 1/12$ numerically coincides with the Nariai
stability margin ($1/3$) times the Bekenstein--Hawking
coarse-graining normalisation ($1/4$); we record this as an
observation, not an independent derivation.

\subsection{Routes to \texorpdfstring{$I_{\rm eff} = (1-q)/2$}{I\_eff = (1-q)/2}}
\label{sec:three_routes}

The interaction is modulated by a thermodynamic efficiency factor
$I_{\rm eff}$ that suppresses coupling at early times and saturates
at de~Sitter equilibrium.  Two independent macroscopic derivations
and a kinematic restatement converge on the same functional form.

\paragraph{(i) Onsager linear response.}
The Kodama--Hayward surface gravity of the apparent
horizon~\cite{Hayward1998} is
\begin{equation}
\kappa_{\rm KH} = -r_A\!\left(H^2 + \frac{\dot{H}}{2}\right)
  = -\frac{cH(1-q)}{2}\,,
\label{eq:kappa_KH}
\end{equation}
using $\dot{H} = -(1+q)H^2$.  At de~Sitter ($q = -1$), this
reaches its maximum $|\kappa_{\rm dS}| = cH$.  In Onsager's
framework~\cite{Onsager1931}, the thermodynamic flux is proportional to the first
power of the driving force; the natural proximity measure is
\begin{equation}
I_{\rm eff} = \frac{|\kappa_{\rm KH}|}{|\kappa_{\rm dS}|} = \frac{1-q}{2}\,.
\label{eq:Ieff_onsager}
\end{equation}

\paragraph{(ii) Padmanabhan's holographic equipartition.}
In Padmanabhan's framework~\cite{Padmanabhan2010,Padmanabhan2012}, the accelerated
expansion is driven by the displacement between the surface degrees
of freedom $N_{\rm sur} = A/\ell_P^2$ and the bulk degrees of
freedom $N_{\rm bulk}$.  The volume emergence rate (which reaches its maximum of unity
at radiation domination and vanishes at de~Sitter), is
\begin{equation}
\Gamma_{\rm emerge} = \frac{N_{\rm sur} - N_{\rm bulk}}{2N_{\rm sur}}
  = \frac{1+q}{2}\,.
\label{eq:Gamma_emerge}
\end{equation}
If the horizon possesses a single finite processing
bandwidth---a direct consequence of the holographic principle---then
emergence and the dark-sector interaction draw from a shared budget.
By capacity conservation,
\begin{equation}
I_{\rm eff} = 1 - \Gamma_{\rm emerge} = \frac{1-q}{2}\,.
\label{eq:Ieff_padmanabhan}
\end{equation}

\paragraph{(iii) Kinematic restatement.}
The result also follows directly from the definition of $q$:
$|\kappa_{\rm KH}|/|\kappa_{\rm dS}| = (1-q)/2$ is an algebraic identity
for the apparent horizon in FLRW spacetime, requiring only the
Friedmann equations and $\dot{H} = -(1+q)H^2$.  This is
the mathematical content of route~(i) stripped of its
thermodynamic interpretation.

Routes (i) and (ii) use different physical
inputs---intensive thermodynamic ratio and extensive holographic
capacity---yet yield the same $(1-q)/2$, with route~(iii)
confirming this as a kinematic identity.
The convergence is non-trivial: one could construct functions of $q$
that respect the de~Sitter and radiation endpoints but differ
elsewhere (e.g.\ the quadratic $(1-q)^2/4$, which gives a
matter-era efficiency of $1/16$ rather than $1/4$; the form
$(1-q^2)/4$ fails outright at the de~Sitter endpoint).  The Lagrangian alternative over-suppresses
clustering at fixed cosmology ($\sigma_8 = 0.678$; see below),
while the MCMC analysis of the companion cosmology paper~\cite{Sayegh2026}, floating the
coupling with the linear form, recovers $\beta = 0.088 \pm 0.014$ in the full-stack diagnostic,
consistent with $1/12$; this is an amplitude check rather than an
independent derivation of the channel structure.

All three routes (two macroscopic derivations and the kinematic restatement), however, are \emph{macroscopic}.  They
establish what $I_{\rm eff}$ is, but not why the microscopic
degrees of freedom of the horizon organise to produce this
specific modulation.

\subsection{The screening problem}
\label{sec:screening_problem}

A quintessence realisation of the interaction exists: conformally
coupled quintessence with potential slope
$\lambda = \sqrt{9/2} \approx 2.12$ reproduces the kernel~\eqref{eq:Q}
exactly, with the product of couplings fixed:
$\beta_c\lambda = 9\beta/2 = 3/8$~\cite{Sayegh2026}.  However,
the minimal scalar action generates a natural modulation
\begin{equation}
I_{\rm eff}^{\rm Lag} = \frac{3}{2-q}
  = \frac{2}{1+\Omega_\Lambda}\,,
\label{eq:Ieff_lag}
\end{equation}
obtained from the quasi-static Klein--Gordon equation
($\phi'' = 0$ in $N = \ln a$), where $\phi' \propto 1/(2-q)$.
This modulation \emph{increases} during matter domination
($I_{\rm eff}^{\rm Lag} \approx 2$ at $q = 1/2$), the opposite
behaviour to the thermodynamic form ($I_{\rm eff} = 1/4$ at
$q = 1/2$).  A standalone Boltzmann-code comparison at fixed
cosmology gives $\sigma_8 = 0.678$ for the Lagrangian modulation
versus $0.756$ for the thermodynamic one
($\Lambda$CDM: $0.812$)---the Lagrangian $\sigma_8$ suppression
relative to $\Lambda$CDM is ${\sim}\,2.4$ times larger than the
thermodynamic one (i.e.\
$(0.812-0.678)/(0.812-0.756) \approx 2.4$).

The gap between the two forms defines a screening function $x(q)$:
\begin{equation}
x(q) \cdot I_{\rm eff}^{\rm Lag} = I_{\rm eff}^{\rm thermo}
\quad\Longrightarrow\quad
x(q) = \frac{(1-q)(2-q)}{6}\,.
\label{eq:x_def}
\end{equation}
This polynomial has suggestive structure: a product of two linear
factors in $q$, normalised so that $x = 1$ at de~Sitter ($q = -1$)
and $x = 0$ at radiation domination ($q = 1$).  the companion cosmology paper identifies
$x(q)$ phenomenologically as the correction needed to reconcile the
Lagrangian with the thermodynamic modulation and notes the connection
to a field-dependent coupling $\beta(\phi)$.  But the factorised
form $(1-q)(2-q)/6$ demands a microscopic explanation.  Substituting
$\langle j \rangle = (1-q)/2$ and noting $2 - q = 2\langle j \rangle + 1$:
\begin{equation}
x = \frac{\langle j\rangle\,(2\langle j\rangle + 1)}{3}\,,
\label{eq:x_spin_preview}
\end{equation}
which is the ratio of the throughput of an effective
spin-$\langle j\rangle$ channel (evaluated at the saddle point of
the boundary partition function) to that of the adjoint
($j = 1$, $d_1 = 3$); $\langle j \rangle$ is a continuous
saddle-point value, not a discrete irrep label.
The remainder of this paper derives this expression from the
SU(2) Chern-Simons theory on the apparent horizon.


% ======================================================================
%  SECTION 3: SU(2) CHERN-SIMONS THEORY ON THE APPARENT HORIZON
% ======================================================================
\section{SU(2) Chern-Simons Theory on the Apparent Horizon}
\label{sec:su2}

The ABCK framework~\cite{Ashtekar1998,Ashtekar2000} provides the
standard loop quantum gravity (LQG) description of black hole
horizons.  Bulk spin-network edges pierce the horizon at discrete
points---\emph{punctures}---each carrying an SU(2) spin label $j$.
The boundary degrees of freedom are governed by an SU(2)
Chern-Simons theory whose level $k$ is set by the horizon area.
For an isolated (stationary) black hole, $k$ is fixed and the
entropy counting reproduces the Bekenstein--Hawking area law.

The cosmological apparent horizon is not an isolated horizon.  Its
area changes with the expansion history, so the Chern-Simons level
becomes dynamical.  This section establishes three consequences of
this distinction: (i)~the level evolves as $dk/d\ln a = 2(1+q)k$;
(ii)~the growing state space partitions into two channel populations
with rigid SU(2) representation assignments; and (iii)~the mean
processing spin of the interaction sector tracks the deceleration
parameter, $\langle j \rangle = (1-q)/2$.


\subsection{The dynamical Chern-Simons level}
\label{sec:dynamical_k}

In the ABCK framework, the Chern-Simons level is determined by the
horizon area~\cite{Ashtekar1998}:
\begin{equation}
k = \frac{A}{4\pi\gamma\,\ell_P^2}\,,
\label{eq:CS_level}
\end{equation}
where $\gamma$ is the Barbero--Immirzi parameter~\cite{Barbero1995,Immirzi1997} and
$\ell_P = \sqrt{\hbar G/c^3}$ is the Planck length.  The level $k$
counts the effective number of quantum-gravitational degrees of
freedom on the boundary; it sets the dimension of the Hilbert space
of the Chern-Simons theory and controls the representation content
available at each puncture.

For the cosmological apparent horizon in flat FLRW spacetime, the
horizon radius is $r_A = c/H$ and the area is $A = 4\pi c^2/H^2$,
giving
\begin{equation}
k = \frac{c^2}{\gamma\,\ell_P^2\,H^2}\,.
\label{eq:k_FLRW}
\end{equation}
Unlike the isolated horizon of a stationary black hole, $H$ is
time-dependent.  The extension of the ABCK symplectic structure to
the dynamical apparent horizon---a marginally trapped tube with
evolving area~\cite{AshtekarKrishnan2003}---is derived from the
Holst action~\cite{Holst1996} variational principle in
Appendix~\ref{sec:dynamical_k_derivation}.
Differentiating~\eqref{eq:k_FLRW} and using the identity
$\dot{H} = -(1+q)H^2$:
\begin{equation}
\frac{dk}{d\ln a} = -\frac{2\dot{H}}{H^2}\,k = 2(1+q)\,k\,.
\label{eq:dk_dlna}
\end{equation}
This is the central kinematic statement: the Chern-Simons level
grows during deceleration ($q > -1$), with a rate proportional to
$(1+q)$, and freezes only at de~Sitter equilibrium ($q = -1$).
The quantum-gravitational state space of the horizon is not static;
it evolves with the expansion history.

The present-day value is enormous:
$k_0 \sim c^2/(\gamma\,\ell_P^2\,H_0^2) \sim 10^{122}/\gamma$,
placing the boundary theory deep in the semiclassical regime.  This
will be important in Section~\ref{sec:kms}: the conformal dimension
gap between representations scales as
$\Delta h \sim j(j+1)/(k+2) \to 0$ for $k \gg 1$, ensuring that
the fluctuation-dissipation theorem applies in the classical limit.


\subsection{Two classes of horizon channels}
\label{sec:two_channels}

The growth of $k$ corresponds physically to the creation of new
punctures on the horizon (in the entropy-maximising regime where
the dominant spin is $j = 1/2$, area growth at fixed $j$ requires
additional punctures).  Each spin-$j$ puncture contributes an
area quantum
\begin{equation}
a_j = 8\pi\gamma\,\ell_P^2\,\sqrt{j(j+1)}\,,
\label{eq:area_quantum}
\end{equation}
and the total area is $A = \sum_p a_{j_p}$.  The rate of area
growth follows from~\eqref{eq:dk_dlna}:
\begin{equation}
\frac{dA}{d\ln a} = \frac{8\pi c^2(1+q)}{H^2}\,,
\label{eq:dA_dlna}
\end{equation}
so the rate of new puncture creation scales as $2(1+q)N$, where
$N = A/(4\ell_P^2)$ is the coarse-grained surface DOF count.  This is the
extensive, microscopic counterpart of Padmanabhan's capacity
argument (Section~\ref{sec:three_routes}): volume emergence
consumes a fraction $\Gamma_{\rm emerge} = (1+q)/2$ of the
horizon's total processing bandwidth, leaving
$I_{\rm eff} = (1-q)/2$ for interaction processing.

The horizon's channels are thus partitioned into two populations,
distinguished by their physical role and---crucially---by their
SU(2) representation content.  We now show that the representation
assignments are forced by the structure of the boundary theory,
not freely chosen.


\subsubsection{Emergence channels: \texorpdfstring{$j = 1/2$}{j = 1/2} (fundamental)}
\label{sec:emergence_channels}

In the ABCK framework, the boundary condition relating the bulk
Ashtekar connection to the Chern-Simons boundary field is
$F = -(\pi\gamma/A)\,\Sigma$, where $F$ is the curvature of the
SU(2) connection and $\Sigma$ is the densitised triad.  The area
operator is the scalar norm
$\hat{A} \propto \sqrt{\hat{\Sigma}\cdot\hat{\Sigma}}$, and its
eigenvalues on a puncture carrying spin $j$ are proportional to
$\sqrt{j(j+1)}$.

The Bekenstein--Hawking entropy counts the number of
\emph{distinct puncture configurations} for a given total area.
Maximising the entropy at fixed $A$ requires maximising the number
of punctures $N$, which in turn requires minimising the area per
puncture.  Since $a_j$ is monotonically increasing in $j$, this
restricts the emergence channels to the fundamental
representation:
\begin{equation}
j_{\rm emerge} = \tfrac{1}{2}\,.
\label{eq:j_emerge}
\end{equation}
This is the standard result of the ABCK entropy
counting~\cite{Ashtekar1998,Ashtekar2000,Ghosh2011}: the dominant
contribution to the microstate sum comes from configurations
in which all punctures carry the minimum spin.  For the emergence
sector of the cosmological horizon, the same logic applies: area
generation is an entropy-maximising process, and the most efficient
way to build area is to create $j = 1/2$ punctures.


\subsubsection{Interaction channels: \texorpdfstring{$j = 1$}{j = 1} (adjoint)}
\label{sec:interaction_channels}

The dark drag is not building static area; it is a continuous
dynamical flux of energy across the boundary.  In the quantum
theory, a transition in the state of the boundary \emph{without
creating a new puncture} requires varying the Chern-Simons
connection $A_a^i$.  Because $A_a^i$ is an
$\mathfrak{su}(2)$-valued gauge field, its phase-space tangent
vectors $\delta A_a^i$ and the corresponding triad fluxes
$\delta\Sigma^i$ take values in the Lie algebra
$\mathfrak{su}(2)$.  In representation theory, the Lie algebra
is the adjoint representation.

This fixes the interaction channel assignment.  The quantum
operator that mediates the transfer of energy across the boundary
is an internal vector transforming as
\begin{equation}
j_{\rm process} = 1\,.
\label{eq:j_process}
\end{equation}
The argument is straightforward: one cannot construct a macroscopic
flux vector from a single spinor ($j = 1/2$).  To exchange
physical momentum---a vector quantity---with the bulk, the horizon
must utilise its $\mathfrak{su}(2)$ flux operators, which carry
$j = 1$.  In the language of angular momentum coupling, the
minimum representation that supports a non-trivial vector
observable is the adjoint.

The distinction between the two channel classes is therefore rigid
and follows from the mathematical structure of the boundary theory:
\begin{itemize}
\item \emph{Emergence} (area generation) uses the scalar norm of
  the triad, optimises entropy, and operates at $j = 1/2$.
\item \emph{Interaction} (dark drag throughput) uses the vectorial
  flux of the triad and operates at $j = 1$.
\end{itemize}
The leading channel assignments are determined by the mathematical
structure of the boundary theory rather than freely chosen:

\paragraph{Adjoint selection rule for the interaction operator.}

To make the operator-level content of this assignment precise:
the exchange term on the horizon is sourced by an SU(2) current
contracted with the boundary connection, schematically
$\int_\Delta J_i\,A^i$.  Both $J_i$ and $A^i$ transform as
vectors under SU(2), i.e.\ in the adjoint (spin-$1$)
representation.  Accordingly, the puncture operator that mediates
\emph{transport} is a rank-1 (vector) tensor operator, and its
representation factor is the adjoint channel with
$\lambda_1 \propto (1-q)$.

This does \emph{not} forbid the geometric state-counting sector
from being dominated by many fundamental $j = 1/2$ punctures; it
only fixes the symmetry type of the operator insertion responsible
for the interaction channel (adjoint), while the storage and
degeneracy sector remains fundamental with
$\lambda_{1/2} \propto (1+q)$.

\paragraph{Effective truncation to $\{1/2, 1\}$.}

Higher-spin representations ($j \geq 3/2$) are mathematically
valid sectors of the SU(2) Chern-Simons Hilbert space, but they
are sub-dominant in the macroscopic dynamics by the two selection
rules above acting in concert.

For the storage sector, the information density of a spin-$j$
puncture is its entropy-to-area ratio:
\begin{equation}
s(j) = \frac{\ln(2j+1)}{\sqrt{j(j+1)}}\,.
\label{eq:info_density}
\end{equation}
This function is strictly monotonically decreasing for all
$j \geq 1/2$: the numerator $\ln(2j+1)$ grows logarithmically
while the denominator $\sqrt{j(j+1)}$ grows linearly, so
$s'(j) < 0$.  Explicitly: $s(1/2) \approx 0.800$,
$s(1) \approx 0.776$, $s(3/2) \approx 0.715$.  Incorporating any
puncture with $j > 1/2$ strictly lowers the total entropy for a
fixed area.  Because the cosmological apparent horizon saturates
the Bekenstein--Hawking bound, its emergence sector must exist in
the maximum-entropy ground state: a condensate of $j = 1/2$
punctures.

For the processing sector, the linear Israel--Darmois matching
couples the bulk current $J_i$ to the boundary connection
$\hat{A}^i$, both of which are $\mathfrak{su}(2)$-valued
(adjoint, $j = 1$).  There are no higher-rank tensor currents at
linear order that could excite $j \geq 3/2$ channels.

Therefore, while an isolated black hole horizon may permit thermal
fluctuations into higher-spin states, the \emph{macroscopic
fluxes} driving the cosmological apparent horizon are rigidly
partitioned: the entropy bound restricts the emergence sector to
$j = 1/2$, and the Lie algebra restricts the leading interaction
channel to $j = 1$.  The binary partition $\{1/2, 1\}$ is the
dominant macroscopic configuration; sub-leading corrections from
higher spins in the storage sector are quantified below.

\paragraph{Entropy estimate for higher-$j$ admixtures.}

The thermodynamic exclusion can be made quantitative.  Consider
the probability that, in the maximum-entropy ensemble of $N$
punctures at fixed total area $A$, at least one puncture carries
$j \geq 3/2$ rather than $j = 1/2$.  Each such substitution
incurs an entropy deficit per puncture of
\begin{equation}
\Delta s
  = s(1/2) - s(j)
  > 0\,,
\label{eq:entropy_deficit}
\end{equation}
with $\Delta s \approx 0.024$ for $j = 1$ and
$\Delta s \approx 0.085$ for $j = 3/2$.  The corresponding
suppression in the microcanonical count goes as
$\exp(-N_{\rm sub}\,\Delta s)$, where $N_{\rm sub}$ is the number
of substituted punctures.  For even a single substitution
($N_{\rm sub} = 1$) in the storage sector, the relative
probability is
\begin{equation}
\frac{P(j \geq 3/2)}{P(j = 1/2)}
  \sim \exp\!\bigl(-\Delta s\bigr)
  \approx 0.92 \quad (j = 3/2)\,,
\label{eq:single_sub}
\end{equation}
which is $\mathcal{O}(1)$ per puncture.  However, the area
released by replacing one $j = 1/2$ puncture with one $j = 3/2$
puncture must be compensated by removing $\Delta N$ other
$j = 1/2$ punctures to maintain the total area.  Since
$a_{3/2}/a_{1/2} = \sqrt{15/4}/\sqrt{3/4} \approx 2.24$, the
substitution removes $\sim 1.24$ punctures from the count.  The
net entropy change for the full configuration is therefore
\begin{equation}
\Delta S
  = \bigl[\ln 4 - 2.24\,\ln 2\bigr]
  \approx -0.17\,,
\label{eq:net_entropy}
\end{equation}
per substitution, and the probability of $n$ such substitutions
in an ensemble of $N \sim 10^{122}$ punctures is suppressed by
\begin{equation}
P(n) \;\propto\;
  \binom{N}{n}\,e^{-n\,|\Delta S|}
  \;\sim\; \left(\frac{N\,e^{-|\Delta S|}}{n}\right)^n.
\label{eq:higher_j_prob}
\end{equation}
The expected number of higher-$j$ punctures is
$\langle n \rangle \sim N\,e^{-|\Delta S|} \sim 0.84\,N$---which
is order-$N$ in absolute terms, i.e.\ an $\mathcal{O}(1)$ fraction of
storage punctures \emph{could} carry $j \geq 3/2$ rather than
$j = 1/2$ based on this entropy estimate alone.  The naive
exponential suppression is thus insufficient on its own to
guarantee a pure $j = 1/2$ storage sector.

This is where the interaction-sector selection rule becomes
decisive.  Even if the storage sector contains a thermal
admixture of $j = 1/2$ and higher-$j$ punctures, the
\emph{interaction operator} is constrained by the adjoint
selection rule: the source coupling $\int J_i\,\hat{A}^i$ excites
only $j = 1$ transitions.  The storage-sector composition
affects the area counting (and hence $k$), but it does not enter
the rate balance that determines $I_{\rm eff}$.  The saddle
allocation~\eqref{eq:saddle} depends only on the
\emph{interaction} rate $\lambda_1 = 1-q$ and the
\emph{emergence} rate $\lambda_{1/2} = 1+q$, both of which are
kinematic quantities independent of the storage-sector
composition.


\subsection{Mean processing spin tracks the deceleration parameter}
\label{sec:mean_spin}

With the channel populations identified, we can compute the mean
spin of the horizon's processing sector.  Emergence channels
carry $j = 1/2$ physically but zero interaction throughput;
Appendix~\ref{app:source_reduction} shows that the homogeneous
Israel--Darmois source couples only to the boundary monopole,
so emergence channels are invisible to it.  We therefore assign
$j = 0$ (the trivial representation) to
emergence channels in the throughput-weighted average, since they build area but carry no interaction
flux.  The mean spin of the full population is then
\begin{equation}
\langle j \rangle
  = j_{\rm process}\cdot I_{\rm eff}
    + 0 \cdot \Gamma_{\rm emerge}
  = 1 \cdot \frac{1-q}{2}
  = \frac{1-q}{2}\,,
\label{eq:mean_j}
\end{equation}
recovering the thermodynamic modulation factor $I_{\rm eff}$ as
the mean SU(2) spin of the interaction sector.  The macroscopic
efficiency of the dark drag equals the quantum number that
characterises the average state of the boundary punctures
participating in the interaction.

The boundary values of $\langle j \rangle$ at cosmologically
distinguished epochs acquire representation-theoretic significance:

\begin{center}
\small
\begin{tabular}{lcccp{3.2cm}}
\toprule
Epoch & $q$ & $\langle j\rangle$ & Eff.\ rep. &
  Interpretation \\
\midrule
Radiation & $+1$ & $0$ & Singlet &
  Horizon transparent \\
Matter & $+\tfrac{1}{2}$ & $\tfrac{1}{4}$ &
  Mixed &
  25\% active at $j\!=\!1$ \\
Accel.\ onset & $0$ & $\tfrac{1}{2}$ & $\langle j\rangle\!=\!\tfrac{1}{2}$ &
  Non-trivial channel opens \\
de~Sitter & $-1$ & $1$ & Adjoint &
  Full saturation \\
\bottomrule
\end{tabular}
\end{center}

Several features of this table deserve comment.

The radiation-era value $\langle j \rangle = 0$ means the singlet
representation: all horizon capacity is consumed by emergence
($\Gamma_{\rm emerge} = 1$), every puncture is building area, and
no channels are available for energy transfer.  The dark sector is
completely decoupled---not because of an ad hoc suppression
mechanism, but because the horizon's quantum state has no room
for adjoint transitions.

The matter-era value $\langle j \rangle = 1/4$ does not violate
SU(2) quantisation: it is a macroscopic statistical average over
the ensemble, corresponding to 25\% of channels processing at
$j = 1$ while the remaining 75\% are engaged in area generation
(equivalently, 50\% active at $j = 1/2$ via the Clebsch--Gordan
decomposition
$\tfrac{1}{2} \otimes \tfrac{1}{2} = 0 \oplus 1$).  The
non-integer value of $\langle j \rangle$ is no more problematic
than a non-integer mean occupation number in a Bose gas.

The acceleration onset at $q = 0$ marks the moment when
$\langle j \rangle = 1/2$: exactly one-half of the horizon's
capacity is allocated to interaction, and the processing sector
reaches a throughput-weighted mean spin of $\langle j\rangle = 1/2$---not
the actual irrep of active channels (which remain $j=1$), but the
population mean at the threshold where exactly half of all channels
are active.  Below this threshold, the interaction is sub-dominant;
above it, the adjoint sector begins to dominate.

The de~Sitter limit $\langle j \rangle = 1$ is the unique fixed
point: emergence ceases ($\Gamma_{\rm emerge} = 0$), the full
processing capacity is allocated to the adjoint ($j = 1$) channel,
and the interaction sector is fully saturated.  The boundary theory
reaches thermal equilibrium with the bulk---but simultaneously,
$Q \propto \Omega_m \to 0$, so the interaction shuts off from
the matter side.  The system converges to zero-dissipation
equilibrium from both channels, a self-consistent infrared
fixed point.

The identification $\langle j \rangle = I_{\rm eff}$ is the
bridge between the macroscopic and microscopic descriptions.
The two macroscopic derivations and kinematic restatement of Section~\ref{sec:three_routes}
established that $I_{\rm eff} = (1-q)/2$; this section has shown
that the same quantity is the mean SU(2) spin of the interaction
sector on the apparent horizon.  The next section derives the
screening function $x(q)$ by evaluating the processing amplitude
at this mean spin.


% ======================================================================
%  SECTION 4: DERIVATION OF THE SCREENING FUNCTION
% ======================================================================
% ======================================================================
%  SECTION 4: DERIVATION OF THE SCREENING FUNCTION
% ======================================================================
\section{Derivation of the Screening Function}
\label{sec:derivation}

Section~\ref{sec:screening_problem} defined the screening function
$x(q) = (1-q)(2-q)/6$ as the ratio needed to reconcile the
Lagrangian and thermodynamic modulations.
Section~\ref{sec:mean_spin} identified the mean interaction spin
$\langle j \rangle = (1-q)/2$.  We first rewrite $x$ in
representation-theoretic language, then derive the same polynomial
from a boundary state-sum calculation using the Chern-Simons
Hilbert space.


\subsection{Representation-theoretic reparameterisation}
\label{sec:algebraic}

The screening function~\eqref{eq:x_def} is
$x(q) = (1-q)(2-q)/6$.  We rewrite this in terms of the mean
interaction spin $\langle j \rangle = (1-q)/2$
(Section~\ref{sec:mean_spin}; note that $\langle j \rangle$ is
the interaction-throughput-weighted spin, not the physical mean
spin of all punctures---see
Section~\ref{sec:macroscopic_screening} for the precise
distinction).  The two factors
are:
\begin{align}
1-q &= 2\langle j\rangle\,,
\label{eq:factor1} \\
2-q &= 2\langle j\rangle + 1 \equiv d_{\langle j\rangle}\,,
\label{eq:factor2}
\end{align}
where $d_{\langle j\rangle} = 2\langle j\rangle+1$ is the
effective degeneracy factor, coinciding with the SU(2)
representation dimension when $\langle j \rangle$ takes a
half-integer value but well-defined as a continuous function for
all $\langle j \rangle$.  Here $\langle j \rangle$ is the
\emph{interaction-throughput-weighted} mean spin (Section~\ref{sec:mean_spin}):
emergence channels are assigned $j^{\rm int} = 0$ rather than
their physical $j = 1/2$, because they carry no interaction flux.
This is the same throughput-weighted average used throughout;
$d_{\langle j\rangle}$ is therefore not the literal dimension of
a spin-$\langle j\rangle$ irrep evaluated at a non-half-integer
argument, but a continuous interpolation of that degeneracy
appropriate to the throughput label.  Substituting:
\begin{equation}
\boxed{x = \frac{\langle j\rangle \cdot d_{\langle j\rangle}}
  {d_{\rm adj}}
  = \frac{\langle j\rangle\,(2\langle j\rangle + 1)}{3}\,,}
\label{eq:x_spin}
\end{equation}
where $d_{\rm adj} = d_1 = 3$ is the dimension of the adjoint
representation, which normalises to the de~Sitter maximum
($\langle j \rangle = 1$, $x = 1$).

The physical content of this factorisation is as follows.  The
macroscopic (population-averaged) processing rate is proportional to
$\langle j \rangle \cdot d_{\langle j \rangle}$; individual active channels
operate at $j=1$ with single-channel rate $\Gamma(1) \propto 1 \cdot 3$,
and $\langle j \rangle \cdot d_{\langle j \rangle}$ is the mean-field product
after averaging over the channel population.  The first factor,
$\langle j \rangle$, carries the flux magnitude: the
expectation value of the triad operator in the coherent state
$|j, j\rangle$ scales linearly with $j$
(Section~\ref{sec:state_sum}).  The second factor,
$d_{\langle j \rangle} = 2\langle j \rangle + 1$, counts the
internal microstates available for transitions within the SU(2)
Chern-Simons Hilbert space at each puncture---the degeneracy of
the magnetic quantum number.  The normalisation to $d_{\rm adj}$
ensures $x \to 1$ when the processing sector is fully saturated
in the adjoint.

The identification~\eqref{eq:x_spin} is remarkable because each
ingredient has an independent origin:
\begin{itemize}
\item $\langle j \rangle = (1-q)/2$ was derived macroscopically
  from two independent macroscopic derivations and a kinematic restatement
  (Section~\ref{sec:three_routes}) and microscopically from the
  channel partition (Section~\ref{sec:mean_spin}).
\item $d_{\langle j \rangle} = 2-q$ is a representation-theoretic
  quantity---the continuous interpolation of the SU(2) dimension
  function $2j+1$, evaluated at the mean spin $\langle j \rangle$
  (coinciding with the literal irrep dimension at half-integer values)---that
  happens to equal the denominator of the Lagrangian
  modulation $I_{\rm eff}^{\rm Lag} = 3/(2-q)$.
\item The normalisation $d_{\rm adj} = 3$ is the numerator of the
  same Lagrangian modulation.
\end{itemize}
The screening function is therefore the ratio of the effective
representation dimension to the maximum representation dimension,
weighted by the spin flux.  The Lagrangian's ``wrong'' modulation
$3/(2-q)$ is not wrong at all: it is the inverse of the
representation-theoretic factor $d_{\langle j \rangle}/d_{\rm adj}$,
and the screening function $x$ simply supplies the missing
microscopic weight.


\subsection{Boundary state-sum derivation}
\label{sec:state_sum}

We now derive the same result directly from the boundary state-sum
of the SU(2) Chern-Simons theory, without assuming the algebraic
identification.

In the quantum theory, the dark drag is a continuous thermodynamic
flux across the apparent horizon.  The transition amplitude
$\mathcal{A}$ for a quantum of vacuum energy to interact with the
bulk is determined by a boundary state calculation.  By analogy
with Fermi's Golden Rule, the macroscopic processing rate per
channel, $\Gamma(j)$, is proportional to the geometric flux
carried by the puncture multiplied by the density of available
microscopic transitions:
\begin{equation}
\Gamma(j) \;\propto\;
  \langle \hat{E}_{\rm normal} \rangle
  \;\times\; \rho(j)\,.
\label{eq:Gamma_j}
\end{equation}
Both factors are strictly determined by the SU(2) representation
theory of the boundary.

\paragraph{Density of states $\rho(j)$.}

The statistical weight of an SU(2) boundary puncture carrying
spin $j$ is given by the Plancherel measure, which for a compact
Lie group counts the number of independent states within the
irreducible representation:
\begin{equation}
\rho(j) = d_j = 2j + 1\,.
\label{eq:plancherel}
\end{equation}
This is a kinematic property of SU(2)---it does not depend on the
dynamics of the Chern-Simons theory, the value of $k$, or the
Barbero--Immirzi parameter.  A spin-$j$ puncture has $2j+1$
magnetic substates available for transitions, and each substate
represents an independent microscopic channel through which energy
can cross the boundary.

\paragraph{Geometric flux $\langle \hat{E}_{\rm normal} \rangle$.}

The physical momentum flux across the horizon is governed by the
densitised triad operator $\hat{E}^i$, which in the Ashtekar
variables is canonically conjugate to the connection $A_a^i$.  The
ABCK boundary condition
$F = -(\pi\gamma/A)\,\Sigma$ selects the outward-pointing triad
component, restricting the flux to the coherent sector aligned with
the horizon normal $\hat{n}$.  For a puncture in the SU(2) coherent
highest-weight state $|j, j\rangle$ (the semiclassical state
aligned with $\hat{n}$), the expectation value of the normal flux
operator scales linearly with the spin:
\begin{equation}
\langle j, j \,|\, \hat{E}_{\rm normal} \,|\, j, j \rangle
  \;\propto\; j\,.
\label{eq:flux_j}
\end{equation}
This linear scaling is standard in angular momentum
theory~\cite{Hall2015}: the highest-weight state $|j, j\rangle$
is the most classical state in the representation, and its
projection onto the quantisation axis is $m = j$.

\paragraph{The processing rate.}

Combining the density of states and the geometric flux, the
thermodynamic throughput of a single horizon channel operating at
spin $j$ is:
\begin{equation}
\Gamma(j) \;\propto\; j\,(2j+1)\,.
\label{eq:Gamma_j_explicit}
\end{equation}
This is a quadratic function of $j$---a fact with important
consequences for the ensemble averaging below.  Normalising to the
maximum processing capacity at de~Sitter equilibrium, where all
interaction channels are saturated in the adjoint representation
($j = 1$, $\Gamma_{\rm max} = 1 \cdot 3 = 3$):
\begin{equation}
x(j) = \frac{\Gamma(j)}{\Gamma_{\rm max}}
  = \frac{j\,(2j+1)}{3}\,.
\label{eq:x_j}
\end{equation}
This is the normalised processing amplitude of a single channel
at spin $j$.  At this stage, $j$ is a free label.  The
macroscopic screening function is obtained by evaluating $x$ at
the physically relevant spin, which requires specifying the
ensemble-averaging prescription.


\subsection{Coherent processing: the zero-mode and the quantum
effective action}
\label{sec:mean_field}

A physical subtlety arises in how the single-channel amplitude
$x(j) = j(2j+1)/3$ is averaged over the macroscopic boundary.
The puncture population is bimodal, with fractions
$\Gamma_{\rm emerge} = (1+q)/2$ at $j = 1/2$ (kinematically
frozen from the interaction perspective) and
$I_{\rm eff} = (1-q)/2$ at $j = 1$.  Two natural prescriptions
yield different results.

\paragraph{Exact microscopic average (throughput-weighted).}

An important distinction arises between two possible averages.
The \emph{physical} mean spin of all punctures in the
$\{1/2,\,1\}$ model is
$\bar{j}_{\rm phys} = (1-f)\cdot\tfrac{1}{2} + f\cdot 1
= (1+f)/2 = (3-q)/4$.  But the screening function
$x(j) = j(2j+1)/3$ measures \emph{interaction throughput}, and
emergence channels carry zero throughput (they build area, not
mediate energy transfer; Appendix~\ref{app:source_reduction} derives
this from the zero-mode projection of the sourced boundary path
integral).  The throughput-weighted average
therefore assigns $x = 0$ to emergence channels, regardless of
their physical spin:
\begin{align}
\langle x(j) \rangle_{\rm throughput}
  &= I_{\rm eff} \cdot x(1) + \Gamma_{\rm emerge} \cdot 0
\nonumber \\
  &= \frac{1-q}{2} \cdot 1 + \frac{1+q}{2} \cdot 0
  = \frac{1-q}{2}\,.
\label{eq:x_exact}
\end{align}
This recovers $I_{\rm eff}$ itself---no new information beyond
the capacity fraction.  The throughput-weighted average simply
counts the proportion of active channels; it discards the
representation-theoretic nonlinearity of the processing amplitude.

\paragraph{Mean-field (VEV) evaluation.}

Evaluating the processing rate at the ensemble-averaged spin
$\langle j \rangle = (1-q)/2$ instead gives:
\begin{equation}
x\bigl(\langle j \rangle\bigr)
  = \frac{\langle j \rangle\,(2\langle j \rangle + 1)}{3}
  = \frac{(1-q)(2-q)}{6}\,.
\label{eq:x_meanfield}
\end{equation}
The two expressions differ for $-1 < q < 1$, and the data
strongly favour the latter~\cite{Sayegh2026}: the
throughput-weighted average over-suppresses $\sigma_8$ by a
factor of ${\sim}\,2.4$ (relative to $\Lambda$CDM), while the
mean-field screening---and only the mean-field
screening---restores agreement with the observed growth rate.

\paragraph{Why the macroscopic source uses the mean-field closure.}

If the interaction were a sum of \emph{microscopic} local
scattering events at the horizon, the relevant rate would be the
puncture-average of a nonlinear function,
$\langle \Gamma(j) \rangle = \sum_j p_j\,\Gamma(j) \propto
\langle j(2j+1) \rangle$.  In contrast, the dark drag kernel
$Q \propto H\rho_{\rm crit}\Omega_\Lambda\Omega_m$ is a strictly
homogeneous ($k = 0$) background current.  The junction-condition
coupling is therefore a \emph{macroscopic} source acting on the
horizon as a whole:
\begin{equation}
\hat{H}_{\rm source}
  = \int_\Delta J_i\,\hat{A}^i
  \;\longrightarrow\;
  J_i\,\overline{A}^{\,i}\,A_\Delta\,,
\label{eq:source_zero_mode}
\end{equation}
where
$\overline{A}^{\,i} \equiv A_\Delta^{-1}\int_\Delta \hat{A}^i$
denotes the coarse-grained (spatially averaged) boundary field
picked out by the monopole source, and $A_\Delta$ is the horizon
area.  Equivalently (and gauge-invariantly), one may regard
$\overline{A}^{\,i}$ as the order parameter controlling the
corresponding averaged boundary flux.  (Higher spatial multipoles,
sourced by cosmological perturbations, enter at first order and
contribute sub-percent corrections to the background interaction.)

The macroscopic response is governed by the quantum effective
action for this coarse-grained order parameter, so the leading
contribution is obtained by evaluating the rate functional at the
background (mean) boundary charge rather than by averaging over
puncture-by-puncture fluctuations.  In the puncture ensemble, this
background charge is proportional to the first moment
$\langle j \rangle$, and hence the macroscopic processing rate
takes the mean-field form:
\begin{equation}
\Gamma_{\rm mac}
  \;\propto\; \Gamma\!\bigl(\langle j \rangle\bigr)
  \;\propto\; \langle j \rangle\,
    \bigl(2\langle j \rangle + 1\bigr)\,,
\label{eq:Gamma_mean_field}
\end{equation}
which is precisely the factorised closure used in the screening
function $x \equiv \langle j \rangle\,\langle d \rangle/3$.
Finally, because the apparent horizon is deep in the thermodynamic
limit ($N \sim A/(4\ell_P^2) \gg 1$), fluctuations of the
coarse-grained monopole are suppressed,
$\delta\overline{A}/\overline{A} = \mathcal{O}(N^{-1/2})
\sim 10^{-61}$, so corrections to the mean-field evaluation are
parametrically small.

\paragraph{Fluctuation correction to the mean-field closure.}

The screening function $x(j) = j(2j+1)/3$ is exactly quadratic in
$j$, with $x''(j) = 4/3$ everywhere.  Define the activity variable
$X_p \in \{0,\,1\}$ for each puncture ($X_p = 0$ for emergence,
$X_p = 1$ for interaction), so that the interaction-weighted spin
is $j_p^{\rm int} = X_p \cdot j_{\rm process} = X_p$.  The
empirical mean is
$\bar{j} = N^{-1}\sum_{p=1}^{N} X_p$, drawn from a binomial
with $f = (1-q)/2$.  The variance is
\begin{equation}
\mathrm{Var}(\bar{j})
  = \frac{f(1-f)}{N}
  = \frac{1-q^2}{4N}\,.
\label{eq:var_jbar}
\end{equation}
By Taylor expansion of $x$ around $\langle j \rangle$, the
expectation of $x(\bar{j})$ over realisations receives a
correction:
\begin{align}
\langle x(\bar{j}) \rangle
  &= x\bigl(\langle j \rangle\bigr)
    + \tfrac{1}{2}\,x''\,\mathrm{Var}(\bar{j})
    + \mathcal{O}(N^{-2})
\nonumber \\
  &= \frac{(1-q)(2-q)}{6}
    + \frac{1-q^2}{6N}\,.
\label{eq:x_fluctuation}
\end{align}
For the cosmological horizon, $N = A/(4\ell_P^2) \sim 10^{122}$,
so the correction is $\sim 10^{-123}$---smaller than the
$\epsilon_2 = \mathcal{O}(1/k)$ corrections already bounded
in~\eqref{eq:eps2_bound}.  The mean-field closure is exact to all
observationally relevant precision.


\subsection{The macroscopic screening factor}
\label{sec:macroscopic_screening}

\paragraph{Discrete mixture with factorised screening moment.}

We model the punctures by a discrete distribution $p_j$ over
SU(2) irreps, $\sum_j p_j = 1$, with first moments
\begin{equation}
\langle j \rangle \equiv \sum_j p_j\,j\,,\qquad
\langle d \rangle \equiv \sum_j p_j\,(2j+1)\,.
\label{eq:discrete_moments}
\end{equation}
The screening function used in this paper is the
\emph{factorised} mean-field moment
\begin{equation}
x \;\equiv\; \frac{\langle j \rangle\,\langle d \rangle}{3}\,,
\label{eq:x_factorised}
\end{equation}
rather than $\langle j(2j+1) \rangle/3$.  The factorisation
reflects the physical independence of the two contributions:
$\langle j \rangle$ measures what fraction of channels process
(the flux), while $\langle d \rangle$ counts the average number
of internal substates available per channel (the degeneracy).
Since $d(j) = 2j+1$ is linear in $j$, the ensemble-averaged
dimension satisfies
$\langle d \rangle = 2\langle j \rangle + 1 \equiv d_{\langle
j \rangle}$ exactly---no Jensen correction arises for the
degeneracy factor individually.  The factorisation matters only
because the \emph{product} $j \cdot d_j = j(2j+1)$ is quadratic
in $j$: $\langle j \cdot d_j \rangle \neq \langle j \rangle
\cdot \langle d \rangle$ in general.  The definition~\eqref{eq:x_factorised} adopts the factorised (mean-field)
closure; the alternative $\langle j(2j+1) \rangle/3$ corresponds
to a different microscopic closure and is not the screening moment
adopted here.

In the two-channel model ($j = 1/2$ storage, $j = 1$ transport)
with adjoint fraction $f \equiv p_1$, the \emph{physical} spin
moments are
\begin{equation}
\bar{j}_{\rm phys} = \tfrac{1}{2}(1 + f)\,,\qquad
\bar{d}_{\rm phys} = 2 + f\,.
\label{eq:moments_phys}
\end{equation}
Imposing $f = (1-q)/2$ gives $\bar{j}_{\rm phys} = (3-q)/4$ and
$\bar{d}_{\rm phys} = (5-q)/2$, which yields
$\bar{j}_{\rm phys}\,\bar{d}_{\rm phys}/3 = (3-q)(5-q)/24$---not
the target $(1-q)(2-q)/6$.

The screening function uses the \emph{interaction-throughput}
moments instead.  Emergence channels carry $j = 1/2$ physically
but contribute zero interaction throughput (they build area, not
mediate flux).  In the throughput-weighted average, emergence
channels are assigned $j^{\rm int} = 0$, giving
\begin{align}
\langle j \rangle
&=0 \cdot (1-f)+1 \cdot f=f,\notag\\
\langle d \rangle
&=1 \cdot (1-f)+3 \cdot f=1+2f.
\label{eq:moments_f}
\end{align}
Imposing the saddle allocation $f = (1{-}q)/2$
from~\eqref{eq:saddle} gives $\langle j \rangle = (1{-}q)/2$ and
$\langle d \rangle = 2 - q$, hence
\begin{equation}
x(q) = \frac{\langle j \rangle\,\langle d \rangle}{3}
  = \frac{(1-q)/2 \;\cdot\; (2-q)}{3}
  = \frac{(1-q)(2-q)}{6}\,.
\label{eq:x_q_final}
\end{equation}
This recovers the screening function defined
in~\eqref{eq:x_def}, now derived from the boundary state-sum
using the interaction-throughput weighting.  The factorisation is
physical: $\langle j \rangle$ measures the fraction of channels
that process, while $\langle d \rangle / d_{\rm adj}$ measures
how efficiently each processes.

The polynomial structure has a clean decomposition into its
representation-theoretic content:
\begin{equation}
x(q) = \underbrace{\frac{1-q}{2}}_{\displaystyle
  \langle j \rangle\;\text{(flux)}}
  \;\times\;
  \underbrace{\frac{2-q}{3}}_{\displaystyle
  \langle d \rangle/d_{\rm adj}\;\text{(degeneracy)}}\,.
\label{eq:x_decomposed}
\end{equation}
The first factor is the linear-response amplitude---the fraction
of the horizon's capacity allocated to interaction, equivalently
the mean spin of the processing sector.  The second factor is
the representation-theoretic weight---the ratio of the effective
degeneracy $d_{\langle j \rangle} = 2\langle j \rangle + 1$ to
the maximum degeneracy $d_{\rm adj} = 3$ at de~Sitter.  The
screening function is therefore the product of \emph{what fraction of
channels process} and \emph{how efficiently each processes}.

The boundary values confirm the interpretation:
\begin{center}
\begin{tabular}{lcccc}
\toprule
Epoch & $q$ & $\langle j \rangle$ &
  $d_{\langle j \rangle}/d_{\rm adj}$ & $x$ \\
\midrule
Radiation & $+1$ & $0$ & $1/3$ & $0$ \\
Matter & $+1/2$ & $1/4$ & $1/2$ & $1/8$ \\
Accel.\ onset & $0$ & $1/2$ & $2/3$ & $1/3$ \\
de~Sitter & $-1$ & $1$ & $1$ & $1$ \\
\bottomrule
\end{tabular}
\end{center}

\noindent At radiation domination, $x = 0$ because
$\langle j \rangle = 0$: no channels process, regardless of
degeneracy.  At matter domination, $x = 1/8$: only a quarter of
channels are active, and each operates at half the degeneracy
of the adjoint.  At acceleration onset, $x = 1/3$: one-half
of channels active at two-thirds degeneracy.  At de~Sitter,
$x = 1$: all channels active at full adjoint degeneracy.  The
progression from singlet to adjoint is monotonic and smooth,
with the degeneracy factor providing additional suppression at
intermediate epochs beyond what the linear capacity fraction alone
would give.

This completes the derivation of the screening function from
the SU(2) boundary theory.  The field-theory translation factor
$x(q)$ is not phenomenological; it is the normalised transition
amplitude of the Chern-Simons theory on the apparent horizon,
evaluated at the thermodynamic saddle point.  The next section
provides an independent quantum-gravitational derivation of
$I_{\rm eff}$ via KMS thermodynamics---complementing the two
macroscopic derivations and kinematic restatement of
Section~\ref{sec:three_routes}---and establishes the structural
protection of the saddle-point result.


% ======================================================================
%  SECTION 5: INDEPENDENT KMS DERIVATION AND STRUCTURAL PROTECTION
% ======================================================================
\section{Independent KMS Derivation and Structural Protection}
\label{sec:kms}

The two macroscopic derivations and kinematic restatement of $I_{\rm eff}$ in
Section~\ref{sec:three_routes} share a common feature: each
produces the factors $(1-q)$ and $(1+q)$ but from different
physical inputs.  The KMS route presented here is a fourth,
independent derivation that operates entirely within the boundary
quantum theory---deriving $(1-q)$ from the fluctuation-dissipation
theorem applied to the Kodama surface gravity and $(1+q)$ from the
kinematic evolution of the Chern-Simons level.  The route also
reveals why the equal-footing normalisation $c_e = c_i$ is not
an assumption but a structural consequence of the boundary theory,
establishing the robustness of the saddle-point allocation against
renormalisation.


\subsection{KMS thermodynamics on the boundary}
\label{sec:kms_derivation}

\paragraph{The interaction rate.}

The Kodama--Hayward surface gravity of the apparent
horizon~\cite{Hayward1998} is a purely geometric quantity:
\begin{equation}
\kappa_{\rm KH}
  = -r_A\!\left(H^2 + \frac{\dot{H}}{2}\right)
  = -\frac{cH(1-q)}{2}\,,
\label{eq:kappa_KH_kms}
\end{equation}
derived from $r_A = c/H$ and $\dot{H} = -(1+q)H^2$, with no
thermodynamic input.  The KMS theorem guarantees that the boundary
quantum theory, evolved by the Kodama flow~\cite{Kodama1980}, is thermal at
temperature
\begin{equation}
T_{\rm KH} = \frac{\hbar\,|\kappa_{\rm KH}|}{2\pi k_B c}
  = \frac{\hbar H(1-q)}{4\pi k_B}\,.
\label{eq:TA_KMS}
\end{equation}
We write $T_{\rm KH}$ for this dynamical Kodama--Hayward
temperature to distinguish it from the de~Sitter-normalised
horizon temperature $T_{\rm dS} = \hbar H/(2\pi k_B)$ of
the companion cosmology paper~\cite{Sayegh2026} (denoted $T_A$ there); the two are
related by $T_{\rm KH} = \tfrac{1-q}{2}\,T_{\rm dS}$ and coincide
at the de~Sitter fixed point.

The use of the KMS transition rate $\lambda_1 = 1-q$ relies on
the apparent horizon maintaining a thermal state at $T_{\rm KH}$.
In the local-equilibrium approximation, this requires that the
boundary's internal relaxation be fast compared to the macroscopic
driving rate $|\dot{q}|/H$.  We now evaluate both timescales.

The deceleration parameter is $q = 1/2 - (3/2)\Omega_\Lambda$.
Differentiating, noting that $\dot{\rho}_\Lambda = 0$ in the
background, the time derivative falls entirely on the critical
density $\rho_c \propto H^2$.  Using $\dot{H}/H = -(1+q)H$ and
$1 + q = (3/2)\Omega_m$:
\begin{equation}
\frac{|\dot{q}|}{H}
  = \frac{9}{2}\,\Omega_\Lambda\,\Omega_m
  \equiv \mathcal{S}(z)\,.
\label{eq:qdot_sweep}
\end{equation}
This reveals a profound consistency: the dimensionless evolution
rate of the horizon's quantum state is exactly the
super-Schwarzschild sweep ratio $\mathcal{S}(z)$ defined in
Section~\ref{sec:kernel_recap}.  Crucially, the derivation above
uses nothing but the Friedmann equations---no horizon
thermodynamics, no temperature, no Schwarzschild bound.  The
quantity $(9/2)\,\Omega_\Lambda\,\Omega_m$ was derived in
the companion cosmology paper~\cite{Sayegh2026} as the ratio
$\dot{M}_\Lambda/(c^3/2G)$, i.e.\ from thermodynamic throughput
saturation.  Here it emerges as $|\dot{q}|/H$,
i.e.\ from pure FLRW kinematics---the kinematic identity also
recorded in the companion cosmology paper~\cite{Sayegh2026}.  The two routes---one
thermodynamic, one kinematic---converge on the same dimensionless
function without any shared assumption.

In the deep matter era ($\Omega_\Lambda \to 0$) and the de~Sitter
limit ($\Omega_m \to 0$), $\mathcal{S}(z) \ll 1$: the horizon
evolves quasi-statically and the dissipative cost is small.
During the transition epoch, $\mathcal{S}$ peaks at
$9/8 = 1.125$, and the dark drag reaches its maximum strength.
The super-Schwarzschild window ($\mathcal{S} > 1$) is not a
threshold that switches the interaction on; the dark drag
$Q \propto H\rho_{\rm crit}\Omega_\Lambda\Omega_m$ operates
continuously whenever both dark energy and matter are present.
Rather, $\mathcal{S} > 1$ identifies the epoch where the driving
is strongest and where the accumulated throughput produces the
observable mapping between the CMB-inferred global $H_0$ and the
locally measured expansion rate.

The dark drag is the horizon's continuous thermodynamic response
to being driven through its configuration space by the cosmic
expansion.  The driving generates irreversible Prigogine entropy
production $d_i S \geq 0$, with the dissipative flux $Q$
self-consistently tracking the driving rate $\mathcal{S}(z)$.
The horizon processes the vacuum energy exactly as fast as it
can---via the drag---and the drag strength is set by the spin
velocity (below).

Two distinct timescales should be distinguished.  The boundary
Chern-Simons theory equilibrates \emph{internally} on a timescale
much shorter than $H^{-1}$: the conformal dimension gap is
\begin{equation}
\Delta h \sim \frac{j(j+1)}{k+2} = \mathcal{O}(k^{-1}) \ll 1\,,
\end{equation}
so the thermally accessible level density is enhanced by
$\Delta h^{-1} \sim k \gg 1$, meaning all states within each
representation are thermally accessible.  In this regime, the
fluctuation-dissipation theorem guarantees an internal relaxation
rate $\Gamma_{\rm int} \sim (k_B T_{\rm KH}/\hbar)\,\Delta h^{-1}
\sim k\,H$, parametrically faster than the macroscopic driving
rate $|\dot{q}|$.  Comparing rates of the same dimension:
$\Gamma_{\rm int}/H \sim k \sim 10^{122}$, while
$|\dot{q}|/H \sim \mathcal{S}(z) = \mathcal{O}(1)$
(where $\mathcal{S}(z)$ is an $\mathcal{O}(1)$ function encoding
the slow variation of $q$ along the cosmic history).
The ratio $\Gamma_{\rm int}/|\dot{q}| \sim k/\mathcal{S}(z) \gg 1$
confirms adiabatic tracking: the thermal state at
$T_{\rm KH}$ therefore tracks $\kappa_{\rm KH}$ instantaneously to
$\mathcal{O}(H/\Gamma_{\rm int}) = \mathcal{O}(1/k)
\sim 10^{-122}$; Appendix~\ref{app:source_reduction} computes the
explicit linear-response correction
$\delta I_{\rm eff} \leq 9/(16k)$, confirming this bound.
The dissipation arises because the continuous
driving of the boundary state through the Chern-Simons
configuration space has an irreducible thermodynamic cost.

\paragraph{Spin velocity and boundary dissipation.}

The continuous dissipation admits a geometric interpretation
within the microscopic SU(2) boundary theory.  The thermodynamic
state of the horizon's processing sector is characterised by the
mean spin $\langle j \rangle = (1-q)/2$.  Its rate of change per
e-fold is
\begin{equation}
\frac{d\langle j \rangle}{d\ln a}
  = -\frac{\dot{q}}{2H}
  = \frac{1}{2}\,\mathcal{S}(z)\,.
\label{eq:spin_velocity}
\end{equation}
The super-Schwarzschild sweep ratio $\mathcal{S}(z)$ is therefore
exactly twice the ``spin velocity'' of the horizon's state vector
through the Chern-Simons configuration space.

The dark drag is the thermodynamic cost of this driving.  In the
quantum theory of driven systems, varying an external parameter at
finite velocity induces irreversible dissipation proportional to
the driving speed.  The dark drag
$Q \propto H\rho_{\rm crit}\Omega_\Lambda\Omega_m \propto
\mathcal{S}(z)$ is the bulk manifestation of this microscopic
dissipative boundary response: the continuous entropy-producing flux generated by
dragging the boundary's SU(2) state through its configuration
space at the spin velocity $\mathcal{S}/2$.  The drag is weakest
in the deep matter era and de~Sitter limit (where
$\mathcal{S} \to 0$), strongest during the transition epoch
(where $\mathcal{S}$ peaks), and always proportional to the
instantaneous driving rate.

For the SU(2) Chern-Simons theory on the horizon, the conformal
dimension gap between representations scales as
$\Delta h \sim j(j+1)/(k+2)$.  In the macroscopic limit
$k \sim 10^{122}/\gamma$ (Section~\ref{sec:dynamical_k}), this
gap is negligible: $\Delta h \sim \mathcal{O}(k^{-1}) \ll 1$,
and the boundary theory
is deep in the classical regime.  The fluctuation-dissipation
theorem then gives a transition rate per puncture proportional to
the temperature:
\begin{equation}
\Gamma_{\rm interact} \;\propto\; T_{\rm KH} \;\propto\; (1-q)\,.
\label{eq:rate_interact}
\end{equation}
This derivation of the $(1-q)$ factor is independent of the
Onsager ratio (which uses $|\kappa_{\rm KH}|/|\kappa_{\rm dS}|$
directly) and of the Padmanabhan capacity split (which subtracts
$\Gamma_{\rm emerge}$ from unity).  It follows from the KMS
condition applied to the Kodama surface gravity---a standard result
in quantum field theory on curved backgrounds, here applied to
the boundary Chern-Simons theory rather than to bulk fields.


\paragraph{The emergence rate.}

The Chern-Simons level evolves as
$dk/d\ln a = 2(1+q)k$~\eqref{eq:dk_dlna}.  Since $k$ counts
the total boundary degrees of freedom, the fractional rate at
which new punctures are created is
\begin{equation}
\Gamma_{\rm emerge} \;\propto\;
  \frac{1}{k}\frac{dk}{dt} \;=\; 2(1+q)\,H\,.
\label{eq:rate_emerge}
\end{equation}
The factor of $2$ is absorbed into the normalisation constants
in the competitive allocation below, so the
\emph{proportionality} $\Gamma_{\rm emerge} \propto (1+q)$ is the
physically relevant content.  This is purely kinematic: it uses
$\dot{H} = -(1+q)H^2$ and the definition of $k$.  No
thermodynamic assumption enters.


\paragraph{Competitive allocation.}

The two processes---puncture creation at rate $\propto (1+q)$ and
state transitions at rate $\propto (1-q)$---compete for the
horizon's finite processing capacity.  The fraction of total
activity allocated to interaction is
\begin{equation}
f_{\rm interact}
  = \frac{c_e\,(1-q)}{c_e\,(1-q) + c_i\,(1+q)}\,,
\label{eq:competitive}
\end{equation}
where $c_e$ and $c_i$ are the absolute normalisations of
the emergence and interaction rates respectively (distinct
from the physical coupling $\beta = 1/12$).  Under the
\emph{equal-footing normalisation} $c_e = c_i$---derived from
the first-order CS kinetic term
(Section~\ref{sec:canonical};
Appendix~\ref{app:source_reduction})---this reduces to
\begin{equation}
f_{\rm interact}
  = \frac{1-q}{(1-q) + (1+q)}
  = \frac{1-q}{2}
  = I_{\rm eff}\,.
\label{eq:f_from_KMS}
\end{equation}

Four derivations, each using different physical inputs, now yield
the same allocation:

\begin{center}
\footnotesize
\begin{tabular}{lccc}
\toprule
Route & $(1{-}q)$ source & $(1{+}q)$ source &
  Norm. \\
\midrule
Onsager & $|\kappa_{\rm KH}|/|\kappa_{\rm dS}|$ &
  (implicit) & de~Sitter ref.\ \\
Padmanabhan & $1 - \Gamma_{\rm emerge}$ &
  $N_{\rm sur} - N_{\rm bulk}$ &
  $N_{\rm sur}$ = total \\
Surf.\ gravity & $\dot{H}/H^2$ ratio &
  (implicit) & Friedmann \\
CS/KMS & $T_{\rm KH}$ via FDT &
  $dk/dt$ kinematics & Equal-footing \\
\bottomrule
\end{tabular}
\end{center}

\noindent The convergence from multiple starting
points---two independent physical derivations (Onsager,
Padmanabhan), a kinematic restatement (surface gravity), and a
quantum-gravitational route
(CS/KMS)---constitutes strong evidence that
$I_{\rm eff} = (1-q)/2$ is a geometric property of the dynamical
apparent horizon, not an artefact of any single derivation
framework.

These routes share one structural element: the equal-footing
normalisation, i.e.\ that emergence and interaction exhaust the
horizon's total capacity with commensurate weighting.  We now show
that this is not an independent assumption but a consequence of the
canonical structure of the boundary theory.


\subsection{Canonical structure: confinement of \texorpdfstring{$k$}{k} and \texorpdfstring{$\gamma$}{gamma}}
\label{sec:canonical}

On the horizon worldtube $\Delta$, the boundary symplectic
structure of the Chern-Simons theory is
\begin{equation}
\Omega_\Delta \;\propto\;
  \frac{k}{4\pi}\int_\Delta \delta A \wedge \delta A\,,
\label{eq:symplectic}
\end{equation}
with $k = k(A_\Delta, \gamma)$ encoding all dependence on the
horizon area and the Barbero--Immirzi parameter.  The
junction-condition source coupling takes the form
$H_{\rm source} = \int_\Delta J_i\,\hat{A}^i$, which is linear in
the connection $A$.

The crucial observation is that the Chern-Simons kinetic term is
\emph{first-order} (topological): it takes the form
$A \wedge dA$, not $F \wedge \star F$.  Canonical normalisation
of $A$ therefore pulls down a single factor $\propto 1/k$ that
multiplies the source uniformly.  All $k$ and $\gamma$ dependence
is confined to a scalar prefactor:
\begin{align}
\beta_{\rm th}\,H_{\rm source} &=
  \eta \sum_p \widehat{\mathcal{O}}_{j_p}\,,
\nonumber \\
  \eta &= \eta(k,\gamma,\text{geometry},T_{\rm KH},\ldots)\,,
\label{eq:eta_prefactor}
\end{align}
where $\widehat{\mathcal{O}}_j$ is a dimensionless puncture
operator whose $j$-dependence is fixed by SU(2) representation
theory alone and $\beta_{\rm th}$ is the inverse temperature.

This confinement is specific to topological field theories.  In a
second-order theory such as Yang--Mills, the kinetic term
$F \wedge \star F$ is quadratic in $A$, so rescaling $A$ affects
the kinetic and source terms differently, generating
$k$-dependent ratios between representations.  In Chern-Simons
theory, the first-order kinetic term prevents this: $k$ enters
as an overall scale, and the relative weights of different
representations are set by group theory, not by the dynamics.


\subsection{Gauge invariance and Schur's lemma}
\label{sec:schur}

Physical horizon states satisfy the Gauss constraint, so all
observables must be gauge-invariant.  The gauge-invariant reduction
of any puncture operator
$\hat{O} \in \mathrm{End}(V_j)$ is obtained by group-averaging
over SU(2)---the \emph{SU(2) twirl}:
\begin{equation}
\mathcal{S}(\hat{O}) \;\equiv\;
  \int_{\mathrm{SU}(2)}\!dg\;
  U_j(g)\,\hat{O}\,U_j(g)^\dagger\,.
\label{eq:twirl_def}
\end{equation}
By construction, $\mathcal{S}(\hat{O})$ commutes with all
$U_j(g)$.  Schur's lemma~\cite{Hall2015} then requires it to be
proportional to the identity on $V_j$:
\begin{equation}
\mathcal{S}(\hat{O}) = c\,\mathbb{I}_j\,.
\label{eq:schur_prop}
\end{equation}
Taking the trace and using
$\mathrm{Tr}(U\hat{O}U^\dagger) = \mathrm{Tr}\,\hat{O}$ fixes
the proportionality constant:
\begin{equation}
\mathrm{Tr}\,\mathcal{S}(\hat{O})
  = \mathrm{Tr}\,\hat{O}
  = c\,d_j
\qquad\Longrightarrow\qquad
\boxed{\mathcal{S}(\hat{O})
  = \frac{\mathrm{Tr}\,\hat{O}}{d_j}\;\mathbb{I}_j\,.}
\label{eq:schur_exact}
\end{equation}
This is exact---no approximation enters.  It is a theorem of
representation theory, valid for any operator on any irreducible
representation of any compact group.

\paragraph{The rank-1 puncture operator.}

The ABCK boundary condition fixes a preferred internal direction
$\hat{n}$ (the horizon normal).  In the semiclassical
coherent-state sector, the source operator is dominated by
transitions through the coherent highest-weight state
$|j,\hat{n}\rangle$---the SU(2) rotation of $|j,j\rangle$ to
align with $\hat{n}$.  The minimal puncture operator consistent
with one active channel aligned with $\hat{n}$ is the rank-1
projector:
\begin{equation}
\hat{\Pi}_{j,\hat{n}} \equiv
  |j,\hat{n}\rangle\langle j,\hat{n}|\,,\qquad
  \mathrm{Tr}\,\hat{\Pi}_{j,\hat{n}} = 1\,.
\label{eq:rank1}
\end{equation}
The sourced puncture Hamiltonian then takes the form
\begin{equation}
\hat{H}_{\rm source}^{(j)}
  = \eta\,\lambda_j\,\hat{\Pi}_{j,\hat{n}}\,,
\label{eq:H_source}
\end{equation}
where $\eta$ is the common prefactor
from~\eqref{eq:eta_prefactor} and $\lambda_j$ is the kinematic
driving for each channel class:
\begin{itemize}
\item \textbf{Processing} ($j = 1$): The Kodama--Hayward surface
  gravity gives $T_{\rm KH} \propto (1-q)$; KMS linear response sets the
  transition rate $\lambda_1 = 1-q$ (after absorbing constants
  into $\eta$).
\item \textbf{Emergence} ($j = 1/2$): The CS level evolution
  $dk/d\ln a = 2(1+q)k$ gives the puncture creation rate
  $\lambda_{1/2} = 1+q$.
\end{itemize}

\paragraph{Explicit Schur reduction.}

Applying the twirl~\eqref{eq:schur_exact} to the sourced
Hamiltonian~\eqref{eq:H_source}:
\begin{equation}
\mathcal{S}\!\left(\hat{H}_{\rm source}^{(j)}\right)
  = \eta\,\lambda_j\,
    \mathcal{S}\!\left(\hat{\Pi}_{j,\hat{n}}\right)
  = \eta\,\lambda_j\,
    \frac{\mathrm{Tr}\,\hat{\Pi}_{j,\hat{n}}}{d_j}\,\mathbb{I}_j
  = \eta\,\frac{\lambda_j}{d_j}\,\mathbb{I}_j\,,
\label{eq:schur_result}
\end{equation}
since $\mathrm{Tr}\,\hat{\Pi}_{j,\hat{n}} = 1$ for any rank-1
projector, regardless of $j$ or $\hat{n}$.  Evaluating for each
channel class:
\begin{align}
\mathcal{S}\!\left(\hat{H}_{\rm source}^{(1/2)}\right)
  &= \eta\,\frac{1+q}{2}\,\mathbb{I}_{1/2}\,,
\nonumber \\
\mathcal{S}\!\left(\hat{H}_{\rm source}^{(1)}\right)
  &= \eta\,\frac{1-q}{3}\,\mathbb{I}_1\,.
\label{eq:schur_both}
\end{align}
The $d_j$ in the denominator arises purely from the Schur
reduction---it is a property of group-averaging over SU(2), not a
modelling choice.  The following lemma establishes the same result
from a weaker starting point.  Note the distinction: the
explicit projector result~\eqref{eq:schur_both} gives
$\eta\,\lambda_j/d_j\,\mathbb{I}_j$, which appears
$j$-dependent through $\lambda_j$.  The lemma below shows that
the \emph{operator itself} (before multiplication by the
kinematic driving $\lambda_j$) Schur-reduces to
$\eta_0\,\mathbb{I}_j$ with $\eta_0$ independent of $j$.  The
two results are consistent: $\lambda_j$ is external to the
operator algebra (set by KMS and the CS level evolution), and
$d_j$ from the partition function multiplicity cancels the $1/d_j$
from Schur reduction in the rate balance.

\paragraph{Lemma (linearity $\Rightarrow$ $j$-independent trace).}
Let the puncture coupling be sourced by a boundary term linear in
the SU(2) connection,
\begin{equation}
\hat{H}_{\rm source}
  = \int_\Delta J_i\,\hat{A}^i
  \quad\Longrightarrow\quad
  \hat{H}_{\rm source}^{(j)} : V_j \to V_j\,,
\label{eq:Hs_linear}
\end{equation}
restricted to a single puncture sector $V_j$.  SU(2) covariance
implies that, within a fixed irrep, any operator linear in
$\hat{A}^i$ decomposes into scalar $\oplus$ vector pieces:
\begin{equation}
\hat{H}_{\rm source}^{(j)}
  = \eta_0\,\mathbb{I}_j
    + \eta_1\,\hat{\mathbf{n}}\!\cdot\!\hat{\mathbf{J}}_j\,,
\label{eq:Hs_decomp}
\end{equation}
for some coefficients $\eta_{0,1}$ and a unit vector
$\hat{\mathbf{n}}$ fixed by the ABCK horizon data.  Taking the
trace on $V_j$:
\begin{equation}
\mathrm{Tr}_{V_j}\!\left(\hat{H}_{\rm source}^{(j)}\right)
  = \eta_0\,\mathrm{Tr}_{V_j}(\mathbb{I}_j)
    + \eta_1\,\hat{n}_i\,\mathrm{Tr}_{V_j}(\hat{J}^i_j)
  = \eta_0\,d_j\,,
\label{eq:trace_jindep}
\end{equation}
since $\mathrm{Tr}_{V_j}(\hat{J}^i_j) = 0$ for all $j$ (the
generators are traceless in every irreducible representation).
Therefore
\begin{equation}
\frac{\mathrm{Tr}_{V_j}\!\left(\hat{H}_{\rm source}^{(j)}\right)}
     {d_j}
  = \eta_0
  \quad\text{is independent of }j\,.
\label{eq:trace_over_d}
\end{equation}
SU(2) twirling (Schur averaging) then yields the universal
reduction:
\begin{equation}
\mathcal{S}\!\left(\hat{H}_{\rm source}^{(j)}\right)
  = \frac{\mathrm{Tr}_{V_j}\!\left(
      \hat{H}_{\rm source}^{(j)}\right)}{d_j}\,\mathbb{I}_j
  = \eta_0\,\mathbb{I}_j\,.
\label{eq:schur_linear_source}
\end{equation}
A representation-dependent contribution $\propto C_2(j)$ would
require a quadratic (or higher) operator---e.g.\
$\hat{\mathbf{J}}^2$ or $(J \!\cdot\! A)^2$---which is absent at
strictly linear order in the source.  This completes the proof:
\emph{any} source coupling that is linear in the connection yields
$j$-independent Schur-reduced weights, regardless of the specific
form of the puncture operator (rank-1 projector, thermal mixture,
or any intermediate case).

The rank-1 projector used in the explicit
evaluation~\eqref{eq:schur_both} is a special case: with
$\hat{\Pi}_{j,\hat{n}} = |j,\hat{n}\rangle\langle j,\hat{n}|$
one has $\eta_0 = 1/d_j$ and $\eta_1 \neq 0$ (the coherent state
breaks the isotropy), but the Schur-reduced result is identical
because the $\eta_1$ term is projected out by the twirl.  The
kinematic driving factors $\lambda_j$ are external to the operator
algebra (they come from the KMS transition rate and the
Chern-Simons level evolution), so the Gibbs weight per channel
class becomes $w_j = \eta_0\,\lambda_j\,d_j$, and the $d_j$ in the
multiplicity cancels $1/d_j$ from the Schur reduction exactly as
before.

The factors $1/d_j$ appearing in~\eqref{eq:schur_both} are the
same representation dimensions that enter the screening function
in Section~\ref{sec:algebraic}: $d_{1/2} = 2$ and $d_1 = 3$.
Their origin is now established by the lemma: they are the unique
output of gauge-invariant reduction for a linear source coupling,
anchored in the tracelessness of SU(2) generators in every
irreducible representation.

\paragraph{Controlled expansion and the first
representation-dependent correction.}

The lemma fixes the Schur-reduced source at \emph{strictly linear}
order in the boundary connection.  More generally, the
puncture-local coupling can be organised as an SU(2)-invariant
expansion:
\begin{equation}
\hat{H}_{\rm source}^{(j)}
  = \eta_0\,\mathbb{I}_j
    + \eta_1\,\hat{\mathbf{n}}\!\cdot\!\hat{\mathbf{J}}_j
    + \delta\hat{H}^{(j)}_{\rm quad} + \cdots\,,
\label{eq:Hs_expansion}
\end{equation}
where the first genuinely new invariant at quadratic order is
\begin{equation}
\delta\hat{H}^{(j)}_{\rm quad}
  = \epsilon_2\,\hat{\mathbf{J}}_j^{\,2}
  = \epsilon_2\,C_2(j)\,\mathbb{I}_j\,.
\label{eq:quad_correction}
\end{equation}
If present, this would deform the puncture weights to
\begin{equation}
w_j \;\propto\;
  d_j\,\exp\!\Bigl[\eta_0\,\lambda_j
    - \epsilon_2\,C_2(j)\Bigr]\,,
\label{eq:w_with_C2}
\end{equation}
making explicit that \emph{all} residual representation dependence
beyond $\lambda_j$ enters only through the coefficient
$\epsilon_2$.

\paragraph{Why $C_2(j)$ terms are $\mathcal{O}(1/k)$.}

On the horizon the boundary dynamics is Chern--Simons at level
$k$, so the only intrinsic, dimensionless control parameter of
the quantum boundary algebra is $1/k$.  In particular,
non-commutativity of the connection is set by the level:
\begin{align}
[\hat{A}^i_a(x),\,\hat{A}^j_b(y)]
  &= \frac{2\pi i}{k}\,\epsilon_{ab}\,\delta^{ij}\,
    \delta^{(2)}(x-y)
\nonumber \\
  &\quad(\text{up to conventions})\,,
\label{eq:AA_comm}
\end{align}
so any effective higher-order operator generated by
ordering or measure corrections must involve at least one
commutator and therefore carries at least one explicit factor of
$1/k$.  The first representation-sensitive invariant beyond the
linear source $\hat{\mathbf{n}}\!\cdot\!\hat{\mathbf{J}}_j$ is
the Casimir term~\eqref{eq:quad_correction}, with
\begin{equation}
\epsilon_2 = \mathcal{O}\!\left(\frac{1}{k}\right)\,.
\label{eq:eps2_from_comm}
\end{equation}
The level is set by the horizon area (ABCK and ENP normalisations
differ only by convention-dependent numerical factors involving
$\gamma$ and $4\pi$):
\begin{equation}
k = \kappa\,\frac{A}{\ell_P^2}\,,
  \qquad\Longrightarrow\qquad
  \frac{1}{k} = \frac{1}{\kappa}\,\frac{\ell_P^2}{A}
  = \mathcal{O}\!\left(\frac{\ell_P^2}{r_A^2}\right)\,,
\label{eq:eps2_bound}
\end{equation}
where $\kappa$ is a convention-dependent constant.
Therefore $\epsilon_2$ cannot shift the saddle allocation at any
observable level.  Appendix~\ref{app:source_reduction} derives the
operator structure (Casimir $C_2(j)$) and parametric scaling
($1/k$) from the second-order Brown--York expansion on the
perturbed horizon worldtube, with a candidate algebraic
coefficient $c_2 \sim 11/8$.  In this sense the $j$-independent trace
universality class is stable: linear terms renormalise
$\eta_0, \eta_1$ (absorbed into $\eta$), while the first new
representation-sensitive invariant appears only at
$\mathcal{O}(1/k) \sim 10^{-122}$.  The expansion makes precise
the claim that the lemma's ``strictly linear'' condition is not
fragile: it is the leading term of a controlled series whose first
correction is identified, bounded by the CS commutator
algebra~\eqref{eq:AA_comm}, and negligible for any macroscopic
horizon.


\subsection{Corrections to the linear-source regime}
\label{sec:linearity}

The preceding analysis identifies three sources of
representation-dependent corrections to the Schur-reduced weights,
each controlled by the same small parameter
$\ell_P^2/r_A^2 \sim 10^{-122}$.  We collect them here for
completeness.

\paragraph{Linear junction conditions.}

The linearity of the source coupling is guaranteed by the
Israel--Darmois junction conditions~\cite{Israel1966}, which
constrain the discontinuity in the extrinsic curvature linearly.
In LQG variables, the modified Gauss constraint
\begin{equation}
\frac{k}{2\pi}F^i + \Sigma^i = J^i
\label{eq:gauss_modified}
\end{equation}
is linear in the connection and triad.  Nonlinear
corrections---e.g.\ $(J \cdot A)^2$ terms---would arise only at
second order in the Israel--Darmois matching and are suppressed by
\begin{equation}
\left(\frac{\ell_P}{r_A}\right)^2
  \sim 10^{-122}
\label{eq:suppression}
\end{equation}
for the macroscopic apparent horizon.

\paragraph{Quantum-group corrections at finite $k$.}

At finite Chern-Simons level $k$, the SU(2) representation theory
is deformed to the quantum group SU(2)$_k$, with dimensions
replaced by $q$-dimensions:
\begin{equation}
[2j+1]_{\mathfrak{q}}
  = \frac{\sin\!\bigl(\pi(2j+1)/(k+2)\bigr)}
         {\sin\!\bigl(\pi/(k+2)\bigr)}\,.
\label{eq:q_dim}
\end{equation}
For $j = 0, \tfrac{1}{2}, 1$ and $k \sim 10^{122}/\gamma$, the
correction is $[d_j]_{\mathfrak{q}} = d_j\bigl(1 - \mathcal{O}(j^2/k^2)\bigr)$.
The $d_j$ cancellation in the Gibbs weight becomes approximate,
with a residual of order $j^2/k^2 \sim 10^{-244}$---even smaller
than the $\mathcal{O}(1/k) \sim 10^{-122}$ bound on nonlinear
junction-condition corrections.

\paragraph{Topological and perturbative stability.}

Beyond the numerical smallness of corrections, the linear-source
structure is perturbatively protected by the \emph{topological}
character of the Chern-Simons theory.  A topological field theory
has no local propagating degrees of freedom: there are no bulk
gluon modes and no loop diagrams that could generate
Casimir-dependent counterterms $\propto C_2(j)$ through
radiative corrections.  In Yang--Mills theory, gluon
self-interactions generate representation-dependent
renormalisation at each loop order; in Chern-Simons theory, the
absence of local excitations prevents this mechanism.

A careful objection notes that the punctures are defects in the
boundary theory, and one might worry that integrating over
puncture configurations---or regularising the measure on the
space of singular connections---could induce effective
$j$-dependence.  However, punctures are boundary data (specified
by the bulk spin-network state), not dynamical fields that
propagate and get integrated out in a path integral.  The
partition function~\eqref{eq:Z} already sums over all puncture
allocations $\{N_{1/2}, N_1\}$; no further integration remains
that could generate $C_2(j)$ terms.  Non-perturbative effects
(Chern-Simons instantons) are exponentially suppressed at
$k \sim 10^{122}$, contributing corrections of order
$e^{-k} \sim e^{-10^{122}}$.

The linear-source form is therefore not merely ``small corrections
suppressed by $10^{-122}$'': it is perturbatively stable to all
orders in the loop expansion and non-perturbatively protected by
the exponential suppression of instantons at macroscopic $k$.
The language ``structural protection'' is appropriate, with the
caveat that it rests on the standard treatment of punctures as
boundary data rather than dynamical defects---a framework that
could in principle be revised by a more complete formulation of
the quantum theory.

\paragraph{Summary of corrections.}

All three sources of representation-dependent corrections---nonlinear
junction terms, quantum-group deformation, and radiative
counterterms---are bounded by $\ell_P^2/r_A^2 \sim 10^{-122}$,
with the quantum-group corrections entering at the even smaller
scale $\mathcal{O}(1/k^2) \sim 10^{-244}$.  They enter the
puncture weights
through the coefficient $\epsilon_2$ of the controlled
expansion~\eqref{eq:quad_correction}, deforming the Gibbs weight
to~\eqref{eq:w_with_C2}.  This is the sharp falsifiability
criterion: the saddle-point allocation $N_1/N = (1-q)/2$ holds if
and only if $\epsilon_2 \cdot C_2(j) \ll 1$ for $j = 1/2, 1$,
i.e.\ if the sourced Chern-Simons theory remains in the
\emph{linear-source universality class}.  For the cosmological
apparent horizon, the bound~\eqref{eq:eps2_bound} ensures this to
$\mathcal{O}(10^{-122})$.


\subsection{Rate balance, partition function, and saddle-point
protection}
\label{sec:partition}

\paragraph{Primary derivation: rate balance.}

The allocation $N_1/N$ follows most directly from steady-state
rate balance among the puncture populations.  The explicit rank-1 projector
calculation~\eqref{eq:schur_result} gives the Schur-reduced
Hamiltonian $\eta\,(\lambda_j/d_j)\,\mathbb{I}_j$, so the
transition amplitude per magnetic substate in representation $j$
is $\eta\,\lambda_j/d_j$.  (The general lemma,
Eqs.~\eqref{eq:Hs_decomp}--\eqref{eq:schur_linear_source},
establishes that any linear source gives a $j$-independent
Schur weight $\eta_0$; the projector is the special case
$\eta_0 = 1/d_j$, so the per-substate amplitude is
$\eta\,\lambda_j/d_j$, not $\eta_0/d_j$ with $\eta_0$ generic.)
Each representation has $d_j$ magnetic substates, so the total
transition rate per puncture in representation $j$ is
\begin{equation}
\Gamma_j
  = d_j \cdot \frac{\eta_0}{d_j} \cdot \lambda_j
  = \eta_0\,\lambda_j\,,
\label{eq:rate_per_rep}
\end{equation}
where the $d_j$ from the substate count cancels the $1/d_j$ from
the Schur reduction \emph{exactly}, with no approximation.
The kinematic factors are $\lambda_1 = 1-q$ (the KMS transition
rate from Section~\ref{sec:kms_derivation}) and
$\lambda_{1/2} = 1+q$ (the Chern-Simons level evolution rate).
In steady state, the allocation is determined by the ratio of
rates:
\begin{equation}
\boxed{\frac{N_1}{N}\bigg|_{\rm steady\,state}
  = \frac{\Gamma_1}{\Gamma_1 + \Gamma_{1/2}}
  = \frac{1-q}{2}
  = I_{\rm eff}\,.}
\label{eq:saddle}
\end{equation}
(The intermediate step:
$\Gamma_1/(\Gamma_1 + \Gamma_{1/2})
= \eta_0(1{-}q)/[\eta_0(1{-}q) + \eta_0(1{+}q)]$.)
The prefactor $\eta_0$ cancels in the ratio.  This cancellation is
\emph{exact}: it requires no expansion, no classical limit, and no
assumption about the magnitude of $\eta_0$ relative to $k_B T_{\rm KH}$.
The $d_j$ cancellation likewise requires no approximation---it is
an algebraic consequence of the Schur lemma applied to the
linear source (the lemma of
Section~\ref{sec:schur}).

\paragraph{Complementary perspective: partition function.}

The same result can be recast in the language of statistical
mechanics.  The effective weight per puncture in representation
$j$, after the exact $d_j$ cancellation, is
\begin{equation}
w_j^{\rm eff} = \eta_0\,\lambda_j\,,
\label{eq:gibbs_weight}
\end{equation}
and the generating function over puncture configurations
$\{N_{1/2}, N_1\}$ with total $N = N_{1/2} + N_1$ is:
\begin{align}
Z &= \sum_{N_1=0}^{N} \binom{N}{N_1}
  \bigl[\eta_0\,(1+q)\bigr]^{N-N_1}
  \bigl[\eta_0\,(1-q)\bigr]^{N_1}
\nonumber \\
&= \eta_0^N\bigl[(1+q) + (1-q)\bigr]^N
= \eta_0^N \cdot 2^N\,.
\label{eq:Z}
\end{align}
This is a binomial generating function for the rate balance,
with weights $\lambda_{1/2} = 1+q$ and $\lambda_1 = 1-q$ (not
Boltzmann weights $e^{-\beta E_j}$; the ``saddle point'' is the
maximum of the rate allocation, not a free energy minimum).
Its saddle point reproduces~\eqref{eq:saddle}.

\paragraph{The $d_j$ cancellation does not require a
high-temperature expansion.}

One might worry that the $d_j$ cancellation relies on
linearising a Boltzmann factor.  If one writes the Gibbs weight
with an explicit multiplicity and chemical potential,
$w_j = d_j\,\exp(\mu_j/k_B T_{\rm KH})$ with
$\mu_j = \eta_0\,\lambda_j/d_j$, then the $d_j$ cancels
the $1/d_j$ in the exponent \emph{only after} expanding
$e^x \approx 1 + x$.  However, this formulation conflates two
distinct physical contributions.  The multiplicity $d_j$ counts
magnetic substates (a kinematic degeneracy), while $1/d_j$ from
the Schur reduction is the per-substate transition amplitude.
In the rate-balance formulation~\eqref{eq:rate_per_rep}, these
cancel at the level of the transition rate itself---before any
thermodynamic weighting.  The result~\eqref{eq:saddle} is
therefore not contingent on a high-temperature or weak-coupling
limit; it holds for \emph{any} value of $\eta_0/(k_B T_{\rm KH})$.

This is the central robustness result of the paper.  The
allocation does not depend on the absolute normalisation $\eta_0$.
Any renormalisation of the overall coupling strength---from quantum
corrections to $k$, shifts in $\gamma$, operator-ordering
ambiguities, coarse-graining, or changes in the horizon
temperature---modifies $\eta_0$ but cannot move the saddle point,
because $\eta_0$ enters as a common factor for both representations
and cancels in the ratio $N_1/N$.

The structural protection rests on four complementary properties:
\begin{enumerate}
\item \emph{First-order kinetic term.}  The Chern-Simons action
  $A \wedge dA$ is first-order in $A$, confining all $k$ and
  $\gamma$ dependence to the scalar prefactor
  $\eta$~\eqref{eq:eta_prefactor}.  In a second-order theory
  (Yang--Mills), the quadratic kinetic term would generate
  $k$-dependent representation ratios; the topological character
  of the boundary theory prevents this.

\item \emph{Schur reduction with $j$-independent trace.}
  Gauge invariance forces the SU(2)
  twirl~\eqref{eq:twirl_def}.  The linearity
  lemma~\eqref{eq:Hs_decomp}--\eqref{eq:schur_linear_source}
  guarantees that Schur reduction of any linear-in-$A$ source
  yields $\eta_0\,\mathbb{I}_j$ with $\eta_0$ independent of $j$,
  because SU(2) generators are traceless in every irrep.  A
  $C_2(j)$-dependent term would require a quadratic source,
  absent at linear order.

\item \emph{Linear junction conditions.}  The Israel--Darmois
  conditions~\cite{Israel1966} guarantee that the source coupling
  is linear in the connection.  Nonlinear corrections would
  introduce $C_2(j)$-dependent terms in the
  Gibbs weight~\eqref{eq:w_with_C2}, breaking the
  $\eta_0$-cancellation.  These are suppressed by
  $(\ell_P/r_A)^2 \sim 10^{-122}$---effectively zero for the
  cosmological horizon.

\item \emph{Exact $d_j$ cancellation and topological stability.}
  The rate-balance derivation~\eqref{eq:rate_per_rep} shows that
  the $d_j$ cancellation is exact---it occurs at the level of
  the transition rate, not through a high-temperature expansion.
  Quantum-group $\mathfrak{q}$-deformation of $d_j$ at finite
  $k$~\eqref{eq:q_dim} introduces corrections of order
  $j^2/k^2 \sim 10^{-244}$, and the topological nature of the
  Chern-Simons theory prevents radiative generation of
  Casimir-dependent counterterms at any loop order.
\end{enumerate}

Together, these four properties establish that the equal-footing
condition $c_e = c_i$ is not an assumption about the magnitude
of the bulk-boundary coupling.  It is a \emph{derived consequence}
of the sourced SU(2) Chern-Simons theory in the linear-source
regime (see Appendix~\ref{app:source_reduction} for the full
path-integral derivation):
the relative weight of the two channel populations is fixed by
kinematics ($\lambda_{1/2} = 1+q$, $\lambda_1 = 1-q$) and
group theory ($d_j$ cancellation), with all dynamical unknowns
absorbed into the universal prefactor $\eta \propto 1/k$.

The following subsubsection carries out this programme explicitly.


\subsubsection{Canonical quantisation and path-integral
verification}
\label{sec:path_integral_proof}

The symmetry arguments and Schur reduction above can be verified
directly from the canonical path integral of sourced SU(2)$_k$
Chern--Simons theory on the horizon.  The strategy is to compute
the \emph{transition rate} per puncture via linear rate weighting
(not a literal Fermi's golden rule, which would be quadratic in
the matrix element; see the note in the rate-balance derivation
above) rather than a Boltzmann weight, thereby obtaining the
allocation from rate balance without any exponential or
high-temperature expansion.

\paragraph{1.\ Sourced CS path integral and Gauss constraint.}

Consider SU(2)$_k$ Chern--Simons theory on
$\Delta \simeq S^2 \times \mathbb{R}$ with $N$ punctures at
$\{z_p\}$, described by coadjoint-orbit variables $g_p(t)$ in
irreps $V_{j_p}$.  A macroscopic, homogeneous (zero-mode)
Israel--Darmois current $J_i$ couples linearly to the connection:
\begin{align}
Z[J] &= \int \mathcal{D}A\;\prod_{p=1}^{N}\mathcal{D}g_p\;
  \exp\!\biggl\{iS_{\rm CS}[A]
\nonumber \\
  &\quad + i\sum_{p=1}^{N} S_{\rm WZ}[g_p, A]
    + i\!\int_\Delta J_i\,A^i\biggr\}.
\label{eq:ZJ_def}
\end{align}
In canonical variables, integrating out $A_0$ enforces the
(distributional) Gauss law on $S^2$:
\begin{equation}
\frac{k}{2\pi}\,F_{12}^i(z)
  + \sum_{p=1}^{N} \mathcal{J}^i_{(p)}\,
    \delta^{(2)}(z - z_p)
  = J^i(z)\,,
\label{eq:Gauss_source}
\end{equation}
where
$\mathcal{J}^i_{(p)} \equiv
  \mathrm{Tr}\!\bigl(\lambda_p\,g_p^{-1}\tau^i g_p\bigr)$
are the puncture charges.

\paragraph{2.\ Single-puncture transition amplitude.}

The dark drag is a continuous macroscopic flux, not a thermal
equilibrium.  The relevant observable is therefore the
\emph{transition rate} at which each puncture mediates energy
transfer, computed via a linear rate-weight argument (proportional to
the matrix element of the source Hamiltonian, not its square;
the quadratic $|M|^2$ of a literal Fermi's golden rule would
give an allocation $\propto (1-q)^2/[(1-q)^2+(1+q)^2]$ rather
than $(1-q)/2$---the linear form is appropriate here because
$\eta_0 \ll k_BT_{\rm KH}$ places the system in the weak-driving,
linear-response regime).  By SU(2)
covariance and the linearity of the junction coupling, the most
general puncture-local source operator at leading order is
(cf.\ the lemma in Section~\ref{sec:schur}):
\begin{equation}
\hat{H}^{(j)}_{\rm source}
  = \eta_0\,\lambda_j\,\mathbb{I}_j
    + \eta_1\,\hat{\mathbf{n}}\!\cdot\!\hat{\mathbf{J}}_j\,,
\label{eq:Hsource_PI}
\end{equation}
where $\hat{\mathbf{n}}$ is fixed by the macroscopic boundary
data.

\paragraph{3.\ Global SU(2) projection (Schur reduction).}

The physical horizon state is invariant under rigid SU(2)
rotations.  The single-puncture source operator must therefore be
averaged over the gauge group:
\begin{equation}
\mathcal{S}\!\left(\hat{H}^{(j)}_{\rm source}\right)
  = \int_{\mathrm{SU}(2)} dg\;
    U_j(g)\,\hat{H}^{(j)}_{\rm source}\,U_j(g)^{-1}\,.
\label{eq:twirl_PI}
\end{equation}
By Schur's lemma, the traceless vector piece
$\eta_1\,\hat{\mathbf{n}}\!\cdot\!\hat{\mathbf{J}}_j$ integrates
to zero (it transforms as a spin-$1$ tensor under the twirl and
has no singlet component), while the scalar piece
$\eta_0\,\lambda_j\,\mathbb{I}_j$ is invariant.  The result is
exact and non-perturbative:
\begin{equation}
\mathcal{S}\!\left(\hat{H}^{(j)}_{\rm source}\right)
  = \eta_0\,\lambda_j\,\mathbb{I}_j\,.
\label{eq:schur_PI_result}
\end{equation}
No trace of the vector term survives; the $\eta_1$ piece is
projected out \emph{before} any exponentiation.  This is the
critical difference from computing the trace of the exponentiated
operator: the SU(2) character
$\chi_j(\eta_1) = \mathrm{Tr}_{V_j}(e^{\eta_1\,
\hat{\mathbf{n}}\cdot\hat{\mathbf{J}}_j})$ is $j$-dependent, but
the twirl acts on the operator itself, not on its exponential.

\paragraph{4.\ Transition rate and the $d_j$ cancellation.}

The Schur-reduced source~\eqref{eq:schur_PI_result} is
$\eta_0\,\lambda_j$ times the identity on $V_j$.  Every magnetic
substate $|j,m\rangle$ therefore receives the \emph{same} scalar
coupling $\eta_0\,\lambda_j$.  Note the bookkeeping convention
here: the $d_j$ cancellation is absorbed into the per-substate
coupling ($\mathrm{Tr}(\hat{H}^{(j)}_{\rm source}) =
\eta_0\,\lambda_j\,d_j$, divided by $d_j$ substates), so
$\eta_0\,\lambda_j$ is the per-substate rate directly.  This
differs from the rate-balance bookkeeping
of Section~\ref{sec:partition}, where the $d_j$ cancellation
is carried explicitly as a separate factor; both conventions
yield the same total rate $\eta_0\,\lambda_j\,d_j$ per puncture
and the same saddle-point allocation~\eqref{eq:saddle}.

In the gauge-invariant sector, the SU(2) twirl erases the
magnetic quantum number $m$: all substates are physically
equivalent.  The transition rate for a puncture in representation
$j$ is therefore determined entirely by the scalar coupling:
\begin{equation}
\Gamma_j^{\rm phys} = \eta_0\,\lambda_j\,.
\label{eq:rate_physical}
\end{equation}
This reproduces equation~\eqref{eq:rate_per_rep} of the main text
from the path-integral framework.  The $d_j$ does not appear in
the physical rate because it has already been absorbed by the
Schur reduction---it is not cancelled by an approximation but by
exact gauge averaging.

\paragraph{5.\ Saddle allocation from rate balance.}

In steady state, the fraction of punctures in the adjoint channel
is determined by the ratio of physical rates:
\begin{align}
\frac{N_1}{N}\bigg|_{\rm saddle}
  &= \frac{\Gamma_1^{\rm phys}}
         {\Gamma_1^{\rm phys} + \Gamma_{1/2}^{\rm phys}}
\nonumber \\
  &= \frac{\eta_0\,(1-q)}
         {\eta_0\,(1-q) + \eta_0\,(1+q)}
  = \frac{1-q}{2}\,.
\label{eq:saddle_PI}
\end{align}
The prefactor $\eta_0$ cancels in the ratio.  This cancellation is
exact: it requires no linearisation of an exponential, no
high-temperature limit, and no assumption about the magnitude of
$\eta_0$ relative to $k_B T_{\rm KH}$.  The allocation depends only on
the kinematic factors $\lambda_1 = 1-q$ and
$\lambda_{1/2} = 1+q$, which are determined by the horizon's
expansion history.

This completes the path-integral verification: the $d_j$
cancellation is enforced by the exact SU(2) Haar twirl (Schur
reduction) acting on the \emph{operator} before exponentiation,
and the saddle allocation is fixed by the linear rate balance of
the kinematic factors $\lambda_j$, immune to microscopic
regularisation artefacts.

\paragraph{6.\ Regularisation, operator ordering, and the
one-loop effective action.}

While the Schur reduction guarantees the exact cancellation of
$d_j$ at the operator level, one must verify that no
regularisation of the puncture Hilbert space generates quantum
anomalies that evade this suppression.

In the canonical quantisation of Chern-Simons theory, the level
$k$ acts as the inverse Planck constant, dictating the
non-commutativity of the connection on the spatial boundary
(cf.~\eqref{eq:AA_comm}).  Because the commutator carries an
explicit factor of $1/k$, any operator-ordering ambiguity in the
quantisation of $\int J_i\,\hat{A}^i$ inherits at least one
factor of $1/k$.  Zeroth-order quantum corrections to the algebra
are forbidden.

To track how representation-dependent corrections emerge, we
examine two sources: the framing anomaly and the quantum-group
state counting.

\emph{Framing anomaly.}  Integrating out gauge fluctuations
$\tilde{A}$ around the zero-mode background $\bar{A}$
introduces a framing phase to the holonomy expectation values.
The regularised Wilson line evaluates
to~(see, e.g.,~\cite{Hall2015} for the representation-theoretic
context):
\begin{equation}
\langle \mathcal{W}_j \rangle_{\rm reg}
  = \langle \mathcal{W}_j \rangle_{\rm cl}\;
    \exp\!\left(\frac{2\pi i\,C_2(j)}{k+2}\right)\,,
\label{eq:framing_anomaly}
\end{equation}
where $C_2(j) = j(j+1)$ is the quadratic Casimir and the shift
$k \to k+2$ reflects the dual Coxeter number of SU(2).  Crucially,
this is a \emph{pure phase}: $|\exp(2\pi i\,C_2/(k+2))|^2 = 1$.
In the rate-balance framework used in this paper, transition rates
involve $|\text{matrix element}|^2$, so the framing phase cancels
identically.  (In a Euclidean partition-function framework, the
Wick rotation makes this factor real---$\exp(-2\pi C_2(j)/(k+2))$
---contributing a genuine $\mathcal{O}(1/k)$ correction to the
Boltzmann weights.  This is consistent with the $\epsilon_2$
bound of~\eqref{eq:eps2_bound}, but the rate-balance derivation
is immune to it.)

\emph{Quantum-group state counting.}  The exact,
non-perturbative regularisation of the puncture Hilbert space is
given by the quantum group SU(2)$_{\mathfrak{q}}$ with deformation
parameter $\mathfrak{q} = \exp(i\pi/(k+2))$ (distinct from the
deceleration parameter $q$ used throughout the paper).  The
regularised density of states is the $\mathfrak{q}$-dimension~\eqref{eq:q_dim}:
\begin{equation}
[d_j]_{\mathfrak{q}}
  = \frac{\sin\!\bigl(\pi d_j/(k+2)\bigr)}
         {\sin\!\bigl(\pi/(k+2)\bigr)}\,.
\end{equation}
Taylor expanding:
\begin{equation}
[d_j]_{\mathfrak{q}}
  \approx d_j\left[
    1 - \frac{\pi^2(d_j^2 - 1)}{6(k+2)^2}\right]
  = d_j\left(
    1 - \frac{2\pi^2\,C_2(j)}{3(k+2)^2}\right)\,,
\label{eq:q_dim_expand}
\end{equation}
using $d_j^2 - 1 = (2j+1)^2 - 1 = 4j(j+1) = 4C_2(j)$.  The
first representation-dependent correction to the state counting
enters at $\mathcal{O}(1/k^2) \sim 10^{-244}$ for the
cosmological horizon---two orders of magnitude below even the
$\mathcal{O}(1/k) \sim 10^{-122}$ bound from operator ordering.
Non-perturbative Chern-Simons instantons contribute at
$\mathcal{O}(e^{-k})$, which is more severely suppressed still.

In summary, three classes of quantum corrections have been
identified and bounded:
\begin{center}
\small
\begin{tabular}{lcc}
\toprule
Source & Scaling & Cosmological value \\
\midrule
Operator ordering & $\mathcal{O}(1/k)$ & $\sim 10^{-122}$ \\
Framing anomaly (rate) & exact $0$ & $0$ \\
Framing anomaly (Euclidean) & $\mathcal{O}(1/k)$ & $\sim 10^{-122}$ \\
$q$-dimension & $\mathcal{O}(1/k^2)$ & $\sim 10^{-244}$ \\
CS instantons & $\mathcal{O}(e^{-k})$ & $\sim e^{-10^{122}}$ \\
\bottomrule
\end{tabular}
\end{center}
\noindent The rate-balance derivation, which is the primary result
of this paper, is immune to the framing anomaly entirely.  The
leading correction to the saddle allocation $I_{\rm eff} = (1-q)/2$
comes from operator ordering at $\mathcal{O}(1/k) \sim 10^{-122}$.
The Schur-reduced $d_j$ cancellation is robust to all orders in
the quantum theory.


% ======================================================================
%  SECTION 6: DISCUSSION
% ======================================================================


% ======================================================================
%  OVER-CAPACITY WRITE SECTOR
% ======================================================================
\section{Over-Capacity Write Sector and Protected Horizon Memory}
\label{sec:memory_write}

\subsection{Hard overflow and the capacity-cycle clock}

Define the dimensionless drive ratio
\begin{equation}
\mathcal S\equiv \frac{\dot Q_{\rm drive}}{\dot Q_{\rm cap}}.
\label{eq:pb_S}
\end{equation}
When $\mathcal S>1$, one capacity unit is processed and the rejected
fraction of the attempted throughput is
\begin{equation}
\Delta_{\rm hard}(\mathcal S)
=\Theta(\mathcal S-1)\left(1-\frac{1}{\mathcal S}\right).
\label{eq:pb_Delta}
\end{equation}
The hyperbolic-tangent gate used in numerical work is a regulator of this
hard definition, not a new physical function.

Let $u^\mu$ be the selected homogeneous cosmological congruence,
$\Theta=\nabla_\mu u^\mu$, and $H_u=\Theta/3$.  We define the scalar
capacity-cycle coordinate by
\begin{equation}
u^\mu\nabla_\mu\lambda_H=H_u,
\qquad
h^\mu{}_{\nu}\nabla^\nu\lambda_H=0,
\label{eq:pb_lambda}
\end{equation}
where $h^\mu{}_{\nu}=\delta^\mu{}_{\nu}+u^\mu u_\nu$.  On FLRW,
$d\lambda_H=Hdt=d\ln a$.  The choice is physically motivated by
\begin{equation}
E_{\rm MS}=Nk_BT_A,
\qquad
\frac{\dot E_{\rm cap}}{E_{\rm MS}}=H,
\label{eq:pb_capacity_clock}
\end{equation}
so a geometric-capacity cycle occurs once per Hubble time.  The
Pendry-limited interaction clock is instead
$\dot E_{\rm QT}/E_{\rm MS}=\beta H=H/12$; it controls the processed
interaction amplitude, not the normalization of the excess history.

\subsection{The additive history functional}

For any interval on the homogeneous branch, define
\begin{equation}
X_H[a_1,a_2]
=\int_{a_1}^{a_2}\Delta[\mathcal S(a)]\,d\lambda_H.
\label{eq:pb_XH}
\end{equation}
This positive functional is additive:
$X_H[a_1,a_3]=X_H[a_1,a_2]+X_H[a_2,a_3]$.  It is a record of the
history, not yet a claim that a physical carrier retains that record.

For flat matter+$\Lambda$ with the hard gate,
$\mathcal S=(9/2)\Omega_\Lambda(1-\Omega_\Lambda)$ and the roots are
$\Omega_\Lambda=1/3,2/3$.  Since
$d\ln a=d\Omega_\Lambda/[3\Omega_\Lambda(1-\Omega_\Lambda)]$,
\begin{align}
X_\Lambda^{\rm hard}
&=\int_{1/3}^{2/3}
\left(1-\frac{2}{9\Omega_\Lambda(1-\Omega_\Lambda)}\right)
\frac{d\Omega_\Lambda}{3\Omega_\Lambda(1-\Omega_\Lambda)}
\notag\\
&=\frac{10}{27}\ln2-\frac{2}{9}
=0.0344989558.
\label{eq:pb_X_closed}
\end{align}
The companion interacting background yields a nearby sample-dependent
value, approximately $X_0\simeq0.0359$.  The exact number in
Eq.~\eqref{eq:pb_X_closed} belongs to the flat hard-gate benchmark; it
must not be confused with the interacting posterior.

\subsection{Capacity slots and the processed/excess instrument}
\label{sec:pb_write_instrument}

The source decomposition can be sharpened into a microscopic counting
problem.  The energy scale per horizon channel is fixed by the exact
identity
\begin{equation}
\varepsilon_A\equiv\frac{E_{\rm MS}}{N}=k_BT_A,
\label{eq:pb_channel_energy}
\end{equation}
where $N=A/(4\ell_P^2)$.  Since
$\dot Q_{\rm cap}=HE_{\rm MS}$, the number of normalized
thermal-energy-equivalent capacity slots in proper time $dt$ is
\begin{equation}
d\nu_{\rm cap}
\equiv\frac{\dot Q_{\rm cap}\,dt}{\varepsilon_A}
=N Hdt=N\,d\lambda_H.
\label{eq:pb_capacity_slots}
\end{equation}
Thus the geometric-capacity clock is not chosen after the fact: one unit of
$d\lambda_H$ exposes each of the $N$ boundary channels to one normalized
capacity opportunity.  This statement concerns the geometric benchmark;
the Pendry factor $\beta$ controls how much of the processed branch is
converted into the dark-sector interaction.

The same interval contains
\begin{equation}
d\nu_{\rm att}=\mathcal S\,N\,d\lambda_H
\label{eq:pb_attempted_slots}
\end{equation}
source attempts.  Saturation permits
\begin{align}
d\nu_{\rm proc}&=\min(\mathcal S,1)N\,d\lambda_H,\notag\\
d\nu_{\rm exc}&=\max(\mathcal S-1,0)N\,d\lambda_H.
\label{eq:pb_slot_split}
\end{align}
Consequently a randomly selected attempted update carries the excess label
with probability
\begin{equation}
p_{\rm exc}
=\frac{d\nu_{\rm exc}}{d\nu_{\rm att}}
=\Theta(\mathcal S-1)\left(1-\frac1{\mathcal S}\right)
=\Delta(\mathcal S).
\label{eq:pb_rejection_probability}
\end{equation}
This is the trajectory-level form of the Kubo/Pendry processed--excess split.
Counting fields and quantum-jump trajectory ensembles provide the standard
microscopic language for such two-outcome open-system records
\cite{Esposito2009,Garrahan2010}.

Choose a coarse-graining cell $\delta\lambda$ satisfying
$N^{-1}\ll\delta\lambda\ll1$, so the background is adiabatically constant
while the cell contains
$M_j=\mathcal S_jN_j\delta\lambda\gg1$ attempted updates.  The previously
derived internal relaxation rate $\Gamma_{\rm int}/H\sim k\gg1$ makes
successive cells Markovian up to finite-$k$ corrections.  After a
Stinespring dilation of the two-outcome instrument, let
$R_{j\alpha}$ be the projector onto the excess outcome of attempt $\alpha$,
with
\begin{equation}
\langle R_{j\alpha}\rangle=\Delta_j,
\qquad
\overline R_j\equiv\frac1{M_j}\sum_{\alpha=1}^{M_j}R_{j\alpha}.
\label{eq:pb_rejection_projector}
\end{equation}
The source is homogeneous, so the boundary monopole couples to the
\emph{intensive} average $\overline R_j$, not to the extensive raw count.
This monopole projection is the microscopic reason the factors $N$ and $k$
do not multiply the mean memory.

\subsection{Controlled translation of the memory zero mode}
\label{sec:pb_controlled_write}

Let $(\widehat\chi,\widehat p_\chi)$ be the homogeneous shift mode and its
canonical generator,
$[\widehat\chi,\widehat p_\chi]=i$.  The minimal local, shift-covariant
recording unitary for one cell is
\begin{equation}
U_j
=\exp[-i\,\delta\lambda\,
       \overline R_j\otimes\widehat p_\chi].
\label{eq:pb_write_unitary}
\end{equation}
It is a controlled translation: because all higher nested commutators
vanish,
\begin{equation}
U_j^\dagger\widehat\chi U_j
=\widehat\chi+\delta\lambda\,\overline R_j
\label{eq:pb_write_shift}
\end{equation}
exactly, without expanding in the memory amplitude.  Taking the expectation
value in the driven boundary state gives
\begin{equation}
\frac{\delta\langle\chi\rangle}{\delta\lambda}
=\langle\overline R_j\rangle
=\Delta_j.
\label{eq:pb_write_mean_discrete}
\end{equation}
The unit coefficient is fixed by two canonical normalizations: the excess
label has eigenvalue one, and $p_\chi$ generates unit translations of the
normalized homogeneous mode.  A formal rescaling of $\chi$ would move this
normalization into every endpoint gate; the present coordinate is the
physical counting normalization in which a fully rejected capacity cycle
advances the memory by one unit.

The full counting statistics makes the deterministic limit explicit.  For
independent attempts at leading order in the source-reduced saddle, the
characteristic function of the cell increment
$\delta\chi_j=\delta\lambda\,\overline R_j$ is
\begin{equation}
Z_j(\xi)
=\left[1-\Delta_j+
\Delta_j e^{i\xi\delta\lambda/M_j}\right]^{M_j}.
\label{eq:pb_write_fcs}
\end{equation}
For $M_j\gg1$,
\begin{equation}
\ln Z_j
=i\xi\Delta_j\delta\lambda
-\frac{\xi^2}{2}
 \frac{\Delta_j(1-\Delta_j)}{\mathcal S_jN_j}
 \delta\lambda
+\mathcal O(M_j^{-2}).
\label{eq:pb_write_fcs_expand}
\end{equation}
Thus the continuum process obeys
\begin{align}
\frac{d\langle\chi\rangle}{d\lambda_H}
&=\Delta(\mathcal S),
\label{eq:pb_write_law}\\
\mathrm{Var}(\chi_0)
&=\int_{\mathcal S>1}
F_{\rm exc}\,
\frac{\Delta(1-\Delta)}{\mathcal S N}\,d\lambda_H.
\label{eq:pb_write_variance}
\end{align}
Here
\begin{equation}
F_{\rm exc}=1+\mathcal O(k^{-1})
\label{eq:pb_write_fano}
\end{equation}
Here $F_{\rm exc}$ is the excess-channel Fano factor; the independent-channel
saddle has $F_{\rm exc}=1$, while connected WZW-current correlations provide
the displayed finite-$k$ correction.  Since $N\sim10^{122}$ throughout the late window,
$\sigma_\chi\lesssim\sqrt{X_H/N_{\min}}\ll10^{-60}$.  The memory is
therefore a sharply peaked classical record even though its microscopic
origin is a quantum trajectory ensemble.

\begin{proposition}[Capacity-cycle write law]
Within the source-reduced Markovian channel saddle, with the processed and
excess outcomes distinguished by the saturation instrument and the
homogeneous memory coupled to their monopole average, the final protected
memory satisfies
\begin{equation}
\boxed{\langle\chi_0\rangle=X_H
=\int\Delta(\mathcal S)\,d\lambda_H.}
\label{eq:pb_write_theorem}
\end{equation}
The stochastic width is $\mathcal O(N^{-1/2})$ and vanishes in the
semiclassical horizon limit.
\end{proposition}

\subsection{Large-$N$ influence action and protected carrier}
\label{sec:pb_write_effective_action}

The cumulant functional in Eq.~\eqref{eq:pb_write_fcs_expand} has the
large-$N$ saddle generated by the first-order action
\begin{equation}
S_{0}[\chi,p_\chi]
=\int d\lambda_H\,p_\chi
\left(\frac{d\chi}{d\lambda_H}-\Delta\right).
\label{eq:pb_counting_action}
\end{equation}
The second cumulant adds a Gaussian Schwinger--Keldysh noise kernel
proportional to
$\Delta(1-\Delta)/(\mathcal S N)$; connected channel corrections enter at
the same finite-size order controlled by $1/k$.  The covariant lift of the
leading saddle is precisely
\begin{equation}
S_{\rm mem}
=\int_{\mathcal M}C_3\wedge
\left[d\chi-\Delta(\mathcal S)d\lambda_H\right].
\label{eq:pb_memory_action}
\end{equation}
Varying $C_3$ yields $d\chi=\Delta d\lambda_H$.  The action is invariant
under $\chi\rightarrow\chi+\chi_c$ and contains no local propagating bulk
mode.  Within this EFT the shift symmetry forbids a local potential
$V(\chi)$, protecting the zero mode from ordinary potential-driven
relaxation.  It does not exclude every ultraviolet or environmental
derivative coupling; an explicit erasure channel would modify the transport
law to
\begin{equation}
u^\mu\nabla_\mu\chi
=H_u\Delta(\mathcal S)-\Gamma_{\rm er}\chi.
\label{eq:pb_memory_transport}
\end{equation}
The protected branch has $\Gamma_{\rm er}=0$.

The derivation also explains why no extra factor appears.  The Pendry
coefficient $\beta$ belongs to the processed interaction branch, whereas the
write instrument counts the normalized excess label.  The factor $1+q$
controls $dN/d\lambda_H=2(1+q)N$, but cancels between the extensive excess
count and the attempted count in $\overline R$; weighting by newly created
area would define the distinct historical accumulator $\chi_A$.  Finally,
$k$ fixes the number and correlations of microscopic channels, but the
monopole average removes it from the mean, leaving only finite-$k$
fluctuation corrections.  Consequently the capacity-cycle clock, rather
than the Pendry or area-growth clock, is selected by the normalized
trajectory observable itself.



% ======================================================================
%  ENDPOINT THEOREM AND COMMON GATE
% ======================================================================

% ======================================================================
%  SOURCE-REDUCED HAMILTONIAN COMPLETION
% ======================================================================
\section{Source-Reduced Hamiltonian Dilation of the Write Instrument}
\label{sec:exact_hamiltonian_dilation}

The counting argument above fixes the drift and its large-$N$ width.  We now
exhibit a Hamiltonian dilation whose reduced dynamics gives the same
processed/excess instrument.  The construction is exact within the
source-reduced, adiabatic CS/WZW boundary theory.  It is not a claim that a
unique non-perturbative quantum-gravity Hamiltonian has been reconstructed.

\subsection{Hilbert bundle and level-changing boundary transition}

Because the Chern--Simons level changes with the apparent-horizon area, the
instantaneous boundary spaces are not literally one fixed finite-dimensional
Hilbert space.  We use the direct-sum dilation
\begin{equation}
 \mathcal H_{\rm B}=\bigoplus_{k\in\mathbb N}\mathcal H_k^{\rm WZW},
 \qquad
 H_{\rm B}=\bigoplus_k P_k H_{\rm WZW}(k)P_k,
 \label{eq:pb_hamiltonian_bundle}
\end{equation}
where $P_k$ projects onto the level-$k$ sector.  The slow geometric growth is
implemented by an isometry
\begin{equation}
 V_{k+\delta k,k}:\mathcal H_k^{\rm WZW}
\rightarrow
 \mathcal H_{k+\delta k}^{\rm WZW},
 \qquad V_{k+\delta k,k}^{\dagger}V_{k+\delta k,k}=P_k.
 \label{eq:pb_level_isometry}
\end{equation}
The relevant adiabatic parameter is the fractional, rather than the
absolute, change of the level per attempted update.  Since
$dk/d\lambda_H=2(1+q)k$ and there are $N\mathcal S$ attempts per
capacity-cycle unit,
\begin{equation}
 \frac{\delta k_{\rm att}}{k}
 =\frac{1}{k}\frac{dk/d\lambda_H}{N\mathcal S}
 =\frac{2(1+q)}{N\mathcal S}\ll1.
 \label{eq:pb_level_adiabaticity}
\end{equation}
Using $k=N/(\pi\gamma)$, the corresponding mean absolute increment is
\begin{equation}
 \delta k_{\rm att}
 =\frac{2(1+q)}{\pi\gamma\mathcal S},
 \label{eq:pb_level_absolute_increment}
\end{equation}
and need not be smaller than unity.  The instantaneous current algebra is
nevertheless adiabatic because its relative change is
$\mathcal O(N^{-1})$.  At exact integer level the isometry is understood as
the nearest allowed level transition, or as an ensemble of such transitions
whose mean is Eq.~\eqref{eq:pb_level_absolute_increment}.

The source-selected transverse current operators are whitened with their KMS
Gram matrix,
\begin{align}
 G_{rs}(k)&\equiv
 \frac{1}{k}\left\langle J^{r}_{1}J^{s}_{-1}\right\rangle_{\rm KMS},
 \qquad r,s\in\{+,-\},
 \label{eq:pb_current_gram}\\
 B_{r,k}&\equiv
 \frac{1}{\sqrt{k}}\sum_s[G(k)^{-1/2}]_{rs}J^s_{-1},
 \label{eq:pb_whitened_current}\\
 A_{r,k}&\equiv V_{k+\delta k,k}B_{r,k}.
 \label{eq:pb_level_transition}
\end{align}
On the driven source subspace the normalization may be chosen so that
$\sum_r\pi_r A_{r,k}^{\dagger}A_{r,k}=P_k$, with
$\pi_++\pi_-=1$.  This whitening is useful beyond the independent-channel
limit: connected WZW correlations change $G_{rs}$ and the noise kernel, but
not the normalized monopole drift.

\subsection{One-attempt collision unitary and Kraus operators}

Introduce an outcome pointer qutrit
\begin{equation}
 \mathcal H_Q=\operatorname{span}\{|0\rangle,|p\rangle,|e\rangle\},
 \label{eq:pb_outcome_qutrit}
\end{equation}
for ready, processed, and excess.  The transverse-current label
$r\in\{+,-\}$ is carried by a separate two-state register.  In a cell where
$\mathcal S$ is adiabatically constant, let
\begin{equation}
 \theta_{\mathcal S}\equiv\arcsin\sqrt{\Delta(\mathcal S)},
 \qquad
 \epsilon_{\mathcal S}\equiv\frac{1}{N\mathcal S},
 \label{eq:pb_collision_parameters}
\end{equation}
and define the pointer translation
$T_{\epsilon}=\exp[-i\epsilon\widehat p_\chi]$.  Let $S_{0p}$ be a
unitary ready-to-processed swap, $S_{0p}|0\rangle=|p\rangle$, and define
\begin{equation}
 Y_{pe}\equiv-i\bigl(|p\rangle\langle e|-|e\rangle\langle p|\bigr),
 \qquad
 R_{pe}(\theta)\equiv e^{-i\theta Y_{pe}},
 \label{eq:pb_qutrit_rotation}
\end{equation}
where $Y_{pe}$ acts only on
$\operatorname{span}\{|p\rangle,|e\rangle\}$.  Then
\begin{equation}
 R_{pe}(\theta_{\mathcal S})S_{0p}|0\rangle
 =\sqrt{1-\Delta}\,|p\rangle+\sqrt{\Delta}\,|e\rangle.
 \label{eq:pb_qutrit_preparation}
\end{equation}
This qutrit preparation, followed by an excess-controlled pointer
translation and the level-changing current transition, gives the isometry
\begin{multline}
 \mathcal V_{\rm att}|\psi_k\rangle|0\rangle
 =\sum_{r=\pm}\sqrt{\pi_r}\bigg[
 \sqrt{1-\Delta}\,A_{r,k}|\psi_k\rangle|p,r\rangle\\
 +\sqrt{\Delta}\,T_{\epsilon_{\mathcal S}}
 A_{r,k}|\psi_k\rangle|e,r\rangle\bigg].
 \label{eq:pb_attempt_isometry}
\end{multline}
The normalization above makes
$\mathcal V_{\rm att}^{\dagger}\mathcal V_{\rm att}=P_k$.  By Stinespring's
theorem the isometry extends to a unitary on the system plus unused ancilla
states.  An explicit collision Hamiltonian is obtained by composing the
four elementary kicks,
\begin{equation}
 U_{\rm att}
 =U_A\,\exp[-i|e\rangle\langle e|\otimes
 \epsilon_{\mathcal S}\widehat p_\chi]\,
 R_{pe}(\theta_{\mathcal S})S_{0p},
 \label{eq:pb_collision_unitary}
\end{equation}
where $U_A$ is a unitary dilation of the isometries $A_{r,k}$ together with
the transverse-current label register.  Equivalently,
$U_{\rm att}=\exp[-iH_{\rm att}\tau_{\rm att}]$ defines a bounded collision
Hamiltonian on the enlarged space.

Tracing the outcome record gives the Kraus operators
\begin{align}
 K_{p,r}&=\sqrt{\pi_r(1-\Delta)}\,A_{r,k},
 \label{eq:pb_processed_kraus}\\
 K_{e,r}&=\sqrt{\pi_r\Delta}\,
 T_{\epsilon_{\mathcal S}}A_{r,k}.
 \label{eq:pb_excess_kraus}
\end{align}
They obey $\sum_r(K_{p,r}^{\dagger}K_{p,r}+
K_{e,r}^{\dagger}K_{e,r})=P_k$.  Thus the processed/excess instrument used
in Sec.~\ref{sec:pb_write_instrument} is not an abstract probability split:
it is the reduced measurement generated by one explicit source-reduced
collision unitary.

\subsection{Continuous-time GKSL limit and exact drift}

Each attempted update consumes a fresh outcome qutrit and a fresh
transverse-label ancilla prepared in the same ready state, initially
uncorrelated with the boundary, the memory pointer, and all previous
ancillas.  After the collision the ancilla is traced out and is never reused.
The ancilla stream therefore contains no inter-ancilla correlations, so the
reduced evolution is Markovian at the collision-model level.  Because
$k(\lambda_H)$ and $\mathcal S(\lambda_H)$ vary adiabatically, the resulting
generator is time-inhomogeneous but time local.

There are $N\mathcal S\,d\lambda_H$ attempts in a capacity interval
$d\lambda_H$.  The collision limit gives the jump operators
\begin{align}
 L_{p,r}&=\sqrt{N\mathcal S\,\pi_r(1-\Delta)}\,A_{r,k},
 \label{eq:pb_processed_jump}\\
 L_{e,r}&=\sqrt{N\mathcal S\,\pi_r\Delta}\,
 T_{1/(N\mathcal S)}A_{r,k}.
 \label{eq:pb_excess_jump}
\end{align}
The reduced density operator satisfies the time-dependent GKSL equation
\cite{Lindblad1976}
\begin{equation}
 \frac{d\rho}{d\lambda_H}
 =-i[H_{\rm B}+H_{\rm LS},\rho]
 +\sum_{r,\alpha=p,e}\mathcal D[L_{\alpha,r}]\rho,
 \label{eq:pb_gksl}
\end{equation}
where
$\mathcal D[L]\rho=L\rho L^{\dagger}-\tfrac12\{L^{\dagger}L,\rho\}$.
Both the instantaneous boundary Hamiltonian and the Lamb-shift correction
act trivially on the memory-pointer factor,
\begin{equation}
 [H_{\rm B}+H_{\rm LS},\widehat\chi]=0.
 \label{eq:pb_pointer_neutral_hamiltonian}
\end{equation}
Consequently the Hamiltonian contribution to the Heisenberg drift is
$i\langle[H_{\rm B}+H_{\rm LS},\widehat\chi]\rangle=0$.  Because
$A_{r,k}$ also commutes with the memory pointer and
$T_{\epsilon}^{\dagger}\chi T_{\epsilon}=\chi+\epsilon$, the Heisenberg
equation gives
\begin{align}
 \frac{d\langle\chi\rangle}{d\lambda_H}
 &=\sum_r N\mathcal S\,\pi_r\Delta\,
 \frac{1}{N\mathcal S}
 =\Delta,
 \label{eq:pb_hamiltonian_exact_drift}\\
 \frac{d\,\mathrm{Var}(\chi)}{d\lambda_H}
 &=\frac{\Delta(1-\Delta)}{N\mathcal S}
 F_{\rm exc}+\mathcal O(N^{-2}).
 \label{eq:pb_hamiltonian_exact_noise}
\end{align}
Equation~\eqref{eq:pb_hamiltonian_exact_drift} is exact in the pointer
translation and shows explicitly why the mean contains no factor of $N$,
$k$, $1+q$, or $\beta$: the attempt rate $N\mathcal S$ cancels the
intensive pointer step $1/(N\mathcal S)$.

\subsection{Correlated jumps and finite-level corrections}

The current Gram matrix in Eq.~\eqref{eq:pb_current_gram} retains all
quadratic correlations.  Whitening therefore preserves the trace-preserving
instrument even when $G_{+-}\neq0$.  Higher connected cumulants appear in
the Fano factor and in multi-jump operators, while the one-point Ward
identity leaves Eq.~\eqref{eq:pb_hamiltonian_exact_drift} unchanged.  In the
SU(2)$_k$ theory the first affine corrections are controlled by the shifted
level $k+2$,
\begin{equation}
 F_{\rm exc}=1+\mathcal O\!\left(\frac{C_A}{k+2}\right),
 \qquad C_A=2,
 \label{eq:pb_finite_k_fano}
\end{equation}
consistent with the correction hierarchy derived in
Sec.~\ref{sec:kms}.  At the cosmological horizon these corrections
are negligible in the mean and of order $10^{-122}$ in normalized
correlators.

\begin{proposition}[Source-reduced Hamiltonian completion]
The direct-sum CS/WZW boundary Hamiltonian, the KMS-whitened transverse
current transition, and the collision unitary
Eq.~\eqref{eq:pb_collision_unitary} provide a completely positive,
trace-preserving Hamiltonian dilation of the processed/excess instrument.
Its continuum limit is Eq.~\eqref{eq:pb_gksl} and its memory drift is exactly
$d\langle\chi\rangle/d\lambda_H=\Delta$.  Correlated jumps and finite-$k$
terms change the noise and higher cumulants but not the normalized mean
within the source-reduced universality class.
\end{proposition}


\section{Endpoint Response, Common-Gate Universality, and Native Ladder Variables}
\label{sec:endpoint_theorem}

\subsection{Exact endpoint-response theorem}

Let $I$ label a physical register and let $D_I$ generate the memory-sector
dilation of two endpoint operators,
\begin{equation}
[D_I,O_{I,\alpha}]=-ih_{I,\alpha}O_{I,\alpha},
\qquad \alpha\in\{a,b\}.
\label{eq:pb_endpoint_weights}
\end{equation}
For a local comparison event, define the finite gate
\begin{equation}
U_I(X_H)=\exp[-ig_I X_HD_I].
\label{eq:pb_endpoint_gate}
\end{equation}
Because each endpoint is an eigen-operator of the adjoint action, the
Baker--Campbell--Hausdorff series closes exactly,
\begin{equation}
U_I^\dagger O_{I,\alpha}U_I
=e^{g_Ih_{I,\alpha}X_H}O_{I,\alpha}.
\label{eq:pb_endpoint_response}
\end{equation}
For the dimensionless ratio
$\mathcal R_I=\Omega_{I,a}/\Omega_{I,b}$ and
$y_I=\ln\mathcal R_I$,
\begin{equation}
\Delta y_I
=g_I(h_{I,a}-h_{I,b})X_H
\equiv\eta_I X_H,
\qquad
\eta_I=g_I\Delta h_I.
\label{eq:pb_eta_factorization}
\end{equation}
This is exact to all orders in $X_H$.  Representation theory fixes the
weight difference; microscopic matching fixes the gate coefficient; and the
experimental path determines whether the gate is engaged.

A canonical local action reproducing the translation is
\begin{equation}
S_{\rm read}
=\sum_I\int_{\gamma_I}p_I(dy_I-r_I\chi\,dn_I),
\label{eq:pb_read_action}
\end{equation}
where $n_I$ changes only during a physical comparator operation.  The
combined memory--readout sector preserves the shift redundancy under
\begin{equation}
\chi\rightarrow\chi+\chi_c,
\qquad
y_I\rightarrow y_I+r_I\chi_c n_I.
\label{eq:pb_combined_shift}
\end{equation}
Absolute node origins are therefore conventional; completed relative
comparisons are physical.

\subsection{Common-gate corollary}

Suppose the temporal and ruler endpoint representations are restrictions of
one normalized generator $D$ and are read by one physical gate,
\begin{equation}
U_{\rm com}(X_H)=\exp[-ig_{\rm com}X_HD].
\label{eq:pb_common_gate}
\end{equation}
Then $g_T=g_R=g_{\rm com}$ and
\begin{equation}
\frac{\eta_T}{\eta_R}
=\frac{\Delta h_T}{\Delta h_R}.
\label{eq:pb_gain_ratio}
\end{equation}
For equal unit endpoint-weight differences and aligned orientation,
\begin{equation}
\eta_T=\eta_R=g_{\rm com},
\qquad
A_{\rm map}=2g_{\rm com}.
\label{eq:pb_common_gain}
\end{equation}
The equality of the two gains is parameter-independent within the
common-gate realization.  At the theorem level only the common amplitude
remains; Section~\ref{sec:common_gate_proof_complete} evaluates the
normalized minimal portal matrix element and fixes $g_{\rm com}=1$, yielding
$A_{\rm map}=2$ in that EFT.

\subsection{Native recipient coordinates}

The local ladder relation can be written
\begin{equation}
5\log_{10}H_0=M_B+5a_B+25.
\label{eq:pb_ladder_relation}
\end{equation}
Define the logarithmic register coordinates
\begin{equation}
y_T=(\ln10)a_B,
\qquad
y_R=\frac{\ln10}{5}M_B.
\label{eq:pb_native_coords}
\end{equation}
Under the equal-weight common gate,
\begin{equation}
\Delta a_B=\frac{g_{\rm com}X_H}{\ln10},
\qquad
\Delta M_B=\frac{5g_{\rm com}X_H}{\ln10},
\label{eq:pb_native_shifts}
\end{equation}
so the ratio
\begin{equation}
\boxed{\Delta M_B=5\Delta a_B}
\label{eq:pb_native_ratio}
\end{equation}
is independent of both $g_{\rm com}$ and $X_H$.  A compressed local-$H_0$
likelihood sees only the sum $2g_{\rm com}X_H$; a native joint
$(M_B,a_B)$ likelihood can test Eq.~\eqref{eq:pb_native_ratio} directly.

An ordinary common photometric zero-point displacement $\zeta$ obeys
$\Delta M_B=\zeta$ and $5\Delta a_B=-\zeta$, cancelling from
Eq.~\eqref{eq:pb_ladder_relation}.  The proposed temporal response therefore
cannot be a second appearance of the same photometric calibration.

\subsection{Emission-epoch atomic comparator: a restricted exclusion}
\label{sec:emission_atomic}

If source and observer atomic standards carry the same temporal weight, the
endpoint theorem gives
\begin{equation}
\Delta y_T^{\rm src/obs}(z_{\rm em})
=g_T[\chi(a_0)-\chi(a_{\rm em})]
=g_TX_{\rm em\to0}(z_{\rm em}),
\label{eq:pb_emission_response}
\end{equation}
not the full stored value $g_TX_0$.  On the flat hard-gate benchmark,
\begin{equation}
X_{\rm em\to0}(z_{\rm em})
=\int_{z_{\rm exit}}^{z_\star(z_{\rm em})}
\frac{\Delta_{\rm hard}[\mathcal S(z)]}{1+z}\,dz,
\label{eq:pb_post_emission}
\end{equation}
with $z_\star$ clipped to the SCR window.  For $\Omega_{m0}=0.32$,
$z_{\rm exit}=0.020414$ and $z_{\rm entry}=0.619806$.
\begin{table}[!htbp]
\centering
\scriptsize
\begin{tabular}{@{}cccc@{}}
\hline\hline
$z_{\rm em}$ & $X_{\rm em\to0}$ & $X_{\rm em\to0}/X_\Lambda^{\rm hard}$
& mag-equivalent \\
\hline
$0.023$ & $3.19\times10^{-6}$ & $9.25\times10^{-5}$ & $0.000007$ \\
$0.050$ & $3.89\times10^{-4}$ & $0.0113$ & $0.00085$ \\
$0.100$ & $2.48\times10^{-3}$ & $0.0718$ & $0.00538$ \\
$0.150$ & $5.79\times10^{-3}$ & $0.1678$ & $0.01257$ \\
\hline\hline
\end{tabular}
\caption{Post-emission loading of an equally charged source--observer
atomic comparator.  The last column is $(5/\ln10)X_{\rm em\to0}$.}
\label{tab:pb_post_emission}
\end{table}
The realization is too small and redshift dependent to produce the constant
full-unit temporal loading.  This excludes that endpoint class; it does not
prove that every possible observer-side phase comparator is neutral.



\section{Gauge-Neutral Current Bilocal as a Microscopic Common Gate}
\label{sec:bilocal_completion}

The conditional endpoint theorem can be sharpened within the same
SU(2)$_k$ Chern--Simons/WZW boundary theory used for the horizon
throughput calculation.  The microscopic comparator is taken to be a
boundary-anchored, gauge-neutral bilocal.  Its two endpoint insertions are
the two physical transverse affine currents that remain after the
source-selected Gauss reduction.

\paragraph{Affine algebra and exact endpoint weights.}
Let $J^a_n$ be the modes of the chiral SU(2)$_k$ currents~\cite{Witten1989,EMSS1989,DiFrancesco1997}.  They obey
\begin{equation}
 [J^a_m,J^b_n]
 =i\epsilon^{ab}{}_{c}J^c_{m+n}
 +k m\,\delta^{ab}\delta_{m+n,0}.
\label{eq:micro2_affine_algebra}
\end{equation}
The Sugawara generator is
\begin{equation}
 L_n=\frac{1}{2(k+2)}\sum_m :J^a_{n-m}J^a_m: .
\label{eq:micro2_sugawara}
\end{equation}
The current Ward identity gives
\begin{equation}
 [L_0,J^a_{-1}]=J^a_{-1},
\label{eq:micro2_current_weight}
\end{equation}
so every affine current has exact boundary scaling weight one,
independent of the Chern--Simons level $k$.

The homogeneous horizon source selects an internal direction $\hat n^a$.
Choosing it as the Cartan axis, define
\begin{equation}
 J^3\equiv \hat n_aJ^a,
 \qquad
 J^\pm\equiv\frac{J^1\pm iJ^2}{\sqrt2}.
\end{equation}
The homogeneous Gauss constraint fixes the longitudinal component $J^3$;
the sourced physical current sector is therefore
$\mathfrak{su}(2)/\mathfrak{u}(1)_{\hat n}$ and contains exactly the two
root currents $J^\pm$.  Their residual charges are
\begin{equation}
 [J^3_0,J^\pm_n]=\pm J^\pm_n.
\label{eq:micro2_root_charge}
\end{equation}
Thus the two surviving currents have opposite gauge charges but the same
positive Weyl weight, $h_+=h_-=1$.

\paragraph{The minimal physical comparator.}
A single $J^+$ or $J^-$ is not invariant under the residual
$U(1)_{\hat n}$.  The lowest-weight non-trivial comparison between neutral
physical standards is a conjugate two-endpoint insertion.  In point-split
form write
\begin{equation}
 \widehat{\mathcal O}_{TR}(x_T,x_R;\gamma)
 =J^+(x_T)\,\widehat W_\gamma[A^3] \,J^-(x_R),
\label{eq:micro2_bilocal}
\end{equation}
where $\widehat W_\gamma[A^3]$ is the residual parallel transporter joining
the two endpoints.  Its gauge transformation cancels the opposite endpoint
charges.  At the topological Wilson-line fixed point~\cite{Gawedzki2004,BachasGaberdiel2004} the transporter carries
no local boundary Weyl weight, so
\begin{equation}
 h(\widehat{\mathcal O}_{TR})=h_++h_-=2.
\label{eq:micro2_total_weight}
\end{equation}
Equivalently, the coincident limit is the neutral normal-ordered current
bilinear $:J^+J^-:$, whose weight is two.  The factor two therefore follows
from residual gauge neutrality plus the exact current Ward identity; it is
not obtained by counting stages of a data-analysis pipeline.

\paragraph{Direct correlator calculation.}
On a locally flat boundary patch the current algebra fixes
\begin{equation}
 \left\langle J^+(z_T)J^-(z_R)\right\rangle_0
 =\frac{C_k}{(z_T-z_R)^2},
\label{eq:micro2_current_correlator}
\end{equation}
where $C_k$ depends on the current normalization but not on the endpoint
separation.  Let the protected homogeneous memory be the canonically
normalized boundary Weyl zero mode,
\begin{equation}
 \sigma_{AB}(\chi)=e^{2\chi}\bar\sigma_{AB}.
\label{eq:micro2_memory_metric}
\end{equation}
A constant shift $\chi_0$ multiplies every physical boundary separation by
$e^{\chi_0}$.  Since the parallel transporter is topological, the normalized
bilocal transfer is therefore
\begin{align}
 \mathcal C_{TR}(\chi_0)
 &\equiv
 \frac{\langle\widehat{\mathcal O}_{TR}\rangle_{\chi_0}}
      {\langle\widehat{\mathcal O}_{TR}\rangle_{0}}
 =e^{-2\chi_0},
\label{eq:micro2_bilocal_response}\\
 \frac{\partial}{\partial\chi_0}\ln\mathcal C_{TR}
 &=-2.
\label{eq:micro2_bilocal_derivative}
\end{align}
The unknown normalization $C_k$ and any overall comparator coupling cancel
from the ratio.  This is the explicit boundary matrix element requested by
the operational matching programme.

The same result follows from the finite Weyl gate.  If $D_{\rm W}$ generates
\eqref{eq:micro2_memory_metric}, then
\begin{equation}
 [D_{\rm W},J^\pm]=-iJ^\pm,
 \qquad
 [D_{\rm W},\widehat{\mathcal O}_{TR}]
 =-2i\widehat{\mathcal O}_{TR},
\end{equation}
and
\begin{equation}
 U_{\rm W}(\chi_0)=e^{-i\chi_0D_{\rm W}}
 \quad\Longrightarrow\quad
 U_{\rm W}^{\dagger}\widehat{\mathcal O}_{TR}U_{\rm W}
 =e^{2\epsilon\chi_0}\widehat{\mathcal O}_{TR}.
\label{eq:micro2_finite_gate}
\end{equation}
Here $\epsilon=\pm1$ is the orientation of the physical transfer ratio.
The response is exact to all orders in $\chi_0$.  There is no sign conflict
between Eqs.~\eqref{eq:micro2_bilocal_response} and
\eqref{eq:micro2_finite_gate}.  The former is a passive Weyl rescaling of
the physical endpoint separation with the operator insertions held fixed,
so a weight-two correlator scales as $e^{-2\chi_0}$.  The latter is the
active transformation of the bilocal operator or of the oriented transfer
coordinate.  The orientation variable specifies whether the physical
transfer uses the correlator or its inverse: $\epsilon=-1$ reproduces the
passive-correlator convention, whereas $\epsilon=+1$ describes the
oppositely oriented transfer ratio.

\paragraph{Local separation into two unit endpoint contributions.}
The point-split form makes the two channels local.  Allow a Weyl source that
can be varied independently near the two endpoints.  The local Ward identity
is
\begin{equation}
 \frac{\delta\ln\langle\widehat{\mathcal O}_{TR}\rangle}
      {\delta\chi(x)}
 =-\delta^{(2)}(x-x_T)-\delta^{(2)}(x-x_R),
\label{eq:micro2_local_ward}
\end{equation}
up to the orientation convention.  Integrating over a neighborhood of one
endpoint gives one unit, while integrating over both gives two.  Hence the
microscopic bilocal supports the decomposition
\begin{equation}
 |\Delta y_T|=\chi_0,
 \qquad
 |\Delta y_R|=\chi_0,
 \qquad
 |\Delta y_T+\Delta y_R|=2\chi_0.
\label{eq:micro2_unit_endpoint_shifts}
\end{equation}
The assignment of $x_T$ to the temporal matter interface and $x_R$ to the
ruler/radiative interface is a matter--boundary portal identification; the
unit endpoint weights themselves are fixed by the current algebra.

\paragraph{Native ladder variables.}
For the operational coordinates
\begin{equation}
 y_T=(\ln10)a_B,
 \qquad
 y_R=\frac{\ln10}{5}M_B,
\end{equation}
the aligned orientation gives
\begin{equation}
 \Delta a_B=\frac{\chi_0}{\ln10},
 \qquad
 \Delta M_B=\frac{5\chi_0}{\ln10},
\end{equation}
and therefore
\begin{equation}
 \boxed{
 \frac{H_{0,\rm ladder}}{H_{0,\rm phys}}
 =e^{2\chi_0}.}
\label{eq:micro2_H0_map}
\end{equation}
Using the capacity-cycle write theorem,
$\langle\chi_0\rangle=X_H$, Eq.~\eqref{eq:micro2_H0_map} becomes the
common-gate ladder map $e^{2X_H}$ at the semiclassical saddle.

\paragraph{Status of the completion.}
The bilocal derives, per unit canonically normalized boundary memory:
(i) exactly two source-accessible current directions after Gauss reduction;
(ii) one exact Weyl unit at each endpoint; (iii) a minimal gauge-neutral
comparator of total weight two; and (iv) an exact exponential response with
no continuous conditional readout coefficient.  Section~\ref{sec:portals_map}
now supplies explicit temporal and radiative matter operators, a single
normalized defect vertex, the serial-versus-parallel proof of one ladder
engagement, and the retarded passive orientation.  Together with the
capacity-cycle write theorem $\chi_0=X_H$, these results complete the
mapping within the minimal local comparator EFT.  Deriving the same portal
operators and processed/excess instrument from the exact Standard-Model
plus time-dependent CS/WZW Hamiltonian remains a UV embedding problem; it is
not an undetermined coefficient in the low-energy map.


% ======================================================================
%  CALIBRATION HOLONOMY
% ======================================================================
\section{Register-Valued Calibration Holonomy}
\label{sec:calibration_holonomy}

\subsection{Why a source-local luminosity modification calibrates out}

Let a universal local effect multiply the standardized luminosity of class
$\alpha$ by $g_\alpha(X)$.  An anchor at known geometric distance $d_A$
has
\begin{equation}
F_A=\frac{g_\alpha(X)L_\alpha^{(0)}}{4\pi d_A^2},
\qquad
L_\alpha^{\rm cal}=4\pi d_A^2F_A=g_\alpha(X)L_\alpha^{(0)}.
\end{equation}
For a target of the same class,
$F_T=g_\alpha(X)L_\alpha^{(0)}/(4\pi d_T^2)$, and therefore
\begin{equation}
d_T^{\rm cal}
=\left(\frac{L_\alpha^{\rm cal}}{4\pi F_T}\right)^{1/2}=d_T.
\label{eq:pb_source_cancel}
\end{equation}
A universal source-local change, or a common multiplicative detector gain,
cancels between calibration and application.  A surviving response must
therefore be environment dependent, propagation dependent, or a relative
transfer between inequivalent physical standards.  The first option predicts
host correlations; the second distorts the distance--redshift relation.  The
baseline adopts the third.

\subsection{Directed graph and loop invariant}

Let $G=(V,E)$ be a graph of physically realized standards and irreducible
comparators.  Each oriented edge carries a register-sensitivity covector
\begin{equation}
\mathbf w_e=(w_{T,e},w_{R,e}),
\qquad
\boldsymbol\eta=(\eta_T,\eta_R).
\label{eq:pb_edge_vectors}
\end{equation}
Its logarithmic response is
\begin{equation}
\delta\ln\mathcal T_e
=(\mathbf w_e\cdot\boldsymbol\eta)X_H.
\label{eq:pb_edge_response}
\end{equation}
For a path $p$,
\begin{equation}
A_p=\sum_{e\in p}\mathbf w_e\cdot\boldsymbol\eta.
\label{eq:pb_path_charge}
\end{equation}
Two independently realizable paths $p,p'$ with the same endpoints define
\begin{equation}
\boxed{
\mathcal H_{pp'}
=\ln\!\left[
\frac{\mathcal T_p/\mathcal T_{p'}}
{\mathcal T_p^{(0)}/\mathcal T_{p'}^{(0)}}
\right]
=(A_p-A_{p'})X_H.}
\label{eq:pb_loop_holonomy}
\end{equation}
Under a node zero-point redefinition,
$\ln\mathcal T_e\rightarrow\ln\mathcal T_e+
\lambda_{t(e)}-\lambda_{s(e)}$, the node terms telescope from every closed
loop.  If $\mathbf w_e=\mathbf r_{t(e)}-\mathbf r_{s(e)}$, all loop
residuals vanish.  Such a response is a convention, not a physical ledger.

An admissible edge is an independently realizable physical operation whose
output changes a relative standard and for which a redundant route can be
constructed in principle.  Passband corrections, software stages, nuisance
parameters, and terminal algebra are not extra edges.  The graph therefore
does not generate the factor two.  The boundary common gate supplies the two
primitive unit-weight endpoints; the graph audits whether a real experimental
path engages them and whether the result survives zero-point redefinitions.

\subsection{Local edge gate and microscopic matching target}

The edge version of the common gate can be written
\begin{equation}
U_e(X_H)=\exp\!\left[-iX_H
\sum_{I=T,R}w_{I,e}g_I D_I\right].
\label{eq:pb_edge_gate}
\end{equation}
A UV or boundary completion should calculate the normalized endpoint matrix
element
\begin{equation}
a_e=
\left.\frac{\partial}{\partial X}
\ln\frac{\langle O_{t(e)}W_e(X)O_{s(e)}\rangle}
{\langle O_{t(e)}W_e(0)O_{s(e)}\rangle}
\right|_{X=0},
\label{eq:pb_matching_target}
\end{equation}
where $W_e$ is the physical parallel transporter.  The current bilocal of
Sec.~\ref{sec:bilocal_completion} evaluates this target for one explicit
boundary-defect realization.

\subsection{Pre-registered closure amplitudes}

For the canonical unit specialization and $X_0=0.0359$, one register unit is
\begin{equation}
e^{X_0}=1.03655,
\qquad
\frac{5X_0}{\ln10}=0.07796\ {\rm mag},
\label{eq:pb_unit_numbers}
\end{equation}
and two aligned units give $e^{2X_0}=1.07444$ and $0.15591$ mag-equivalent.
A useful pre-registration scheme is
\begin{center}
\begin{tabular}{ccc}
\toprule
$\Delta\mathbf W$ & loop prediction & role \\
\midrule
$(0,0)$ & $0$ & systematic control \\
$(1,0)$ & $X_0$ & temporal closure \\
$(0,1)$ & $X_0$ & ruler closure \\
$(1,1)$ & $2X_0$ & full common-gate closure \\
\bottomrule
\end{tabular}
\end{center}
Equal-sensitivity loops are valuable controls but cannot distinguish the
ledger from a neutral theory.  A decisive experiment requires a
pre-registered unequal-register loop.



% ======================================================================
%  MATTER PORTALS, COMMON DEFECT, AND SINGLE ENGAGEMENT
% ======================================================================
\section{Local Matter-Comparator Portals and the Completed Ladder Map}
\label{sec:portals_map}

The boundary theory fixes the memory history and the response of the two
minimal current endpoints.  To turn those endpoint residues into an
astronomical observable one must specify what is physically compared.  This
section gives an explicit local comparator EFT.  Its scope is deliberately
operational: the memory does not change an atomic Hamiltonian throughout
spacetime, alter the luminosity emitted by a star, or act on every propagating
photon.  It acts during the preparation of a relative calibration record.
The matter systems supply the two observables being compared; canonical
pointer coordinates carry the transferable zero points.

\subsection{Comparator observables and a local measurement record}
\label{sec:portal_pointer}

Let $n(\tau)$ be a switching coordinate that is constant except during one
completed physical comparison and is normalized by
\begin{equation}
 \int_{\gamma_{\rm cal}}dn=1 .
 \label{eq:portal_switch_unit}
\end{equation}
The value of $n$ is not a software flag.  It labels the state preparation of
the shared calibration standard.  Introduce two pointer pairs
\begin{equation}
 [\widehat y_I,\widehat p_J]=i\delta_{IJ},
 \qquad I,J\in\{T,R\},
 \label{eq:portal_pointer_algebra}
\end{equation}
and two positive matter-comparator observables $\widehat{\mathcal R}_T$ and
$\widehat{\mathcal R}_R$.  Define
\begin{equation}
 \widehat Z_T\equiv\ln\widehat{\mathcal R}_T,
 \qquad
 \widehat Z_R\equiv-\ln\widehat{\mathcal R}_R.
 \label{eq:portal_matter_logs}
\end{equation}
The opposite sign in $Z_R$ is the usual magnitude convention: a smaller
radiative amplitude corresponds to a larger logarithmic ruler coordinate.

The minimal local comparator action is
\begin{multline}
 S_{\rm cmp}=\int_{\gamma_{\rm cal}}d\tau\,
 \Big\{
 p_T[\dot y_T-\dot n\,Z_T]
 +p_R[\dot y_R-\dot n\,Z_R]\\
 -g_{\rm com}\,\chi\,\dot n\,(p_T+p_R)
 \Big\}.
 \label{eq:portal_comparator_action}
\end{multline}
The first line is a standard von-Neumann-type premeasurement: it records the
matter ratios in two pointer coordinates.  The second line is the common
boundary defect.  Varying with respect to the pointer momenta gives
\begin{equation}
 \dot y_T=\dot n(Z_T+g_{\rm com}\chi),
 \qquad
 \dot y_R=\dot n(Z_R+g_{\rm com}\chi).
 \label{eq:portal_pointer_eom}
\end{equation}
For a short read event, $Z_I$ and $\chi$ are adiabatically constant.  The
memory-dependent part of the completed record is therefore
\begin{equation}
 \boxed{
 \delta_\chi y_T=\delta_\chi y_R=g_{\rm com}\chi_0 .}
 \label{eq:portal_equal_pointer_shift}
\end{equation}
Equivalently, the Stinespring unitary acting on the two pointers is
\begin{equation}
 U_{\rm read}(\chi_0)
 =\exp[-ig_{\rm com}\chi_0(p_T+p_R)],
 \label{eq:portal_stinespring_gate}
\end{equation}
so that $U_{\rm read}^\dagger y_IU_{\rm read}=y_I+g_{\rm com}\chi_0$.
The matter systems themselves are not continuously driven by $\chi$; the
portal translates the relative calibration record when the comparison is
made.

\subsection{Temporal phase portal}
\label{sec:temporal_portal_complete}

A physical clock comparison is a comparison of phases, not of a coordinate
label.  Let $A$ be a local two-level reference with Hamiltonian
\begin{equation}
 H_A=\frac{\hbar\Omega_A}{2}\sigma_z,
 \qquad
 \Theta_A(\tau)=\int^\tau\Omega_A(\tau')d\tau',
 \label{eq:portal_atomic_clock}
\end{equation}
whose phase can be read by a Ramsey comparison~\cite{Ramsey1950}.  The
memory-neutral reference is a processed-sector boundary phase $\Theta_K$.
It is pinned to the geometrically preferred capacity-cycle flow by the
first-order action
\begin{equation}
 S_K=\int_{\gamma_{\rm cal}}\Pi_K
 \left(d\Theta_K-\omega_K d\lambda_H\right),
 \label{eq:portal_kodama_clock_action}
\end{equation}
where $d\lambda_H=H_u d\tau$ on the local cosmological congruence.  The
Kodama construction supplies a preferred dynamical time flow in spherical
symmetry~\cite{Kodama1980,AbreuVisser2010}; Eq.~\eqref{eq:portal_kodama_clock_action}
turns that flow into a physical boundary phase record.  Because it belongs
to the processed sector, $\Theta_K$ is invariant under the memory shift and
therefore satisfies
\begin{equation}
 [D,\Omega_K]=0,
 \qquad
 \Omega_K\equiv\frac{d\Theta_K}{d\tau}.
 \label{eq:portal_kodama_neutral}
\end{equation}
The temporal matter ratio recorded by the comparator is
\begin{equation}
 \mathcal R_T\equiv\frac{\Omega_A}{\Omega_K},
 \qquad
 Z_T=\ln\mathcal R_T.
 \label{eq:portal_temporal_ratio_complete}
\end{equation}
The local phase endpoint is the weight-one $J^+$ endpoint of the current
bilocal, while the processed-sector phase reference is a singlet.  The
local Ward residue therefore gives $\Delta h_T=1$.  The portal is local at
readout: it does not imply a universal variation of atomic transition
frequencies, nor does it alter the spectroscopic redshift of every source.
Indeed, the equal-charge source--observer realization remains excluded by
Table~\ref{tab:pb_post_emission}.

The observable coordinate associated with this phase comparison is
\begin{equation}
 y_T=(\ln10)a_B .
 \label{eq:portal_temporal_native}
\end{equation}
This does not identify the regression variable $a_B$ itself with a
microscopic clock.  Rather, $a_B$ is the native global zero point in which a
completed local phase comparison appears.

\subsection{Radiative-amplitude portal}
\label{sec:radiative_portal_complete}

For a standardized optical band, Glauber photodetection theory identifies
the detected intensity with the normally ordered first-order field
correlator~\cite{Glauber1963a,Glauber1963b}.  Define
\begin{equation}
 \mathcal I_L\equiv
 \int d\nu\,W(\nu)
 \left\langle E^{(-)}_\nu E^{(+)}_\nu\right\rangle,
 \qquad
 \mathcal A_L\equiv\mathcal I_L^{1/2},
 \label{eq:portal_glauber_amplitude}
\end{equation}
where $W(\nu)$ includes the passband and the ordinary light-curve
standardization.  Let $\mathcal A_G$ be a geometrically anchored reference
amplitude prepared and read in the same radiometric chain.  The ruler
comparison is
\begin{equation}
 \mathcal R_R\equiv\frac{\mathcal A_L}{\mathcal A_G},
 \qquad
 Z_R=-\ln\mathcal R_R.
 \label{eq:portal_radiative_ratio_complete}
\end{equation}
A common multiplicative detector gain cancels in $\mathcal R_R$.  Likewise,
a universal source-local luminosity change cancels between calibration and
application as shown in Eq.~\eqref{eq:pb_source_cancel}.  The portal acts on
the relative transfer record after these common-mode quantities have
cancelled; it is not opacity, a colour term, or a change in stellar
emission.

The radiative endpoint is the conjugate $J^-$ endpoint.  Its affine weight
has unit magnitude, and its residual charge is opposite to the phase
endpoint.  The native ruler coordinate is
\begin{equation}
 y_R=-\ln\mathcal R_R
 =\frac{\ln10}{5}M_B+{\rm const.}
 \label{eq:portal_ruler_native_complete}
\end{equation}
so a positive pointer translation corresponds to a fainter calibrated
absolute magnitude.

\subsection{Why the two portals are one physical common gate}
\label{sec:common_gate_proof_complete}

The two matter records are attached to the two endpoints of the single
boundary bilocal of Eq.~\eqref{eq:micro2_bilocal}.  To make the portal
operators explicit, take
\begin{equation}
\begin{aligned}
 \mathcal Q_T^+&\equiv \sigma_+e^{-i\Theta_K},\\
 \mathcal Q_R^-&\equiv
 \frac{b_G^\dagger b_L}{\sqrt{(N_G+1)N_L}},
 \qquad N_X\equiv b_X^\dagger b_X.
\end{aligned}
 \label{eq:portal_matter_transition_ops}
\end{equation}
where $\mathcal Q_T^+$ creates one relative atomic--boundary phase
excitation and $\mathcal Q_R^-$ transfers one normalized radiative-amplitude
quantum from the standardized mode $b_L$ to the geometric reference mode
$b_G$.  On each nonzero fixed-occupation sector the square-root
normalization makes the transfer a unit partial isometry.  Their Hermitian
conjugates implement the reverse comparison.

A convenient microscopic portal vertex is
\begin{multline}
 S_{\rm port}^{\rm micro}=\lambda\int d\tau\,\dot n(\tau)
 \Big[
 \mathcal Q_T^+J^+(x_T)W_\gamma J^-(x_R)\mathcal Q_R^-\\
 +{\rm h.c.}\Big].
 \label{eq:portal_uv_vertex}
\end{multline}
The matter transitions are dressed with compensating residual charges, so
the complete operator is gauge neutral and has one Wilson coefficient
$\lambda$.  A single-endpoint vertex is forbidden by the residual
$U(1)_{\hat n}$ charge.  At the lowest boundary dimension there is only one
neutral operator with the required endpoint quantum numbers, namely the
bilocal in Eq.~\eqref{eq:portal_uv_vertex} and its Hermitian conjugate.  A
second coefficient multiplying an antisymmetric endpoint response would
require a second linearly independent neutral operator or an explicit
endpoint-breaking spurion.  Neither exists in the minimal operator basis.
Integrating out the current sector therefore produces one common
translation coefficient, while the local Ward identity supplies equal
unit residues at its two ends:
\begin{equation}
 \boxed{g_T=g_R\equiv g_{\rm com}.}
 \label{eq:portal_common_gate_proved}
\end{equation}
This is stronger than the statement that two portals happen to couple to the
same scalar: they are the two residues of one irreducible gauge-neutral
defect insertion.  Higher-dimensional endpoint-asymmetric operators are a
well-defined UV correction rather than an additional leading coupling.

The normalized response of a \emph{completed} minimal-charge comparison is
also fixed.  Consider the portal matrix element conditional on the
transition,
\begin{equation}
 \mathcal C_{\rm port}(\chi)
 =\frac{\langle f|\mathcal V_{\rm port}(\chi)|i\rangle}
        {\langle f|\mathcal V_{\rm port}(0)|i\rangle},
 \qquad
 \mathcal V_{\rm port}\equiv
 \mathcal Q_T^+J^+W_\gamma J^-\mathcal Q_R^- .
 \label{eq:portal_normalized_matrix_element}
\end{equation}
The overall Wilson coefficient $\lambda$, the normalized matter matrix
elements, and the current two-point normalization cancel between numerator
and denominator.  What remains is the affine Ward residue.  The write
sector has already fixed $\chi$ as the canonically normalized translation
coordinate, the two root currents have minimal unit weight, and the
completed switch obeys Eq.~\eqref{eq:portal_switch_unit}.  Therefore
\begin{equation}
 \left.\frac{\partial}{\partial\chi}
 \ln\mathcal C_{\rm port}(\chi)\right|_{\chi=0}=2,
 \qquad
 \boxed{g_{\rm com}=1}
 \label{eq:portal_unit_common_gate}
\end{equation}
for the two-endpoint gate, with one unit assigned locally to each endpoint
by Eq.~\eqref{eq:micro2_local_ward}.  A non-unit $g_{\rm com}$ would require
a non-minimal endpoint representation, a different normalization of the
already-fixed memory coordinate, or an explicit endpoint-asymmetric
counterterm.  Thus unit common matching is a theorem of the minimal
normalized comparator EFT.  It is not yet a theorem that the Standard Model
realizes this particular portal vertex.

\subsection{Retarded passivity and the endpoint orientation}
\label{sec:orientation_complete}

The current algebra fixes the magnitudes of the two endpoint residues but
not the arrow assigned to a physical transfer.  The arrow is supplied by
the retarded comparator instrument.  The write variable is non-negative for
a positive excess history, $\chi_0=X_H\ge0$, and the forward portal
transition in Eq.~\eqref{eq:portal_uv_vertex} is
\begin{equation}
 \mathcal Q_R^-\longrightarrow \mathcal Q_T^+ .
 \label{eq:portal_forward_arrow}
\end{equation}
The arrow removes one normalized radiative-reference quantum and creates one
relative phase-record excitation.
The reverse operator is the Hermitian-conjugate comparison and belongs to
the inverse path.  Retarded boundary conditions and non-negative entropy
production select the forward jump for the calibration path; they do not
allow the reverse path to be silently substituted after the data are seen.

The radiative part of the forward comparison is passive.  In the absence of
an external radiative energy source, a quantum-limited local comparison may
preserve or attenuate a coherent amplitude but cannot noiselessly amplify
it.  The canonical passive Gaussian channel is the beam-splitter dilation
\cite{HolevoWerner2001,Caves1982},
\begin{equation}
 b_L^{\rm cal}=e^{-\chi_0}b_L^{\rm phys}
 +\sqrt{1-e^{-2\chi_0}}\,v,
 \label{eq:portal_beamsplitter_dilation}
\end{equation}
where $v$ is the environment mode of the local comparator.  Consequently
\begin{equation}
 \frac{\mathcal A_L^{\rm cal}}{\mathcal A_L^{\rm phys}}
 =e^{-\chi_0}.
 \label{eq:portal_passive_radiative_sign}
\end{equation}
This is a local record-transfer channel, not attenuation accumulated by
astronomical photons in flight.

The same oriented bilocal transfers the conjugate unit to the temporal
phase record.  With $\mathcal R_T=\Omega_A/\Omega_K$ defined in the forward
phase-counting direction,
\begin{equation}
 \frac{\mathcal R_T^{\rm cal}}{\mathcal R_T^{\rm phys}}
 =e^{+\chi_0}.
 \label{eq:portal_temporal_sign}
\end{equation}
The opposite signs in the two physical ratios are the two ends of one
oriented transfer: the radiative amplitude is the source endpoint and the
relative phase record is the receiving endpoint.  With the logarithmic
coordinates of Eq.~\eqref{eq:portal_matter_logs}, both native pointers shift
in the same direction,
\begin{equation}
 \boxed{\Delta y_T=\Delta y_R=+\chi_0.}
 \label{eq:portal_signed_pointer_shifts}
\end{equation}
Reversing the comparison replaces
$\mathcal V_{\rm port}\leftrightarrow\mathcal V_{\rm port}^\dagger$,
reverses the Wilson line, and changes $n:1\to0$.  It therefore gives
$-\chi_0$ and
\begin{equation}
 \mathcal T_{p^{-1}}=\mathcal T_p^{-1}.
 \label{eq:portal_inverse_path}
\end{equation}
The sign is thus fixed within the retarded passive portal EFT.  It is not
fixed by the SU(2) representation weights alone; a different microscopic
portal that violates passivity or reverses the thermodynamic arrow would be
a different theory.

\subsection{Single engagement in the actual distance ladder}
\label{sec:single_engagement_complete}

The SH0ES ladder is a single simultaneous fit linking geometric anchors,
primary distance indicators, calibrator SNe~Ia, and Hubble-flow SNe~Ia
\cite{Riess2022}.  Its relevant global SN zero points can be written
schematically as
\begin{align}
 m_{B,i}^{0}&=\mu_i+M_B+\epsilon_i,
 &&i\in\mathcal C,
 \label{eq:portal_calibrator_block}\\
 m_{B,j}^{0}&=5\log_{10}\mathcal D_L(z_j;\mathbf q)-5a_B+\epsilon_j,
 &&j\in\mathcal H,
 \label{eq:portal_hubble_block}
\end{align}
where $\mathcal C$ and $\mathcal H$ denote the calibrator and Hubble-flow
samples, and $\mathcal D_L$ is the dimensionless low-redshift distance
function with the overall Hubble normalization removed.  The exact
Pantheon+/SH0ES analysis retains the full covariance and solves the ladder
simultaneously~\cite{Riess2022,Brout2022}; the essential point here is that
$M_B$ and $a_B$ are each one shared global latent coordinate, not one
parameter per supernova.

Let $\mathbf y=(y_T,y_R)^T$ and let $\mathbf y^{(0)}$ be the neutral
calibration state.  The common defect prepares the shifted state through the
kernel
\begin{equation}
 K_{\chi}(\mathbf y|\mathbf y^{(0)})
 =\delta^{(2)}\!\left[
 \mathbf y-\mathbf y^{(0)}-\chi_0
 \begin{pmatrix}1\\1\end{pmatrix}
 \right].
 \label{eq:portal_gate_kernel}
\end{equation}
The ladder contains many objects but only one serial conversion of the
shared standard.  Calibrator hosts provide parallel estimates of the same
$M_B$ coordinate, and Hubble-flow supernovae provide parallel estimates of
the same $a_B$ coordinate.  Parallel replicas multiply likelihood factors
and reduce statistical uncertainty; they do not compose transfer maps and
therefore do not add powers of $e^{\chi_0}$.  Only serial crossings of
inequivalent standard sectors contribute to the holonomy.

The complete likelihood has the form
\begin{multline}
 \mathcal L(D|\chi_0)=
 \int d^2y\,d^2y^{(0)}\,
 K_\chi(\mathbf y|\mathbf y^{(0)})\,\pi_0(\mathbf y^{(0)})\\
 \times
 \prod_{i\in\mathcal C}\mathcal L_i(D_i|y_R)
 \prod_{j\in\mathcal H}\mathcal L_j(D_j|y_T).
 \label{eq:portal_hierarchical_likelihood}
\end{multline}
There is one and only one kernel $K_\chi$.  Increasing the number of
calibrators or Hubble-flow SNe multiplies likelihood factors but does not
reapply the defect.  The physical switch therefore has the topological
intersection number
\begin{equation}
 \nu_{\rm cal}=\int_{p_{\rm ladder}}dn=1.
 \label{eq:portal_single_intersection}
\end{equation}
Subdividing the analysis, changing the number of objects, or rewriting the
normal equations leaves $\nu_{\rm cal}$ unchanged.

\begin{proposition}[Single-engagement ladder]
Within the global hierarchical ladder
Eqs.~\eqref{eq:portal_calibrator_block}--\eqref{eq:portal_hierarchical_likelihood},
the common boundary defect acts once on the shared two-coordinate
standardization state.  Consequently the ladder sensitivities are
\begin{equation}
 (w_T,w_R)_{\rm ladder}=(1,1),
 \label{eq:portal_ladder_weights_derived}
\end{equation}
independently of the number of anchors, calibrator SNe, or Hubble-flow SNe.
\end{proposition}

This result does not count $M_B$ and $a_B$ as two unrelated pipeline edges.
They are the two endpoint coordinates of one prepared calibration state.  A
common photometric zero-point perturbation instead excites the antisymmetric
mode and therefore cancels, as shown next.

\subsection{Physical and conventional normal modes}
\label{sec:portal_normal_modes}

In the native coordinates,
\begin{equation}
 \ln H_0=y_T+y_R+{\rm const.}
 \label{eq:portal_h0_vector}
\end{equation}
The common gate excites
\begin{equation}
 \delta\mathbf y_+=\chi_0
 \begin{pmatrix}1\\1\end{pmatrix},
 \label{eq:portal_common_mode}
\end{equation}
whereas a common photometric zero-point shift $\zeta$ gives
\begin{equation}
 \delta\mathbf y_-=\frac{\ln10}{5}\zeta
 \begin{pmatrix}-1\\+1\end{pmatrix}.
 \label{eq:portal_photometric_mode}
\end{equation}
The latter is a coboundary mode:
$(1,1)\cdot\delta\mathbf y_-=0$.  The former is the physical loop mode:
$(1,1)\cdot\delta\mathbf y_+=2\chi_0$.  Thus the two-channel response is
orthogonal to the familiar photometric cancellation and cannot be produced
by applying one zero-point error twice.

\subsection{Signed native map}
\label{sec:signed_native_map}

Using $\chi_0=X_H$, Eqs.~\eqref{eq:portal_signed_pointer_shifts} and
\eqref{eq:pb_native_coords} give
\begin{equation}
 \boxed{
 \Delta a_B=+\frac{X_H}{\ln10},
 \qquad
 \Delta M_B=+\frac{5X_H}{\ln10}.}
 \label{eq:portal_signed_native_shifts}
\end{equation}
The parameter-independent common-gate relation remains
\begin{equation}
 \boxed{\Delta M_B=5\Delta a_B.}
 \label{eq:portal_signed_native_ratio}
\end{equation}
At fixed apparent magnitude, the calibrator distance modulus changes by
$\Delta\mu=-\Delta M_B$, so
\begin{equation}
 \frac{d_{\rm cal}}{d_{\rm phys}}=e^{-X_H}.
 \label{eq:portal_distance_sign}
\end{equation}
The temporal endpoint raises the Hubble-flow intercept by the same
logarithmic amount.  The exact ladder relation then yields
\begin{equation}
 \boxed{
 \Delta\ln H_0=2X_H,
 \qquad
 \frac{H_{0,\rm ladder}}{H_{0,\rm phys}}=e^{2X_H}.}
 \label{eq:portal_complete_h0_map}
\end{equation}
For $X_0=0.0359$, the two endpoint shifts are each $0.07796$ mag-equivalent
and the total factor is $1.07444$.  The reverse comparison gives the inverse
map, as required by loop orientation.

\subsection{Scope of the portal completion}
\label{sec:portal_scope}

Equations~\eqref{eq:portal_atomic_clock}--\eqref{eq:portal_complete_h0_map}
complete the four operational tasks in a definite low-energy theory: they
construct temporal and radiative matter comparators, place them on one
normalized gauge-neutral defect, show that the global distance ladder
crosses that defect once, and fix the signed map from the retarded passive
orientation.  Sections~\ref{sec:exact_hamiltonian_dilation},
\ref{sec:concrete_matter_realizations}, and
\ref{sec:native_joint_likelihood} further supply the source-reduced
Hamiltonian dilation, a concrete metrological realization and counterterm
basis, and the native statistical test.  The completion is microscopic in
the open-system EFT sense.  A different non-perturbative embedding may
select another portal or break the common-gate symmetry, but it cannot be
represented as an uncomputed normalization inside the minimal theory.

\subsection{Prior-blind prediction and remaining observational work}
\label{sec:portal_statistical_interface}

Let the companion cosmology paper determine a non-ladder posterior
$p(\theta|D_A)$, including the physical expansion rate and the exposure.
The completed portal EFT predicts the push-forward
\begin{multline}
 p(H_{0,\rm ladder}|D_A)
 =\int d\theta\,p(\theta|D_A)\\
 \times\delta\!\left[
 H_{0,\rm ladder}-H_{0,\rm phys}(\theta)
 e^{2X_H(\theta)}
 \right].
 \label{eq:pb_pushforward}
\end{multline}
This distribution should be frozen before confronting the withheld local
ladder.  A no-ladder run establishes prior-blindness of the mapped number;
it does not replace a native recipient test.

The decisive analysis is a joint likelihood in $(M_B,a_B)$ that tests the
individual shifts in Eq.~\eqref{eq:portal_signed_native_shifts} and the
normalization-independent ratio
Eq.~\eqref{eq:portal_signed_native_ratio}.  Section~\ref{sec:native_joint_likelihood}
provides that likelihood and the accompanying executable implementation.
Section~\ref{sec:estimator_reductions} then performs the dimensionless
reductions for sirens, chronometers, time-delay lenses, and megamasers.  A
publication number still requires the official released arrays and a
separately frozen non-ladder native reference; the paper does not substitute
synthetic values for that empirical step.



% ======================================================================
%  CONCRETE PORTALS, COUNTERTERMS, NATIVE LIKELIHOOD, AND CLOSURE TESTS
% ======================================================================
\section{Concrete Matter Realizations and Endpoint Counterterms}
\label{sec:concrete_matter_realizations}

The abstract matter operators of Sec.~\ref{sec:portals_map} admit a
laboratory realization with existing metrological primitives.  This section
constructs one such realization and classifies the deformations that could
invalidate common-gate universality.  It is a concrete low-energy embedding,
not a claim that the present astronomical apparatus already implements the
coherent defect without further engineering.

\subsection{Temporal portal: a caesium Ramsey phase against the neutral boundary phase}

The SI second is defined by fixing the unperturbed ground-state hyperfine
transition frequency of $^{133}$Cs to
$9\,192\,631\,770\,\mathrm{Hz}$~\cite{BIPMSI2026}.  A concrete charged phase
operator is therefore
\begin{equation}
 Q_T^+
 =|F=4,m_F=0\rangle\langle F=3,m_F=0|\,e^{-i\Theta_K},
 \label{eq:pb_cs_phase_operator}
\end{equation}
where the clock transition is read by a Ramsey sequence and $\Theta_K$ is
the processed-sector phase introduced in
Eq.~\eqref{eq:portal_atomic_clock}.  The measured quantity is the
interferometric phase difference
\begin{equation}
 \Phi_T(\tau)=\arg\langle Q_T^+(\tau)Q_T^-(0)\rangle,
 \qquad
 \Omega_A-\Omega_K=\frac{d\Phi_T}{d\tau}.
 \label{eq:pb_ramsey_phase_difference}
\end{equation}
The memory-neutrality condition is an operator statement,
$[D,\Omega_K]=0$, rather than a coordinate-time convention.  The charged
clock coherence has one endpoint residue, while the processed-sector phase
has none.  The normalized common-gate matrix element therefore gives
$\Delta\ln(\Omega_A/\Omega_K)=+\chi_0$ in the forward orientation.

This construction separates two questions that were previously conflated.
The SI realization identifies a physical local phase standard; the boundary
operator $e^{-i\Theta_K}$ identifies the independent neutral reference.
An ordinary comparison of two equally charged atomic clocks would instead
read only their differential exposure and is the emission-side realization
excluded in Sec.~\ref{sec:emission_atomic}.

\subsection{Radiative portal: SI-traceable optical power against a geometric reference}

Let $b_L$ denote the science optical mode after the standard source
standardization and let $b_G$ denote a geometrically anchored reference mode
propagated through the same local optical train.  The normalized endpoint
transition is
\begin{equation}
 Q_R^-=\frac{b_G^{\dagger}b_L}
 {\sqrt{(N_G+1)N_L}},
 \label{eq:pb_radiometric_transition_concrete}
\end{equation}
which is the operator already used in
Eq.~\eqref{eq:portal_uv_vertex}.  A concrete realization uses electrical
substitution radiometry: a cryogenic radiometer equates absorbed optical
power to calibrated electrical power and supplies an SI-traceable reference
amplitude~\cite{NISTPOWR,Gentile1996}.  A transfer detector may then place
$b_L$ and $b_G$ in the same balanced readout.  In that arrangement a common
detector responsivity $G_{\rm det}$ multiplies both modes and cancels from
\begin{equation}
 \mathcal R_R
 =\left[\frac{\langle b_L^{\dagger}b_L\rangle}
 {\langle b_G^{\dagger}b_G\rangle}\right]^{1/2}.
 \label{eq:pb_concrete_radiometric_ratio}
\end{equation}
The gate therefore acts on a relative local amplitude record, not on stellar
emission and not on photons accumulated along the cosmological line of
sight.

The astronomical application does not require a cryogenic radiometer at
every telescope.  The radiometer supplies a physical realization of the
reference sector and its traceability chain.  Transfer standards and
instrument calibration carry that reference to the detector.  Whether the
actual distance-ladder apparatus couples to the boundary defect is an
empirical portal-selection question; the operator and its null tests are no
longer undefined.

\subsection{Counterterm basis and radiative stability}

Let $\mathcal O_+=\mathcal O_T+\mathcal O_R$ and
$\mathcal O_-=\mathcal O_T-\mathcal O_R$ denote the common and asymmetric
endpoint combinations.  The most general local portal through first order
in the memory and one derivative can be organized as
\begin{multline}
 S_{\rm end}=\int d\tau\,\dot n\bigg[
 g_+\chi\mathcal O_+ +g_-\chi\mathcal O_-
 +\frac{c_+}{\Lambda_*}(u\!\cdot\!\nabla\chi)\mathcal O_+\\
 +\frac{c_-}{\Lambda_*}(u\!\cdot\!\nabla\chi)\mathcal O_-
 +d_+\chi^2\mathcal O_+ +d_-\chi^2\mathcal O_-+\cdots\bigg].
 \label{eq:pb_counterterm_basis}
\end{multline}
The minimal common-gate sector has an endpoint-exchange symmetry
$\mathcal P_{TR}:T\leftrightarrow R$, under which
$\mathcal O_+$ is even and $\mathcal O_-$ is odd.  Consequently $g_-=c_-=d_-=0$
unless an explicit asymmetric spurion is present.  Common wave-function and
vertex renormalizations multiply the numerator and denominator of the
normalized conditional matrix element and cancel.  Descendant-current
mixing is suppressed by $(k+2)^{-1}$, derivative operators by
$H/\Lambda_*$, and nonlinear memory terms by $X_H$.

Writing $\epsilon_-$ for the renormalized endpoint-odd spurion gives the
model-independent estimate
\begin{equation}
 \left|\frac{\eta_T}{\eta_R}-1\right|
 \lesssim 2|\epsilon_-|
 +\mathcal O\!\left(\frac{H}{\Lambda_*},
 \frac{1}{k+2},X_H\right).
 \label{eq:pb_gain_ratio_bound}
\end{equation}
Thus equality of the gains is technically stable when the common-bilocal
exchange symmetry is exact.  A measured violation of
$\Delta M_B=5\Delta a_B$ would measure the asymmetric spurion or reject the
single-defect construction.

\begin{table}[t]
\caption{Endpoint operators through the first nontrivial orders.  ``Cancels''
means that the term is removed from the normalized conditional matrix
element; ``forbidden'' refers to the exact common-gate exchange symmetry.}
\label{tab:pb_counterterms}
\begin{ruledtabular}
\begin{tabular}{@{}p{0.39\columnwidth}p{0.19\columnwidth}p{0.32\columnwidth}@{}}
Operator & Status & Observable effect\\
\hline
common external-leg/vertex factor & cancels & event rate only\\
$\chi\mathcal O_+$ & retained & common gain $g_+$\\
$\chi\mathcal O_-$ & forbidden & unequal gains\\
$(u\cdot\nabla\chi)\mathcal O_+/\Lambda_*$ & suppressed & transient correction\\
current descendants & $\mathcal O[(k+2)^{-1}]$ & finite-level correction\\
$\chi^2\mathcal O_+$ & $\mathcal O(X_H)$ & nonlinear map\\
\end{tabular}
\end{ruledtabular}
\end{table}

\section{Native Joint \texorpdfstring{$(M_B,a_B)$}{(M_B,a_B)} Likelihood}
\label{sec:native_joint_likelihood}

% ============================================================
%  PAPER B: frozen non-ladder posterior + four figures
%  (numbers from chains_test/modelB_A2_noprior_v5_acc, 4 chains,
%   271,137 post-burn-in samples, thinned x677 for the background
%   re-integration; extraction self-certified against the stored
%   H0_local to a maximum relative deviation of 7.7x10^-5)
% ============================================================

\subsection{Frozen non-ladder reference posterior}
\label{sec:pb_frozen_posterior}

All Paper-B statistical statements use one frozen reference: the
Model~B ($\beta = 1/12$, $A_{\rm map} = 2$) posterior for
Planck~2018 (TT,TE,EE $+$ lowE $+$ lensing) $+$ DESI~DR2 BAO $+$
Pantheon$+$ shape $+$ DES~Y3, with no local-ladder likelihood of any
kind.  Each retained sample carries
$(H_{0,\rm phys},\,X_H,\,\Omega_m,\,z_{\rm entry},\,z_{\rm exit})$,
where $X_H = \int \Delta\,d\ln a$ is re-integrated from the interacting
background and $H_{0,\rm phys} = H(0)$ is the computed late-time
expansion rate.  The extraction is certified sample-by-sample against
the pipeline's independently stored mapped value: the maximum relative
deviation of $H_{0,\rm phys}e^{2X_H}$ from the stored $H_0^{\rm local}$
is $7.7\times10^{-5}$.

The frozen posterior gives
\begin{equation}
\langle X_H\rangle = 0.03592,
\qquad
\sigma(X_H) < 1\times10^{-5},
\label{eq:pb_frozen_XH}
\end{equation}
with window crossings $z_{\rm entry} = 0.640 \pm 0.008$ and
$z_{\rm exit} = 0.014 \pm 0.005$, and
$H_{0,\rm phys} = 68.49 \pm 0.24~{\rm km\,s^{-1}\,Mpc^{-1}}$.
The vanishing posterior width of $X_H$ is structural rather than
accidental: in $\Omega_\Lambda$ variables the $\beta$-fixed flow is
autonomous, so the exposure integral is an invariant of the model ---
the window \emph{location} shifts with the sampled cosmology (hence the
finite spread in $z_{\rm entry}$) while the accumulated exposure does
not, to the precision set by trace radiation.  Consequently the
correlation ${\rm Corr}(X_H, H_{0,\rm phys})$ is noise-dominated and
carries no physical content; the prior-blind push-forward
$H_{0,\rm phys}\,e^{2X_H}$ is a rigid rescaling by
$e^{2\langle X_H\rangle} = 1.0745$ that adds no posterior width,
yielding $73.59 \pm 0.26~{\rm km\,s^{-1}\,Mpc^{-1}}$ with no ladder
information used at any stage.

% ---------------- FIGURE 1 ----------------
\begin{figure}[t]
\centering
\includegraphics[width=\columnwidth]{pB_fig1_capacity_history.pdf}
\caption{Capacity history over the frozen non-ladder posterior
(16--84\% bands over 401 thinned samples).  \emph{Top:} the throughput
ratio $\mathcal S(z)$, crossing unity at
$z_{\rm entry} = 0.640 \pm 0.008$ and $z_{\rm exit} = 0.014 \pm 0.005$
(the $\Omega_\Lambda = \tfrac13$ and $\tfrac23$ roots).
\emph{Middle:} the excess fraction $\Delta(z)$.
\emph{Bottom:} the cumulative exposure $X_H(<z)$; its band is
unresolvable at plot precision because $X_H$ is an invariant of the
$\beta = 1/12$ flow (Eq.~\ref{eq:pb_frozen_XH}).}
\label{fig:pb_capacity_history}
\end{figure}

% ---------------- FIGURE 2 ----------------
\begin{figure}[t]
\centering
\includegraphics[width=\columnwidth]{pB_fig2_pushforward.pdf}
\caption{The frozen joint posterior.  \emph{Left:}
$(H_{0,\rm phys}, X_H)$ samples; the exposure is posterior-rigid while
$H_{0,\rm phys}$ carries the full cosmological width.  \emph{Right:}
the prior-blind push-forward $H_{0,\rm phys}e^{2X_H}$ (no ladder data
used), landing at $73.59 \pm 0.26~{\rm km\,s^{-1}\,Mpc^{-1}}$ as a pure
prediction of the common-gate map.}
\label{fig:pb_pushforward}
\end{figure}

% ---------------- FIGURE 3 ----------------
\begin{figure}[t]
\centering
\includegraphics[width=\columnwidth]{pB_fig3_native_plane.pdf}
\caption{The native calibration plane $(5\Delta a_B, \Delta M_B)$.
The common gate populates only the symmetric direction
$\Delta M_B = 5\Delta a_B$ (solid), with the frozen prediction at
$(0.07796, 0.07796)$~mag for $X_H = 0.0359$.  Conventional photometric
zero-point perturbations excite the antisymmetric null direction
(dashed) and cancel in $\ln H_0$.  The discriminant
$D_{\rm com} = \Delta M_B - 5\Delta a_B$ measures displacement off the
gate line; the common-gate realization predicts $D_{\rm com} = 0$
exactly.}
\label{fig:pb_native_plane}
\end{figure}



The public Pantheon+/SH0ES release provides the simultaneous ladder as a
linear system with data vector $Y$, design matrix $L$, and covariance $C$
\cite{PantheonDataRelease,Riess2022}.  In the release convention,
\begin{equation}
 Y=\theta L+\varepsilon,
 \qquad \langle\varepsilon\varepsilon^T\rangle=C.
 \label{eq:pb_shoes_linear_system}
\end{equation}
The generalized least-squares solution and covariance are
\begin{align}
 F&=LC^{-1}L^T,
 \label{eq:pb_shoes_fisher}\\
 \widehat\theta&=F^{-1}LC^{-1}Y,
 \qquad
 \mathrm{Cov}(\widehat\theta)=F^{-1}.
 \label{eq:pb_shoes_gls}
\end{align}
Let $h_5\equiv5\log_{10}H_0$.  The native transformation is
\begin{equation}
 z\equiv
 \begin{pmatrix}M_B\\ a_B\end{pmatrix},
 \qquad
 a_B=\frac{h_5-M_B-25}{5},
 \label{eq:pb_native_transform}
\end{equation}
with covariance obtained from the Jacobian
$J=\bigl(\begin{smallmatrix}1&0\\-1/5&1/5\end{smallmatrix}\bigr)$.

Let $z_{\rm phys}$ and $C_{\rm phys}$ be a native reference frozen from the
non-ladder analysis before the local data are inspected, and define
\begin{equation}
 \delta z=\widehat z-z_{\rm phys},
 \qquad
 C_\delta=C_{\widehat z}+C_{\rm phys}.
 \label{eq:pb_native_delta}
\end{equation}
For the common gate,
\begin{equation}
 s\equiv\frac{1}{\ln10}
 \begin{pmatrix}5\\1\end{pmatrix},
 \qquad \delta z=X_Hs.
 \label{eq:pb_native_shift_vector}
\end{equation}
The minimum-variance native estimate of the exposure is therefore
\begin{equation}
 \widehat X
 =\frac{s^TC_\delta^{-1}\delta z}
 {s^TC_\delta^{-1}s},
 \qquad
 \sigma_X^2=(s^TC_\delta^{-1}s)^{-1}.
 \label{eq:pb_native_xhat}
\end{equation}
The normalization-independent common-gate null is
\begin{equation}
 D_{\rm com}\equiv
 \begin{pmatrix}1&-5\end{pmatrix}\delta z
 =\Delta M_B-5\Delta a_B,
 \label{eq:pb_native_D}
\end{equation}
with
\begin{equation}
 \chi^2_\perp
 =\frac{D_{\rm com}^2}
 {(1,-5)C_\delta(1,-5)^T}.
 \label{eq:pb_native_chi2_perp}
\end{equation}
This statistic tests the gain ratio without using the predicted magnitude of
$X_H$.  The component along $s$ tests the amplitude after the non-ladder
exposure posterior is frozen.

An executable implementation accompanies this paper.  It reads the official
$Y/L/C$ products, performs a Cholesky-stable GLS solve, transforms the
covariance to $(M_B,a_B)$, evaluates
Eqs.~\eqref{eq:pb_native_xhat}--\eqref{eq:pb_native_chi2_perp}, and includes
a synthetic exact-recovery test.  No numerical native result is quoted here
because the official covariance files and an independently frozen native
reference must be supplied to the executable; placeholder values are not
substituted for an empirical analysis.

\begin{proposition}[Native statistical implementation]
Equations~\eqref{eq:pb_shoes_linear_system}--\eqref{eq:pb_native_chi2_perp}
provide a complete likelihood-level test of the two endpoints.  The
compressed $H_0$ likelihood constrains only the projection along $s$, while
$D_{\rm com}$ separately tests common-gate universality.
\end{proposition}

\section{Estimator-Level Register Reductions}
\label{sec:estimator_reductions}

The gauge-neutral minimal portal has an important selection consequence.  A
single charged endpoint is not a physical observable; the lowest-dimensional
invariant object is the full bilocal.  Therefore the baseline method classes
are neutral $(w_T,w_R)=(0,0)$ or full-gate $(1,1)$.  An isolated one-unit
class requires a different gauge-completed defect and is an extension of the
minimal theory, not an automatic consequence of using time or distance.

\subsection{Electromagnetic supernova ladder}

The geometric-anchor/primary-calibrator/SN hierarchy prepares one shared
absolute radiative and phase standard, so the common bilocal is engaged once:
\begin{equation}
 (w_T,w_R)_{\rm SN}=(1,1),
 \qquad A_{\rm SN}=2.
 \label{eq:pb_sn_weights}
\end{equation}
Adding anchors or subdividing the calibrator chain does not change this
intersection number.  A direct geometric distance to an SN host also engages
the same complete absolute-standard transfer and is therefore $A=2$ in the
minimal common-gate theory, rather than an intermediate one-edge rung.
Pantheon-like Hubble-diagram shape information with $M_B$ marginalized is
neutral.

\subsection{Gravitational-wave standard sirens}

A standard siren obtains luminosity distance from the dimensionless strain
and waveform phasing and combines it with a host or statistical redshift;
it does not use the electromagnetic absolute-luminosity transfer
\cite{Schutz1986,Abbott2017}.  Detector length and time standards enter the
calibration of one dimensionless strain channel and do not form the
$Q_T^+WQ_R^-$ bilocal.  The minimal assignment is therefore
\begin{equation}
 (w_T,w_R)_{\rm GW}=(0,0),
 \qquad A_{\rm GW}=0.
 \label{eq:pb_gw_weights}
\end{equation}
A future experiment that uses a siren distance to calibrate an SN absolute
standard would be a different composite path: the siren route itself is
neutral, while the subsequent SN transfer crosses the full common gate once.

\subsection{Cosmic chronometers}

Cosmic chronometers estimate
\begin{equation}
 H(z)=-\frac{1}{1+z}\frac{dz}{dt_*},
 \label{eq:pb_cc_estimator}
\end{equation}
where $dt_*$ is inferred from differential stellar-population ages rather
than from the local Ramsey--Kodama comparator
\cite{JimenezLoeb2002,Moresco2012}.  The age-scale and redshift calibrations
occur within one stellar-population model and do not instantiate the
radiative endpoint.  Hence the strict minimal assignment is
\begin{equation}
 (w_T,w_R)_{\rm CC}=(0,0).
 \label{eq:pb_cc_weights}
\end{equation}
A charged stellar-age portal is possible, but it is an additional endpoint
model and predicts specific SPS and population-dependent residuals.

\subsection{Time-delay cosmography}

Strong-lens time delays determine
\begin{equation}
 D_{\Delta t}=\frac{c\,\Delta t}{\Delta\phi},
 \label{eq:pb_timedelay_estimator}
\end{equation}
where $\Delta\phi$ is the dimensionless Fermat-potential difference
\cite{Birrer2024}.  A common local time-unit rescaling is accompanied by the
SI length realization through the exact value of $c$ and cancels from the
dimensionless raw likelihood.  No electromagnetic luminosity standard is
transported.  Thus
\begin{equation}
 (w_T,w_R)_{\Delta t}=(0,0)
 \label{eq:pb_timedelay_weights}
\end{equation}
in the minimal portal sector.

\subsection{Megamaser distances}

Disk megamasers infer a geometric length from angular positions, line-of-
sight velocities, and accelerations.  Schematically,
\begin{equation}
 D\propto\frac{v^2}{a\,\theta}.
 \label{eq:pb_maser_estimator}
\end{equation}
Under a common time-standard rescaling,
$v\rightarrow e^{-\xi}v$ and $a\rightarrow e^{-2\xi}a$, so $v^2/a$ is
invariant.  The Hubble-flow use of the distance then combines it with a
dimensionless redshift; it does not cross the common bilocal
\cite{Herrnstein1999,Pesce2020}.  Therefore
\begin{equation}
 (w_T,w_R)_{\rm maser}=(0,0).
 \label{eq:pb_maser_weights}
\end{equation}

\begin{table*}[t]
\caption{Estimator-level register reductions in the minimal gauge-neutral
common-gate theory.  ``Gate incidence'' counts physical bilocal crossings,
not appearances of dimensional quantities in an analysis formula.}
\label{tab:pb_estimator_reductions}
\begin{ruledtabular}
\begin{tabular}{lllll}
Method & Dimensionless/raw core & Cancellation or selection & Gate incidence & $(w_T,w_R)$\\
\hline
SN absolute ladder & calibrated flux + Hubble-flow intercept & one shared absolute-standard state & one & $(1,1)$\\
SN shape only & relative magnitudes, $M_B$ marginalized & absolute mode projected out & zero & $(0,0)$\\
GW standard siren & strain waveform + redshift & no EM luminosity transfer & zero & $(0,0)$\\
cosmic chronometer & $-(1+z)^{-1}dz/dt_*$ & differential SPS age scale & zero & $(0,0)$\\
time-delay lens & $c\Delta t/\Delta\phi$ & common time/length traceability & zero & $(0,0)$\\
megamaser & $v^2/(a\theta)$ & common time powers cancel & zero & $(0,0)$\\
\end{tabular}
\end{ruledtabular}
\end{table*}

The table is a falsifiable selection rule of the minimal portal.  A robust
one-unit result would not be absorbed by retuning $X_H$; it would require an
additional gauge-completed endpoint representation.

\section{Pre-Registered Unequal-Path Closure Experiment}
\label{sec:closure_experiment}

% ---------------- FIGURE 4 ----------------
\begin{figure}[t]
\centering
\includegraphics[width=\columnwidth]{pB_fig4_closure_loop.pdf}
\caption{The pre-registered unequal-path closure loop.  A bright-siren
distance (gravitational-wave amplitude, $\Delta\mathbf W = (0,0)$: no
standardization-sector crossing) and a ladder SN distance
($\Delta\mathbf W = (1,1)$: both endpoints of the common gate engaged)
share endpoints and large-scale environment; node zero points
telescope.  The predicted holonomy is
$\mathcal H_{\rm GW\text{-}SN} = 2X_H = 7.44\%$ in distance
($0.15591$~mag-equivalent); a null result rejects occurrence of the
common gate on the calibration path.}
\label{fig:pb_closure_loop}
\end{figure}


A same-sensitivity null can test systematics but cannot distinguish the
common gate from a neutral theory.  The cleanest unequal-path experiment is
a bright-siren--SN-host closure.

\subsection{Two paths in the same volume}

Select a set of low-redshift binary-neutron-star or neutron-star--black-hole
events with identified hosts.  Require either an SN~Ia in the same host or a
pre-registered host-matched SN sample with the same peculiar-velocity
volume.  Construct two independently analyzed paths:
\begin{align}
 p_{\rm N}:&\quad
 \text{GW strain}+z_{\rm host}
 \longrightarrow H_0^{\rm GW},
 \quad A_{\rm N}=0,
 \label{eq:pb_closure_neutral_path}\\
 p_{\rm L}:&\quad
 d_L^{\rm GW}\longrightarrow M_B^{\rm SN}
 \longrightarrow a_B^{\rm flow},\notag\\
 &\hspace{2.3cm}\longrightarrow H_0^{\rm SN|GW},
 \quad A_{\rm L}=2.
 \label{eq:pb_closure_ladder_path}
\end{align}
The neutral path never establishes an electromagnetic absolute luminosity
standard.  The second path uses the same geometric siren distance to prepare
the SN standard and therefore crosses the full common gate once.

The pre-registered loop statistic is
\begin{equation}
 \mathcal H_{\rm GW-SN}
 \equiv\ln\frac{H_0^{\rm SN|GW}}{H_0^{\rm GW}}
 =2X_H.
 \label{eq:pb_gw_sn_holonomy}
\end{equation}
For $X_H=0.0359$,
\begin{equation}
 \mathcal H_{\rm GW-SN}=0.0718,
 \qquad
 \frac{H_0^{\rm SN|GW}}{H_0^{\rm GW}}=1.07444,
 \label{eq:pb_gw_sn_numeric}
\end{equation}
which is $0.15591$ mag-equivalent.  Reversing the two routes must change the
sign and invert the ratio.

\subsection{Protocol and precision target}

The protocol should be frozen before unblinding:
\begin{enumerate}
\item fix the redshift and host-quality cuts, preferably within the same
$0.023\lesssim z\lesssim0.15$ flow volume;
\item freeze the non-ladder posterior for $X_H$ and do not refit it to the
closure data;
\item use independent GW and SN calibration teams and blind the relative
zero point;
\item propagate the shared peculiar-velocity covariance rather than adding
the two route errors as independent;
\item pre-specify dust, host-population, inclination, and selection controls;
\item require the internal common-gate null
$\Delta M_B-5\Delta a_B=0$ in the SN path;
\item verify the inverse-path sign and a same-sensitivity null control.
\end{enumerate}
To reject the neutral loop at $5\sigma$, the uncertainty in
$\mathcal H_{\rm GW-SN}$ must be below
\begin{equation}
 \sigma_{\mathcal H}<\frac{2X_H}{5}=0.01436,
 \label{eq:pb_closure_precision}
\end{equation}
corresponding to $1.44\%$ in the ratio.  A $3\sigma$ test requires $2.39\%$.
As an orientation estimate, independent $10\%$ siren distances would require
roughly $(0.10/0.01436)^2\simeq49$ effective events if the SN-route error
were negligible, and of order one hundred if the two paths contribute
comparable variance.  The exact forecast must include inclination and
peculiar-velocity covariance.

\begin{proposition}[Convention-independent closure test]
The statistic in Eq.~\eqref{eq:pb_gw_sn_holonomy} compares two physical
paths with identical endpoints and common large-scale environment.  Node
zero points telescope, common peculiar-velocity modes can be included in the
joint covariance, and the minimal theory predicts a fixed two-unit
nonclosure.  A result consistent with zero rejects the occurrence of the
common gate on the SN path; a one-unit result selects an additional endpoint
sector not present in the minimal bilocal theory.
\end{proposition}



% ======================================================================
%  DISCUSSION OF STATUS AND COMPLETION
% ======================================================================
\section{Scope of the Completion and Excluded Realizations}
\label{sec:completion_tasks}

The construction is modular.  The processed screening sector does not
require the over-capacity recipient, while the endpoint sector can be
judged independently of the cosmological growth fit.  Within the explicit
source-reduced comparator EFT, the operational chain is now closed:
\begin{equation}
\begin{aligned}
 \Delta(\mathcal S)&\longrightarrow X_H=\chi_0
 \longrightarrow\{\mathcal R_T,\mathcal R_R\}\\
 &\longrightarrow\{a_B,M_B\}
 \longrightarrow H_{0,\rm ladder}.
\end{aligned}
 \label{eq:pb_completed_chain}
\end{equation}
The write instrument fixes the stored memory; the affine Ward identity fixes
one unit at each endpoint; the unique minimal gauge-neutral bilocal fixes a common gate;
the global ladder topology gives one serial engagement; and the retarded
passive branch fixes the sign.  No continuous mapping coefficient remains
in this minimal sector.

The five construction tasks previously left open now have explicit minimal
realizations.  The processed/excess instrument is generated by the
Hamiltonian dilation of Sec.~\ref{sec:exact_hamiltonian_dilation}; the
caesium phase and SI-radiometric endpoints and their counterterms are given
in Sec.~\ref{sec:concrete_matter_realizations}; the public linear-system
likelihood and null statistic are given in
Sec.~\ref{sec:native_joint_likelihood}; non-ladder methods are reduced in
Sec.~\ref{sec:estimator_reductions}; and
Sec.~\ref{sec:closure_experiment} specifies a concrete unequal-path loop.
The remaining frontier is therefore empirical and foundational rather than
an unspecified formula:
\begin{enumerate}
\item run the native likelihood with the official released covariance and a
pre-frozen non-ladder native reference;
\item build or identify an apparatus in which the proposed matter endpoints
couple to the common boundary defect, and bound the asymmetric spurion
$\epsilon_-$;
\item execute the bright-siren--SN closure experiment;
\item determine whether the full non-perturbative Standard Model plus
CS/WZW theory uniquely selects the minimal source-reduced dilation rather
than another symmetry-allowed universality class.
\end{enumerate}
These are tests of whether nature realizes the constructed model.  They do
not reintroduce a free write coefficient, endpoint gain ratio, engagement
count, sign, likelihood definition, or method-selection rule into the
minimal EFT.

Three alternative realizations remain excluded from the baseline.  First,
a universal source-local luminosity change calibrates out by
Eq.~\eqref{eq:pb_source_cancel}.  Second, an accumulated photon optical depth
would produce a redshift-dependent Hubble residual, distance-duality
violation, or colour dependence.  Third, a strictly homogeneous scalar
metric perturbation can be absorbed into a redefinition of time and scale
factor, while a genuine bulk modification would alter BAO and supernova
shape observables.  The historical relation $\Phi=\Psi$ is therefore only
geometric motivation for equal temporal and spatial responses; in the
present theory their equality follows from equal endpoint weights under one
common gate.

The previous area-growth ODE is not discarded but reclassified.  Since
$r_A=c/H$,
\begin{equation}
1+q=\frac{d\ln r_A}{d\ln a}
=\frac12\frac{d\ln A}{d\ln a},
\end{equation}
so
\begin{equation}
d\chi_A=\Delta\,d\ln r_A
=\frac12\Delta\,d\ln A.
\label{eq:pb_area_accumulator}
\end{equation}
It is an area-growth accumulator, whereas $X_H$ is a capacity-cycle
accumulator.  On the flat hard-gate benchmark,
$X_H=(4/3)\chi_A$.  They are different functionals, not rival predictions
for one observable.


\section{Discussion}
\label{sec:discussion}


\subsection{One boundary, two responses}
\label{sec:one_boundary_two_responses}

The screening and memory sectors should not be identified with one another.
The processed sector is a continuous, equilibrated response that enters the
background and perturbation equations through $Q$.  The over-capacity sector
is a history of the rejected fraction and enters only through a later
relative-standard readout.  Their common origin is the same driven boundary,
but their clocks and observables differ.  In particular, the factor
$\beta=1/12$ belongs to the processed Pendry throughput; it is not
automatically a factor in the capacity-cycle memory.  Conversely, the
common-gate endpoint weight two does not alter the dark-sector growth
calculation.

This separation is experimentally useful.  Failure of the calibration
recipient leaves the screening and growth predictions intact.  Failure of
the screening polynomial leaves the calibration holonomy without a derived
processed interaction.  Paper A and the present paper can therefore be
tested independently while sharing the same apparent-horizon drive.



\subsection{Relation to existing LQG horizon frameworks}
\label{sec:lqg_context}

The ABCK framework~\cite{Ashtekar1998,Ashtekar2000} was developed
for \emph{isolated} horizons: stationary boundaries with fixed
area and therefore fixed Chern-Simons level $k$.  The entropy
calculation proceeds by counting the number of puncture
configurations consistent with a given total area, yielding
$S = A/(4\ell_P^2)$ for an appropriate choice of $\gamma$.
Subsequent refinements by Engle, Noui, and
Perez~\cite{Engle2010} and by Ghosh and Perez~\cite{Ghosh2011}
incorporated the SU(2) structure more carefully, showing that
the entropy counting is robust across different quantisation
prescriptions and does not require a fixed $\gamma$.

The present work extends this framework in a specific direction:
from isolated horizons with static $k$ to the cosmological
apparent horizon with dynamical $k$.  The extension is physically
motivated by the fact that the apparent horizon in FLRW spacetime
is not an isolated horizon---its area changes with the expansion
history---and technically enabled by the observation that the
Chern-Simons formulation remains well-defined for time-dependent
$k$, since the level enters the symplectic structure as a
parameter rather than a dynamical variable in the boundary
theory's own equations of motion.

What is new is not the Chern-Simons description itself but
the \emph{dynamics} it acquires on the cosmological horizon.
The growth rate $dk/d\ln a = 2(1+q)k$ generates a continuous
flow of new punctures that competes with the existing puncture
population for the horizon's processing capacity.  This
competition---absent for isolated horizons, where $dk/dt = 0$
and all punctures are in thermal equilibrium---is the microscopic
origin of the capacity-sharing mechanism that produces
$I_{\rm eff} = (1-q)/2$.  In this sense, the thermodynamic
modulation is a genuinely cosmological effect: it exists because
the horizon is expanding.  In the stationary (de~Sitter) limit,
$I_{\rm eff} \to 1$ (full saturation, not zero), so the modulation
itself does not vanish---rather, the \emph{competition between
emergence and interaction} vanishes ($\Gamma_{\rm emerge} \to 0$)
and the isolated-horizon framework is recovered with the
interaction fully active.


\subsection{Why the mean-field prescription is selected}
\label{sec:mean_field_discussion}

Section~\ref{sec:mean_field} showed that the exact
throughput-weighted average
$\langle x(j) \rangle_{\rm throughput} = I_{\rm eff}$ recovers
the capacity fraction trivially, while the mean-field evaluation
$x(\langle j \rangle) = (1-q)(2-q)/6$ encodes the additional
degeneracy factor $\langle d \rangle = 2-q$.  The phenomenological comparison favours the
mean-field closure: the Lagrangian modulation produces
$\sigma_8 = 0.678$, the throughput-weighted average produces no
correction beyond $I_{\rm eff}$ itself (leaving the
over-suppression unresolved), and only the mean-field screening
restores
$\sigma_8 = 0.756$, and close to the low-redshift clustering amplitude used in the
companion analysis~\cite{Sayegh2026}.

The coarse-grained order-parameter argument of
Section~\ref{sec:mean_field} provides the microscopic reason for
this selection.  The dark drag is a homogeneous background flux
($k = 0$ spatial mode) that couples to the spatially averaged
boundary connection $\overline{A}^{\,i}$, not to individual
puncture fluctuations.  The macroscopic response is therefore
governed by the quantum effective action for this coarse-grained
order parameter, yielding the mean-field rate
$\Gamma(\langle j \rangle)$ rather than the puncture average
$\langle \Gamma(j) \rangle$.  The thermodynamic limit
($N \sim 10^{122}$, monopole fluctuations suppressed by
$\mathcal{O}(N^{-1/2}) \sim 10^{-61}$) renders fluctuations
negligible for observational purposes.


\subsection{The linear-source universality class}
\label{sec:universality}

The structural protection of the saddle
point~(Section~\ref{sec:partition}) defines what we term the
\emph{linear-source universality class}: the set of sourced
Chern-Simons theories in which the junction-condition coupling is
linear in the connection, so that the Schur reduction yields the
clean factorisation $w_j = \eta\,\lambda_j$ with no residual
$j$-dependent corrections.

Within this class, the saddle-point allocation $I_{\rm eff}$
depends only on the kinematic ratio
$\lambda_1/\lambda_{1/2} = (1-q)/(1+q)$ and is insensitive to:
\begin{itemize}
\item the absolute value of the Barbero--Immirzi parameter
  $\gamma$;
\item quantum corrections to the Chern-Simons level $k$;
\item operator-ordering ambiguities in the quantisation of the
  boundary theory;
\item coarse-graining or renormalisation-group flow of the
  boundary couplings;
\item the precise value of the horizon temperature $T_{\rm KH}$.
\end{itemize}
Departure from this class would require nonlinear source terms
with $C_2(j)$-dependent corrections~\eqref{eq:w_with_C2},
suppressed by $(\ell_P/r_A)^2 \sim 10^{-122}$.  The universality
is therefore not merely a technical convenience but a physical
prediction: any microscopic theory of the cosmological horizon
that belongs to this class---whether loop quantum gravity,
spin-foam models, or group field theory---must reproduce the same
$I_{\rm eff} = (1-q)/2$ and the same screening polynomial
$x = \langle j \rangle(2\langle j \rangle + 1)/3$.

The class also admits a sharp boundary.  If one could engineer a
horizon (or a laboratory analogue) where the source coupling is
quadratic in the connection---e.g.\ through strong gravitational
self-interaction at Planckian curvature---the $C_2(j)$ terms would
become significant, the $d_j$ cancellation would fail, and the
saddle point would shift.  This provides a concrete criterion for
when the holographic equilibrium framework would break down: at
energy scales where the Israel--Darmois junction conditions receive
$\mathcal{O}(1)$ nonlinear corrections.


\subsection{The vacuum energy hierarchy as a kinematic identity}
\label{sec:cc_problem}

The traditional cosmological constant problem asks why
$\rho_\Lambda/\rho_{\rm Pl} \sim 10^{-122}$, where
$\rho_{\rm Pl} = c^7/(\hbar G^2)$.  The standard framing treats
this as a fine-tuning anomaly.  However, it assumes the bulk
degrees of freedom are unbounded.  The holographic principle
dictates that the state count of a region is bounded by its
boundary area: $N = A/(4\ell_P^2) = \pi r_A^2/\ell_P^2$
(a coarse-grained surface DOF count, not a literal LQG
puncture number; see Section~\ref{sec:emergence_channels}).

The hierarchy is not a fine-tuned constant; it is a strict
kinematic identity of the Friedmann equations.  The critical
density is $\rho_c = 3H^2/(8\pi G)$.  Substituting
$r_A = c/H$ and $\ell_P^2 = \hbar G/c^3$, and using
$(\ell_P/r_A)^2 = \pi/N$:
\begin{equation}
\frac{\rho_c}{\rho_{\rm Pl}}
  = \frac{3}{8\pi}\left(\frac{\ell_P}{r_A}\right)^2
  = \frac{3}{8N}\,.
\label{eq:rho_c_N}
\end{equation}
Because the vacuum energy is a fraction of the critical density
($\rho_\Lambda = \Omega_\Lambda\,\rho_c$), its exact value
relative to the Planck scale is
\begin{equation}
\frac{\rho_\Lambda}{\rho_{\rm Pl}}
  = \frac{3\,\Omega_\Lambda}{8N}\,.
\label{eq:pure_kinematic_hierarchy}
\end{equation}
This is a pure kinematic identity.  It requires no inequalities,
no UV cutoffs, and no black hole avoidance bounds.  (the companion cosmology paper~\cite{Sayegh2026}
additionally derives the exact inequality form: the Misner--Sharp
processing bound is equivalent to
$\rho_\Lambda/\rho_{\rm Pl} \le 1/[12N(1-\Omega_\Lambda)]
\Leftrightarrow \mathcal{S} \le 1$, saturated at the roots
$\Omega_\Lambda \in \{1/3,\,2/3\}$; the identity here is its
saturation value.)  The
suppression factor $10^{-122}$ is the statement that
$r_A \sim 10^{61}\,\ell_P$---the apparent horizon is
$\sim 10^{61}$ Planck lengths across.  In the quantum-gravitational description
developed above, the same coarse-grained count $N$ sets the
dimension of the boundary Hilbert space and enters the
partition function~\eqref{eq:Z}; the literal LQG puncture number
is $j$-label-dependent and differs by an $\mathcal{O}(1)$ factor,
but the Chern-Simons level $k \propto N/\gamma$ controls both.  The suppression is
controlled by the Chern-Simons level: $N \sim k\gamma$, so
$1/N \sim 1/k$---the same control parameter that bounds the
quadratic correction $\epsilon_2$
in~\eqref{eq:eps2_bound}.  The cosmological constant problem and
the structural protection of $I_{\rm eff}$ share a common origin
in the finiteness of the boundary Hilbert space.

The deep structure emerges from the parallel between two
dimensionless kinematic ratios:
\begin{itemize}
\item \emph{The scale of the vacuum:}
  $1/N \sim (\ell_P/r_A)^2$.  This governs the energy hierarchy.
\item \emph{The rate of evolution:}
  $\mathcal{S}(z) = (9/2)\,\Omega_\Lambda\,\Omega_m$.  This
  governs the thermodynamic drag.
\end{itemize}
These two identities are connected by the late-time attractor.
The sweep ratio $\mathcal{S}(z)$ drives the universe toward the
saturation fixed point at $\Omega_\Lambda = 2/3$
(Section~\ref{sec:kernel_recap}).  At the attractor,
equation~\eqref{eq:pure_kinematic_hierarchy} simplifies to
\begin{equation}
\frac{\rho_\Lambda}{\rho_{\rm Pl}}
  = \frac{1}{4N}\,.
\label{eq:hierarchy_attractor}
\end{equation}
The logic here inverts the traditional presentation.
the companion cosmology paper~\cite{Sayegh2026} derived the super-Schwarzschild window
by \emph{setting} the sweep rate equal to $c^3/2G$ and solving
for the roots.  But $\mathcal{S}(z)$ is a kinematic identity that
requires no bound: it is simply $|\dot{q}|/H$.  The drag
$Q \propto \mathcal{S}$ drives $\Omega_\Lambda$ toward $2/3$.
At the attractor, $\mathcal{S} = 1$, and the sweep rate
\emph{evaluates to}
\begin{equation}
\dot{M}_\Lambda\big|_{\rm att}
  = \tfrac{9}{4}\cdot\tfrac{2}{3}\cdot\tfrac{1}{3}
    \cdot\frac{c^3}{G}
  = \frac{c^3}{2G}\,.
\label{eq:schwarzschild_emergent}
\end{equation}
The Schwarzschild rate is not an input; it is the equilibrium
throughput at the fixed point.  The channel structure
(Section~\ref{sec:pendry}) then answers a different question:
\emph{why} is $\mathcal{S} = 1$ the attractor?  Because the $N$
channels each carry a per-channel capacity of one Gibbons--Hawking
quantum per Hubble time, $k_B T_{\rm dS}\,H$
(the companion cosmology paper~\cite{Sayegh2026}), and these aggregate to exactly
$c^5/2G$---the Misner--Sharp processing bound.  The attractor
exists where the macroscopic sweep demand matches this aggregate
geometric capacity.  (The Heisenberg-saturated \emph{interaction}
throughput is the fraction $\beta = 1/12$ of it---the same split
that fixes the coupling.)

The vacuum energy is therefore not a fine-tuned parameter set at
the Big Bang; it is the dynamical fixed point of
$\mathcal{S}(z)$.  The holographic
interaction provides the feedback mechanism: the dark drag
transfers energy from the vacuum to cold dark matter, throttling
$\Omega_\Lambda$ toward the fixed point.  The scaling $\rho_\Lambda \sim \rho_{\rm Pl}/N$ has been noted
in the context of the CKN bound~\cite{CKN1999} and holographic
vacuum energy stabilisation~\cite{Thomas2000}.
the companion cosmology paper~\cite{Sayegh2026} went further, deriving the $1/N$
suppression from throughput saturation and identifying the
attractor at $\Omega_\Lambda = 2/3$ as the dynamical mechanism
that promotes the holographic inequality to an equality.  The
advance here is that the $1/N$ hierarchy emerges as a kinematic
identity of the Friedmann equations
(equation~\ref{eq:pure_kinematic_hierarchy}), requiring no
bounding argument, and that the Schwarzschild rate $c^3/2G$ is
revealed as the equilibrium throughput at the fixed point
(equation~\ref{eq:schwarzschild_emergent}) rather than an
externally imposed limit.

The identity~\eqref{eq:pure_kinematic_hierarchy} reframes the
cosmological constant as an infrared boundary condition---the
inverse of the horizon's state count---rather than an ultraviolet
fine-tuning of local field couplings.  The absolute value of
$\rho_\Lambda$ is set by the size of the universe in Planck
units; the question of why this size takes its observed value
remains open, but it is now a question about the global boundary,
not about local vacuum fluctuations.

Unlike holographic dark energy models that identify $L$ with the
future event horizon (which requires
acausality~\cite{Hsu2004,Li2004}; see~\cite{WangWangLi2017} for a review), this framework uses the
strictly causal apparent horizon $r_A = c/H$.  The drag preserves
the baseline equation of state $w = -1$ while producing a
fit-equivalent CPL surrogate
$(w_0, w_a) \approx (-0.995, -0.023)$~\cite{Sayegh2026}.


\subsection{Channel complementarity and the Pendry bound}
\label{sec:pendry}

The canonical structure of the boundary theory predicts that the
macroscopic interaction strength
$\beta = 1/12$~\cite{Sayegh2026} originates from the Heisenberg
uncertainty principle on the horizon phase space.

On the boundary, the derivation of $\beta$ uses the bulk
LQG canonical pair $(A, \Sigma)$ rather than the boundary CS
symplectic structure $\Omega_\Delta \propto k\int\delta A\wedge\delta A$
used in Section~\ref{sec:kms}.  This is a deliberate change of
variables: in the semiclassical regime $k \gg 1$, the CS
symplectic form reduces to the LQG canonical pair, so both
descriptions are equivalent at leading order.  Concretely, the
SU(2) connection $A$ and the densitised triad $\Sigma$ satisfy
\begin{equation}
[\hat{A},\,\hat{\Sigma}] = i\hbar\,,
\label{eq:uncertainty_AS}
\end{equation}
leading to the standard uncertainty principle
$\Delta A \cdot \Delta\Sigma \geq \hbar/2$.

We map this geometric uncertainty to a thermodynamic energy-time
uncertainty using the Hamiltonian dynamics of the horizon.  The
energy of the apparent horizon scales linearly with its area,
which is generated by the triad flux $\Sigma$.  The local energy
variation associated with a puncture is therefore
$\Delta E \propto \mathcal{H}\,\Delta\Sigma$, where
$\mathcal{H}$ is the characteristic frequency of the boundary.
Hamilton's equations give
$\dot{A} = \partial E/\partial\Sigma = \mathcal{H}$, so for a
state transition occurring over a time $\Delta t$ the connection
shifts by $\Delta A = \mathcal{H}\,\Delta t$.  Substituting into
the uncertainty relation, the frequency $\mathcal{H}$ cancels
exactly:
\begin{equation}
(\mathcal{H}\,\Delta t)
  \left(\frac{\Delta E}{\mathcal{H}}\right) \geq \frac{\hbar}{2}
  \quad\Longrightarrow\quad
  \Delta E \cdot \Delta t \geq \frac{\hbar}{2}\,.
\label{eq:mandelstam_tamm}
\end{equation}
This is the Mandelstam--Tamm quantum speed limit.  The
cancellation of $\mathcal{H}$ is significant: the result depends
only on the symplectic structure, not on the specific dynamics of
the boundary theory.  It defines the minimum time
$\Delta t_{\rm min} = \hbar/(2\Delta E)$ required for a quantum
channel to process an energy fluctuation $\Delta E$, giving a
maximum power per channel of
\begin{equation}
\dot{E}_{\rm max}
  \approx \frac{\Delta E}{\Delta t_{\rm min}}
  = \frac{2}{\hbar}\,(\Delta E)^2\,.
\label{eq:power_limit}
\end{equation}

Each puncture constitutes an effective $(1{+}1)$-dimensional
quantum channel through the horizon via its boundary edge modes.
Although the bulk Chern-Simons theory has no local propagating
degrees of freedom (Section~\ref{sec:linearity}), the
punctures support boundary edge modes that do propagate along
the $(1{+}1)$-dimensional intersection of the spin-network edge
with the horizon; these are the physical channels that carry
energy.  Because the apparent horizon boundary is in KMS
equilibrium (at the Kodama--Hayward temperature $T_{\rm KH}$),
the energy fluctuations of a single edge-mode channel are
thermally distributed.  The coupling \emph{constant} $\beta$ is
defined at the de~Sitter fixed point, where
$T_{\rm KH} = T_{\rm dS} = \hbar H/(2\pi k_B)$; away from the
fixed point the thermodynamic state enters \emph{once}, through
the channel allocation $I_{\rm eff}$ of the saddle
point~\eqref{eq:saddle}, not through the per-channel spectrum.
(Evaluating the quadratic law below at $T_{\rm KH}$ instead would
impose the modulation $[(1-q)/2]^2$, double-counting the
$(1-q)/2$ already supplied by the channel allocation; internal
consistency of the two routes therefore requires the fixed-point
evaluation.  The MCMC consistency of the linear form,
$\beta_{\rm obs} = 0.088 \pm 0.014$ in the companion
full-stack diagnostic~\cite{Sayegh2026}, is consistent with this reading;
the comparison is an amplitude diagnostic rather than an independent
microscopic measurement.)  Integrating the
saturated uncertainty limit~\eqref{eq:power_limit} over the
$(1{+}1)$-dimensional Bose-Einstein spectrum at $T_{\rm dS}$---the
calculation performed explicitly by Pendry~\cite{Pendry1983}---yields
the exact one-dimensional Stefan-Boltzmann law:
\begin{equation}
\dot{E}_{\rm max}
  = g\,\frac{\pi}{12\hbar}\,(k_B T_{\rm dS})^2\,,
\label{eq:pendry_derived}
\end{equation}
where $g$ counts the independent polarisation degrees of freedom
of the edge mode.

Here, the representation theory of
Section~\ref{sec:interaction_channels} becomes decisive.  The
interaction channels operate in the adjoint
representation ($j = 1$ in SU(2) language) from the gauge-theoretic
requirement that horizon perturbations $\delta A \in
\mathfrak{su}(2)$ transform in the adjoint.  The internal
SU(2) label $j = 1$ should not be confused with a spacetime
spin; it counts the number of independent generator directions
in $\mathfrak{su}(2)$.  For the edge mode, each adjoint-sector
puncture carries $g = 2$ independent transverse polarisation
states---this follows from the dimension of the adjoint minus
the one longitudinal (gauge) direction, not from identifying
$j = 1$ with a spacetime vector boson.  This assignment was not
chosen to reproduce $\beta = 1/12$; it was derived independently
from the channel partition in Section~\ref{sec:two_channels}.
Substituting $g = 2$ into~\eqref{eq:pendry_derived} gives the
exact channel throughput $\dot{E} = \pi(k_B T_{\rm dS})^2/(6\hbar)$ that
normalises the Schwarzschild processing rate to yield
$\beta = 1/12$~\cite{Sayegh2026}.

Therefore, the interaction coupling $\beta = 1/12$ is not a
phenomenological input; it is the macroscopic expression of the
saturated quantum-gravitational uncertainty principle
$[\hat{A},\,\hat{\Sigma}] = i\hbar$ for $j = 1$ processing
channels on the apparent horizon.

\paragraph{The super-Schwarzschild window as a quantum speed-limit
violation.}

This result sharpens the interpretation of the
super-Schwarzschild window identified in
the companion cosmology paper~\cite{Sayegh2026}.  The aggregate interaction-processing capacity of
the horizon is $N_1 = I_{\rm eff}\,N$ active interaction
channels, each operating at the Pendry bound (saturated Heisenberg
uncertainty principle $\Delta E \cdot \Delta t = \hbar/2$).
The remaining $N - N_1 = \Gamma_{\rm emerge}\,N$ channels are
allocated to area growth (emergence) and do not contribute to
the dark-drag processing.  Every active channel is already
operating at the quantum speed limit; there is no reserve
processing capacity within the interaction sector.  The Schwarzschild rate
$c^3/2G$ is therefore not a phenomenological ceiling: it is the
geometric (Misner--Sharp) capacity of the boundary---derived in
the companion cosmology paper~\cite{Sayegh2026} as a processing bound from compactness
and causality alone---of which the Pendry-saturated interaction
sector delivers the fixed fraction $\beta = 1/12$, with no
adjustable headroom in either quantity.

During the super-Schwarzschild epoch ($\mathcal{S}(z) > 1$), the
vacuum drive exceeds the geometric capacity $c^3/2G$.  The processed dark
drag continues to operate in its own channel, while the normalized rejected
fraction is written into the protected memory described in
Section~\ref{sec:pb_write_instrument}.  It does not modify the homogeneous
expansion history or accumulate along photon trajectories.  The stored
record is read only when the local matter comparators of
Section~\ref{sec:portals_map} prepare a relative calibration state.  The
proposed Hubble offset is therefore an endpoint-holonomy signal of
throughput demand beyond the boundary capacity, not a bulk distance
perturbation.

\subsection{Holographic self-consistency and the late-time
attractor}
\label{sec:self_consistency}

The maximum boundary entropy production rate of the $N$
channels, each operating at the per-channel \emph{capacity}
$k_B T_{\rm dS}\,H$---one Gibbons--Hawking quantum per channel
per Hubble time (the companion cosmology paper~\cite{Sayegh2026})---is
$\sigma_{\rm max} = N\,H$ in units of $k_B$.  The balance below is
set by this capacity, not by the Pendry-limited interaction
throughput ($\beta = 1/12$ of it, which would instead give
$q = -23/24$).  The
Bekenstein--Hawking entropy growth rate is
$\dot{S}_{\rm BH} = 2(1+q)\,N\,H$.  These are equal if and only
if $q = -1/2$ ($\Omega_m = 1/3$).
the companion cosmology paper~\cite{Sayegh2026} shows that this self-consistency
condition selects the bound-enforcing completion: the
universe is driven to $\Omega_m = 1/3$ as a global attractor
(eigenvalue $-2/3$, basin of attraction $(0,\,1)$) because this
is the unique state in which the boundary can sustain maximum
entropy production while the horizon creates new
entropy-carrying area at the matching rate.  The coincidence
problem is thereby resolved by the same SU(2) boundary theory
that provides the screening function and coupling derived here.


% ======================================================================
%  SECTION 7: FALSIFIABILITY AND PREDICTIONS
% ======================================================================


\section{Falsifiability and Predictions}
\label{sec:falsifiability}

The microscopic and operational sectors admit independent tests.

\paragraph{Screening polynomial.}
A full LQG or spin-foam state sum in the linear-source universality class
should reproduce
$x(q)=(1-q)(2-q)/6$.  An order-unity Casimir-dependent displacement of the
saddle would falsify the structural-protection argument.

\paragraph{Transverse-current count.}
The source-reduced boundary theory predicts two physical Ohmic current
channels.  A microscopic calculation yielding one or more than two
unconstrained channels would change the Pendry coefficient and invalidate
$\beta=1/12$.

\paragraph{Common-gate native ratio.}
The strongest normalization-independent mapping prediction is
\begin{equation}
\Delta M_B-5\Delta a_B=0.
\label{eq:pb_native_null}
\end{equation}
Failure of this relation in a native likelihood rejects equal endpoint
weights or the common-gate topology even if the summed $H_0$ amplitude is
near the expected value.

\paragraph{Emission-side tilt.}
The equally charged source--observer atomic realization predicts the
redshift dependence in Table~\ref{tab:pb_post_emission}.  Its absence rules
out that endpoint class; its detection would select an emission-differential
model rather than the constant full-memory gate.

\paragraph{Loop nonclosure.}
The pre-registered bright-siren--SN experiment of
Sec.~\ref{sec:closure_experiment} predicts
$\mathcal H_{\rm GW-SN}=2X_H$.  At $X_H=0.0359$ this is $7.44\%$ or
$0.15591$ mag-equivalent.  A null rejects occurrence of the common gate on
the SN path, while a one-unit result requires an additional gauge-completed
endpoint sector.  The residual must be independent of host environment,
redshift, telescope, and analysis pipeline after standard corrections.

\paragraph{Neutral-method selection.}
The minimal gauge-neutral bilocal predicts $A=0$ for pure standard sirens,
cosmic chronometers, time-delay cosmography, and megamaser Hubble diagrams,
while the absolute electromagnetic SN ladder has $A=2$.  A robust isolated
$A=1$ class is not accommodated by retuning $X_H$; it would falsify the
minimal gate topology and require an additional endpoint representation.

\paragraph{Prior-blind ladder prediction.}
Paper A should freeze $p(H_{0,\rm phys},X_H|D_A)$ using non-ladder data.
Paper B then predicts the ladder posterior by Eq.~\eqref{eq:pb_pushforward}
without retuning the growth model.  This separates the test of the
operational recipient from the evidence for the processed dark-sector
interaction.



\section{Conclusions}
\label{sec:conclusion}

A single dynamical SU(2) boundary theory can organise the cosmological
apparent horizon into processed and over-capacity responses.  In the
processed sector, the dynamical Chern--Simons level and source-reduced state
sum give
\begin{equation}
I_{\rm eff}=\frac{1-q}{2},\qquad
x(q)=\frac{(1-q)(2-q)}6,\qquad
\beta=\frac1{12}.
\end{equation}
These quantities are fixed inputs to the companion cosmological growth
calculation.

In the excess sector, the hard rejected fraction defines the additive
capacity-cycle history $X_H$.  The source-reduced processed/excess quantum
instrument turns that history into a microscopic record: its controlled
translation gives $d\langle\chi\rangle/d\lambda_H=\Delta$, its full counting
statistics suppresses the memory width by $1/N$, and its large-$N$ covariant
saddle is the shift-symmetric three-form action.  Thus
$\langle\chi_0\rangle=X_H$ follows from the explicit write realization.

The read sector has now been completed within a definite local
matter--boundary EFT.  A Ramsey-readable atomic phase is compared with a
processed-sector boundary phase, and a Glauber radiative amplitude is
compared with a geometric radiometric reference.  The two records are
attached to the conjugate endpoints of the single gauge-neutral bilocal
$J^+W_\gamma J^-$.  The current Ward identity fixes one unit at each
endpoint, the normalized conditional matrix element removes the overall
portal overlap, and the uniqueness of the minimal gauge-neutral operator leaves only the common
translation.  Hence $g_T=g_R=g_{\rm com}=1$ in the minimal normalized
sector.

The global hierarchical ladder contains one serial preparation of the
shared $(a_B,M_B)$ calibration state.  Additional anchors, calibrator
supernovae, and Hubble-flow supernovae are parallel measurements and do not
reapply the defect.  Retarded passive orientation fixes
\begin{equation}
\Delta a_B=+\frac{X_H}{\ln10},\qquad
\Delta M_B=+\frac{5X_H}{\ln10},
\end{equation}
so that
\begin{equation}
\boxed{\Delta M_B=5\Delta a_B},\qquad
\boxed{\frac{H_{0,\rm ladder}}{H_{0,\rm phys}}=e^{2X_H}}.
\end{equation}
The reverse comparator gives the inverse map, and a common photometric
zero-point mode remains orthogonal to this response and cancels.

Calibration-loop holonomy supplies the convention-independent observable:
only path nonclosure survives arbitrary node-zero-point redefinitions.  The
paper closes the write normalization, the temporal and radiative portals,
common-gate universality, one engagement, endpoint sign, source-reduced
Hamiltonian dilation, counterterm basis, native likelihood, non-ladder
selection rules, and one concrete closure protocol within a minimal
microscopic comparator EFT.  The remaining frontier is to run the released
native data, realize and bound the portal experimentally, execute the loop,
and determine whether the full non-perturbative Standard Model plus CS/WZW
theory uniquely selects this universality class.  Failure of those empirical
or UV tests would reject this microscopic realization while leaving the
processed screening and growth calculation as a separately testable branch.


\section*{Data and Code Availability}
The native common-gate likelihood implementation, its dependency-free
primary-image FITS reader, documentation, and synthetic recovery test are
provided with this manuscript as
\texttt{shoes\_native\_common\_gate.py}.  It is designed for the official
Pantheon+/SH0ES $Y/L/C$ release~\cite{PantheonDataRelease}.  The repository
data themselves are not redistributed, and no placeholder covariance is
used as a scientific result.  The companion cosmology code and chains are
described in Ref.~\cite{Sayegh2026}.


\appendix


\section{Derivation of the dynamical Chern-Simons level}
\label{sec:dynamical_k_derivation}

The identification of the cosmological apparent horizon as
carrying an SU(2) Chern-Simons theory with a time-dependent level
$k(t)$ follows the variational structure of the ABCK
framework~\cite{Ashtekar1998,Ashtekar2000}, extended from isolated
(constant-area) horizons to the dynamical trapping
horizon~\cite{AshtekarKrishnan2003}.  The key dynamical boundary
condition (the slice-by-slice generalisation of the ABCK
horizon constraint) is motivated by spherical symmetry rather
than derived from a fully displayed first-principles variational
calculation; the complete derivation is in ABCK~\cite{Ashtekar1998}.

The starting point is the Holst action~\cite{Holst1996} for General Relativity on
a manifold $\mathcal{M}$ bounded by an inner worldtube $\Delta$
representing the apparent horizon.  To ensure a well-defined
variational principle, the variation of the action yields a bulk
term (which vanishes on-shell) and a boundary term evaluated on
$\Delta$.  The boundary symplectic potential $\Theta_\Delta$ is
read directly from this variation:
\begin{equation}
\Theta_\Delta
  = \frac{1}{8\pi G \gamma}
    \int_\Delta \Sigma_i \wedge \delta A^i\,,
\label{eq:symplectic_potential}
\end{equation}
where $A^i$ is the SU(2) Ashtekar--Barbero connection,
$\Sigma_i$ is the densitised triad 2-form, and $\gamma$ is the
Barbero--Immirzi parameter.

For an isolated black hole horizon, the area $A$ is strictly
constant, and the ABCK boundary condition imposes
$F^i = -(\pi\gamma/A)\,\Sigma^i$.  However, the cosmological
apparent horizon is a marginally trapped tube embedded in a highly
symmetric FLRW bulk.  Spherical symmetry dictates that the induced
connection and the triad remain proportional at any instant $t$
along the foliation, but the area
$A(t) = 4\pi c^2/H(t)^2$ is a dynamical function of cosmic time.
The instantaneous boundary condition, evaluated slice-by-slice, is
therefore
\begin{equation}
\Sigma^i(t) = -\frac{A(t)}{\pi\gamma}\,F^i(t)\,.
\label{eq:dynamical_bc}
\end{equation}

\noindent\textit{Background vs.\ fluctuations.}  The constraint
\eqref{eq:dynamical_bc} fixes the \emph{background} (classical)
relationship between $\Sigma^i$ and $A^i$ on each slice.
After imposing it, the independent boundary DOF is $A^i$ alone,
and the symplectic structure becomes the Chern-Simons form
$\Omega_\Delta \propto \int_{S^2} \delta A \wedge \delta A$ at the
end of this appendix.  The canonical commutation relation
$[\hat{A},\hat{\Sigma}] = i\hbar$ used in
Section~\ref{sec:pendry} refers to \emph{fluctuations}
$(\delta A, \delta\Sigma)$ around this background, which remain
independent because the constraint \eqref{eq:dynamical_bc} is
imposed on the mean field, not on each fluctuation mode
separately.  Both pictures are consistent: $A$ and $\Sigma$ are
proportional at the classical level; their quantum fluctuations
are conjugate.

Substituting into the symplectic potential gives
\begin{equation}
\Theta_\Delta
  = -\frac{A(t)}{8\pi^2 G \gamma^2}
    \int_\Delta F_i \wedge \delta A^i\,,
\end{equation}
where the integral is still over the 3D worldtube $\Delta$ and
$F_i \wedge \delta A^i$ is a 3-form (as required for a worldtube
integral).  Taking the exterior derivative in field space
($\Omega = \delta\Theta$) and restricting to a constant-$t$ spatial
slice $S^2_t \subset \Delta$ (a standard step in the ABCK
analysis~\cite{Ashtekar1998}) yields the symplectic 2-form on $S^2$.
Since $\delta F_i = d\,\delta A_i$ and $d\,\delta A_i \wedge \delta A^i$
reduces to $\delta A_i \wedge \delta A^i$ after integration by
parts on $S^2$, one obtains  Introducing $\ell_P^2 = \hbar G/c^3$, and using the
trace convention $\mathrm{tr}(T^a T^b) = \tfrac{1}{2}\delta^{ab}$
in the fundamental representation of $\mathfrak{su}(2)$ (consistent
with the Chern-Simons normalisation below):
\begin{equation}
\Omega_\Delta
  = \frac{A(t)}{8\pi^2 \gamma^2 \ell_P^2}
    \int_{S^2} \mathrm{tr}\bigl(
      \delta A \wedge \delta A\bigr)\,,
\label{eq:symplectic_2form}
\end{equation}
where the trace is in the fundamental representation of
$\mathfrak{su}(2)$ with the convention above.  Comparing with the
canonical symplectic 2-form for an SU(2) Chern-Simons theory at
level $k$ on a 2-manifold,
\begin{equation}
\Omega_{\rm CS}
  = \frac{k}{4\pi}
    \int_{S^2} \mathrm{tr}\bigl(
      \delta A \wedge \delta A\bigr)\,,
\end{equation}
written in the same trace convention, equating the prefactors
gives
\begin{equation}
k_{\rm schematic}(t) = \frac{A(t)}{2\pi\gamma^2\,\ell_P^2}\,.
\label{eq:k_schematic}
\end{equation}
This schematic matching captures the $A/\ell_P^2$ scaling but
does not yet account for the precise relation between the
component form $F_i \wedge \delta A^i$ appearing in
$\Theta_\Delta$ and the self-dual-projected trace form
$\mathrm{tr}(\delta A \wedge \delta A)$ of the Chern-Simons
theory.  The complete derivation, which properly tracks the
self-dual projector and the inner-product normalisation of the
boundary condition~\eqref{eq:dynamical_bc}, is given in
Ashtekar--Baez--Corichi--Krasnov~\cite{Ashtekar1998} and yields
\begin{equation}
\boxed{k(t) = \frac{A(t)}{4\pi\gamma\,\ell_P^2}\,.}
\label{eq:k_derived}
\end{equation}

The explicit time-dependence of $k(t)$ does not violate
Hamiltonian consistency because $k$ is a parameter of the
symplectic structure, not a dynamical variable of the boundary
theory itself.  A subtlety arises when comparing this continuous
evolution to the exact finite-level formulas (quantum-group
$\mathfrak{q}$-dimensions, framing phases) used at fixed $k$ in
Section~\ref{sec:path_integral_proof}: because
$k \sim A/(4\pi\gamma\ell_P^2) \sim 10^{122}/\gamma \gg 1$
throughout the cosmological era of interest, the rate of change
$dk/d\ln a = 2(1+q)k$ produces a fractional shift
$\delta k/k = 2(1+q)\,\delta(\ln a) \ll 1$ over any
observationally relevant time interval.  The exact-$k$ formulas
can therefore be evaluated at the instantaneous $k(t)$ without
adiabatic corrections; the continuous-$k$ and exact-level
descriptions are consistent to the same $\mathcal{O}(1/k)$
precision that bounds all other quantum corrections.  The apparent horizon is an \emph{open quantum
system} driven by the bulk FLRW expansion.  In the canonical
quantum theory, the continuous growth of $k(t)$ corresponds to
the continuous enlargement of the Chern-Simons Hilbert space
$\mathcal{H}_{k(t)}$, mirroring the physical creation of new
horizon punctures at the kinematic rate
$dk/d\ln a = 2(1+q)k$.  The Chern-Simons level $k$ and the
coarse-grained surface DOF count $N = A/(4\ell_P^2)$ are
proportional---$k = N/(\pi\gamma)$ from~\eqref{eq:k_derived}---but
not identical; $k$ is the level of the boundary gauge theory
while $N$ counts effective horizon DOF.  Throughout the main
text, $1/k$ and $1/N$ serve as equivalent control parameters
for the smallness of quantum corrections, up to the
$\mathcal{O}(1)$ factor $\pi\gamma$, since $1/k=\pi\gamma/N$.  This derivation motivates and extends to the dynamical setting the
treatment of $k$ as a dynamical clock used throughout the main
text (Section~\ref{sec:dynamical_k}).  The schematic formula
\eqref{eq:k_schematic} establishes the $A/\ell_P^2$ scaling
and the path to the exact result; the numerical prefactor in
the boxed formula~\eqref{eq:k_derived} (which differs from
$k_{\rm schematic}$ by a factor of $2\gamma$, reflecting the
self-dual projector and inner-product normalisation) is taken
from the full ABCK calculation~\cite{Ashtekar1998} rather than
derived here.


% ======================================================================
%  ACKNOWLEDGEMENTS
% ======================================================================
\begin{acknowledgments}
Numerical computations were performed using a modified version of
the \textsc{class} Boltzmann code~\cite{Blas2011} and the
\textsc{cobaya} MCMC sampler~\cite{Torrado2021}.
\end{acknowledgments}


% ======================================================================
%  REFERENCES
% ======================================================================


\section{Source Reduction, Transverse Currents, and Equal-Footing Normalisation}
\label{app:source_reduction}

This appendix records the source-reduction results that were previously
presented separately and are used in the main screening derivation.

\subsection{WZW spectral function and the two-channel count}

The SU(2)$_k$ current OPE is
\begin{equation}
J^a(z)J^b(w)\sim\frac{k\delta^{ab}}{(z-w)^2}
+\frac{i\epsilon^{ab}{}_{c}J^c(w)}{z-w}.
\label{eq:app_current_ope}
\end{equation}
The double pole gives an Ohmic retarded response,
\begin{equation}
G_R^{ab}(\omega)=-\frac{k}{2}\delta^{ab}i\omega,
\qquad
\rho^{ab}(\omega)=k\delta^{ab}\omega,
\label{eq:app_ohmic}
\end{equation}
so every physical current direction is a gapless absorber.  The saturated
thermal flux is the one-dimensional Pendry result
\begin{equation}
P_{\rm ch}^{\rm max}(T)
=\frac{g\pi(k_BT)^2}{12\hbar}.
\label{eq:app_pendry}
\end{equation}

The homogeneous junction source selects an internal axis $\hat n^a$,
\begin{equation}
J^a=J_\parallel\hat n^a+J_\perp^a,
\qquad \hat n_aJ_\perp^a=0.
\end{equation}
Variation with respect to $A_0^a$ imposes the boundary Gauss law
\begin{equation}
G^a=\frac{k}{2\pi}\epsilon^{AB}F^a_{AB}+J_0^a=0.
\label{eq:app_gauss}
\end{equation}
In the homogeneous sector it fixes the longitudinal
$U(1)_{\hat n}$ component.  The physical current fluctuations live in
$\mathfrak{su}(2)/\mathfrak{u}(1)_{\hat n}$, of dimension two.  Thus
\begin{equation}
\boxed{g_{\rm pol}=2.}
\end{equation}
Projection of Eq.~\eqref{eq:app_ohmic} gives
$\rho_\perp^{ab}=k\omega(\delta^{ab}-\hat n^a\hat n^b)$, whose physical
trace is $2k\omega$.

\subsection{Monopole selection and throughput weighting}

A homogeneous Israel--Darmois source contains only the $\ell=0$ boundary
monopole.  Expanding the sourced path integral in spherical harmonics,
nonzero modes integrate to a determinant independent of the channel label at
linear order.  The remaining zero-mode action has the schematic form
\begin{align}
S_{\rm eff}(f)
={}&N\left[f\ln f+(1-f)\ln(1-f)\right]-N\ln\eta
\notag\\
&-N\left[f\ln\lambda_1+(1-f)\ln\lambda_{1/2}\right].
\label{eq:app_zero_mode_action}
\end{align}
where $f=N_1/N$, $\lambda_1=1-q$, and
$\lambda_{1/2}=1+q$.  The first-order Chern--Simons kinetic term supplies
one universal prefactor $\eta\propto1/k$ for both populations.  Extremizing
gives
\begin{equation}
\frac{f}{1-f}=\frac{1-q}{1+q},
\qquad
f=\frac{1-q}{2}.
\label{eq:app_equal_footing}
\end{equation}
All dependence on the unknown overall source normalization cancels.  This is
the path-integral origin of equal footing.

The linear source is a rank-one adjoint insertion.  Since
$\mathrm{End}(V_j)=\oplus_{\ell=0}^{2j}V_\ell$ contains one copy of
$V_1$, the Wigner--Eckart theorem fixes the source operator, up to one
reduced matrix element, to $\hat n\cdot\mathbf J_j$.  Tracelessness of the
SU(2) generators removes residual $j$ dependence from the linear trace.
Emergence channels create area but carry no sourced monopole throughput;
they therefore enter the throughput-weighted mean as the trivial response,
whereas the adjoint channels carry one unit.

\subsection{Controlled corrections}

The Ohmic response gives an internal relaxation rate
$\Gamma_{\rm int}\sim kH$.  Since $|\dot q|/H\le9/8$, the first lag of the
instantaneous KMS allocation is bounded parametrically by
\begin{equation}
|\delta I_{\rm eff}|\lesssim\frac{9}{16k}.
\end{equation}
The first nonlinear source term is quadratic and introduces a Casimir
correction,
\begin{equation}
w_j=\eta\lambda_j\left[1+\epsilon_2C_2(j)+\cdots\right],
\qquad \epsilon_2=\mathcal O(1/k).
\end{equation}
For a present cosmological horizon, $1/k\sim\ell_P^2/r_A^2\sim10^{-122}$.
Thus the linear-source universality class is structurally protected at any
observable precision, while remaining falsifiable in a full microscopic
state-sum calculation.



\section{Entropy Bookkeeping and Exposure Provenance}
\label{app:exposure_provenance}

The flux split implies an exact entropy-equivalent identity.  With one
horizon temperature for the two parts,
\begin{equation}
d\Sigma_{\rm tot}=\frac{\delta Q_{\rm tot}}{T_H},
\qquad
d\Sigma_{\rm def}=\frac{\delta Q_{\rm exc}}{T_H},
\end{equation}
so
\begin{equation}
\frac{d\Sigma_{\rm def}}{d\Sigma_{\rm tot}}
=\frac{\delta Q_{\rm exc}}{\delta Q_{\rm tot}}
=\Delta(\mathcal S).
\label{eq:app_entropy_identity}
\end{equation}
Temperature and surface-gravity normalizations cancel because the ratio
compares two pieces of the same instantaneous throughput.

Three variables should not be conflated:
\begin{center}
\begin{tabular}{p{0.18\columnwidth}p{0.34\columnwidth}p{0.36\columnwidth}}
\toprule
variable & definition & role \\
\midrule
$X_H$ & $\int\Delta\,d\ln a$ & capacity-cycle history \\
$\chi_A$ & $\int\Delta\,d\ln r_A$ & area-growth history \\
$X_S$ & $\Sigma_{\rm def}/\Sigma_{\rm tot}$ & instantaneous normalized deficit \\
\bottomrule
\end{tabular}
\end{center}
The $X_S$ product rule contains dilution because its denominator grows.  The
history variables do not.  Physical erasure belongs in the carrier equation
through $\Gamma_{\rm er}$, not in the mathematical definition of an integral.

Finally, if a transfer factor is continuous, local in the accumulated
history, and multiplicative under path concatenation,
$C(X_1+X_2)=C(X_1)C(X_2)$ with $C(0)=1$, then the Cauchy functional equation
for $\ln C$ gives
\begin{equation}
C(X)=e^{s\eta X}.
\end{equation}
The exponential map is therefore fixed by additive exposure and
multiplicative transfer composition; the physical content lies in the
endpoint coefficient and orientation.


\begin{thebibliography}{99}

\bibitem{Sayegh2026}
B.~Sayegh,
``Horizon Thermodynamics as a Dark-Sector Regulator: A Derived Interaction
and Late-Time Growth Suppression,'' companion cosmology paper (Paper~A,
2026),
\href{https://doi.org/10.5281/zenodo.21157831}{doi:10.5281/zenodo.21157831}.

\bibitem{Ashtekar1998}
A.~Ashtekar, J.~Baez, A.~Corichi, and K.~Krasnov,
Phys.\ Rev.\ Lett.\ \textbf{80}, 904 (1998).

\bibitem{Ashtekar2000}
A.~Ashtekar, J.~C.~Baez, and K.~Krasnov,
Adv.\ Theor.\ Math.\ Phys.\ \textbf{4}, 1 (2000).

\bibitem{AshtekarKrishnan2003}
A.~Ashtekar and B.~Krishnan,
Phys.\ Rev.\ D \textbf{68}, 104030 (2003).

\bibitem{Engle2010}
J.~Engle, K.~Noui, and A.~Perez,
Phys.\ Rev.\ Lett.\ \textbf{105}, 031302 (2010).

\bibitem{Ghosh2011}
A.~Ghosh and A.~Perez,
Phys.\ Rev.\ Lett.\ \textbf{107}, 241301 (2011);
Erratum: \textbf{108}, 169901 (2012).

\bibitem{Hayward1998}
S.~A.~Hayward,
Class.\ Quant.\ Grav.\ \textbf{15}, 3147 (1998).

\bibitem{Kodama1980}
H.~Kodama,
Prog.\ Theor.\ Phys.\ \textbf{63}, 1217 (1980).

\bibitem{Padmanabhan2012}
T.~Padmanabhan,
Res.\ Astron.\ Astrophys.\ \textbf{12}, 891 (2012).

\bibitem{Onsager1931}
L.~Onsager,
Phys.\ Rev.\ \textbf{37}, 405 (1931).

\bibitem{Pendry1983}
J.~B.~Pendry,
J.\ Phys.\ A: Math.\ Gen.\ \textbf{16}, 2161 (1983).

\bibitem{Israel1966}
W.~Israel,
Nuovo Cimento B \textbf{44}, 1 (1966);
Erratum: \textbf{48}, 463 (1967).

\bibitem{Bekenstein1973}
J.~D.~Bekenstein,
Phys.\ Rev.\ D \textbf{7}, 2333 (1973).

\bibitem{Hawking1975}
S.~W.~Hawking,
Commun.\ Math.\ Phys.\ \textbf{43}, 199 (1975).

\bibitem{GibbonsHawking1977}
G.~W.~Gibbons and S.~W.~Hawking,
Phys.\ Rev.\ D \textbf{15}, 2738 (1977).

\bibitem{Jacobson1995}
T.~Jacobson,
Phys.\ Rev.\ Lett.\ \textbf{75}, 1260 (1995).

\bibitem{Hall2015}
B.~C.~Hall,
\textit{Lie Groups, Lie Algebras, and Representations},
2nd ed.\ (Springer, 2015).

\bibitem{Planck2020}
Planck Collaboration (N.~Aghanim et al.),
Astron.\ Astrophys.\ \textbf{641}, A6 (2020).

\bibitem{Riess2025}
A.~G.~Riess et al.,
Astrophys.\ J.\ Lett.\ \textbf{992}, L34 (2025).

\bibitem{DESY3_2022}
DES Collaboration (T.~M.~C.~Abbott et al.),
Phys.\ Rev.\ D \textbf{105}, 023520 (2022).

\bibitem{DESIDR2_2025}
DESI Collaboration (A.~G.~Adame et al.),
arXiv:2503.14738 (2025).

\bibitem{Blas2011}
D.~Blas, J.~Lesgourgues, and T.~Tram,
J.\ Cosmol.\ Astropart.\ Phys.\ \textbf{2011}(07), 034.

\bibitem{Torrado2021}
J.~Torrado and A.~Lewis,
J.\ Cosmol.\ Astropart.\ Phys.\ \textbf{2021}(05), 057.

\bibitem{EuclidForecast2020}
Euclid Collaboration (A.~Blanchard et al.),
Astron.\ Astrophys.\ \textbf{642}, A191 (2020).

\bibitem{LSST2019}
LSST Science Collaboration (P.~A.~Abell et al.),
arXiv:0912.0201 (2009).

\bibitem{CKN1999}
A.~G.~Cohen, D.~B.~Kaplan, and A.~E.~Nelson,
Phys.\ Rev.\ Lett.\ \textbf{82}, 4971 (1999).

\bibitem{Hsu2004}
S.~D.~H.~Hsu,
Phys.\ Lett.\ B \textbf{594}, 13 (2004).

\bibitem{Li2004}
M.~Li,
Phys.\ Lett.\ B \textbf{603}, 1 (2004).

\bibitem{Thomas2000}
S.~Thomas,
Phys.\ Rev.\ Lett.\ \textbf{89}, 081301 (2002).


\bibitem{Rovelli1996}
C.~Rovelli,
Phys.\ Rev.\ Lett.\ \textbf{77}, 3288 (1996).

\bibitem{KaulMajumdar2000}
R.~K.~Kaul and P.~Majumdar,
Phys.\ Rev.\ Lett.\ \textbf{84}, 5255 (2000).

\bibitem{Domagala2004}
M.~Domagala and J.~Lewandowski,
Class.\ Quantum Grav.\ \textbf{21}, 5233 (2004).

\bibitem{Meissner2004}
K.~A.~Meissner,
Class.\ Quantum Grav.\ \textbf{21}, 5245 (2004).

\bibitem{ENPP2010}
J.~Engle, K.~Noui, A.~Perez, and D.~Pranzetti,
Phys.\ Rev.\ D \textbf{82}, 044050 (2010).

\bibitem{Frodden2013}
E.~Frodden, A.~Ghosh, and A.~Perez,
Phys.\ Rev.\ D \textbf{87}, 121503(R) (2013).

\bibitem{Bianchi2012}
E.~Bianchi,
arXiv:1204.5122.

\bibitem{Holst1996}
S.~Holst,
Phys.\ Rev.\ D \textbf{53}, 5966 (1996).

\bibitem{Barbero1995}
J.~F.~Barbero~G.,
Phys.\ Rev.\ D \textbf{51}, 5507 (1995).

\bibitem{Immirzi1997}
G.~Immirzi,
Class.\ Quantum Grav.\ \textbf{14}, L177 (1997).

\bibitem{RovelliSmolin1995}
C.~Rovelli and L.~Smolin,
Nucl.\ Phys.\ B \textbf{442}, 593 (1995).

\bibitem{AshtekarLewandowski2004}
A.~Ashtekar and J.~Lewandowski,
Class.\ Quantum Grav.\ \textbf{21}, R53 (2004).

\bibitem{Witten1989}
E.~Witten,
Commun.\ Math.\ Phys.\ \textbf{121}, 351 (1989).

\bibitem{EMSS1989}
S.~Elitzur, G.~Moore, A.~Schwimmer, and N.~Seiberg,
Nucl.\ Phys.\ B \textbf{326}, 108 (1989).

\bibitem{Padmanabhan2010}
T.~Padmanabhan,
Mod.\ Phys.\ Lett.\ A \textbf{25}, 1129 (2010).

\bibitem{Verlinde2011}
E.~Verlinde,
J.\ High Energy Phys.\ \textbf{04} (2011) 029.

\bibitem{Wang2016}
B.~Wang, E.~Abdalla, F.~Atrio-Barandela, and D.~Pav\'on,
Rep.\ Prog.\ Phys.\ \textbf{79}, 096901 (2016).

\bibitem{Valiviita2008}
J.~Valiviita, E.~Majerotto, and R.~Maartens,
J.\ Cosmol.\ Astropart.\ Phys.\ \textbf{07} (2008) 020.

\bibitem{DiValentino2021}
E.~Di Valentino et al.,
Class.\ Quantum Grav.\ \textbf{38}, 153001 (2021).

\bibitem{DiValentino2020}
E.~Di Valentino, A.~Melchiorri, O.~Mena, and S.~Vagnozzi,
Phys.\ Rev.\ D \textbf{101}, 063502 (2020).

\bibitem{Yang2018}
W.~Yang, S.~Pan, E.~Di Valentino, R.~C.~Nunes, S.~Vagnozzi, and D.~F.~Mota,
J.\ Cosmol.\ Astropart.\ Phys.\ \textbf{09} (2018) 019.

\bibitem{WangWangLi2017}
S.~Wang, Y.~Wang, and M.~Li,
Phys.\ Rep.\ \textbf{696}, 1 (2017).




\bibitem{Ramsey1950}
N.~F.~Ramsey,
Phys.\ Rev.\ \textbf{78}, 695 (1950).

\bibitem{AbreuVisser2010}
G.~Abreu and M.~Visser,
Phys.\ Rev.\ D \textbf{82}, 044027 (2010).

\bibitem{Glauber1963a}
R.~J.~Glauber,
Phys.\ Rev.\ \textbf{130}, 2529 (1963).

\bibitem{Glauber1963b}
R.~J.~Glauber,
Phys.\ Rev.\ \textbf{131}, 2766 (1963).

\bibitem{HolevoWerner2001}
A.~S.~Holevo and R.~F.~Werner,
Phys.\ Rev.\ A \textbf{63}, 032312 (2001).

\bibitem{Caves1982}
C.~M.~Caves,
Phys.\ Rev.\ D \textbf{26}, 1817 (1982).

\bibitem{Riess2022}
A.~G.~Riess et al.,
Astrophys.\ J.\ Lett.\ \textbf{934}, L7 (2022).

\bibitem{Brout2022}
D.~Brout et al.,
Astrophys.\ J.\ \textbf{938}, 110 (2022).

\bibitem{BrownYork1993}
J.~D.~Brown and J.~W.~York,
Phys.\ Rev.\ D \textbf{47}, 1407 (1993).

\bibitem{CaiKim2005}
R.-G.~Cai and S.~P.~Kim,
JHEP \textbf{02}, 050 (2005).

\bibitem{MisnerSharp1964}
C.~W.~Misner and D.~H.~Sharp,
Phys.\ Rev.\ \textbf{136}, B571 (1964).

\bibitem{DiFrancesco1997}
P.~Di~Francesco, P.~Mathieu, and D.~S\'en\'echal,
\emph{Conformal Field Theory} (Springer, New York, 1997).

\bibitem{Kubo1957}
R.~Kubo,
J.\ Phys.\ Soc.\ Jpn.\ \textbf{12}, 570 (1957).

\bibitem{Esposito2009}
M.~Esposito, U.~Harbola, and S.~Mukamel,
Rev.\ Mod.\ Phys.\ \textbf{81}, 1665 (2009).

\bibitem{Garrahan2010}
J.~P.~Garrahan and I.~Lesanovsky,
Phys.\ Rev.\ Lett.\ \textbf{104}, 160601 (2010).

\bibitem{Gawedzki2004}
K.~Gawedzki, I.~Todorov, and P.~Tran-Ngoc-Bich,
Commun.\ Math.\ Phys.\ \textbf{247}, 217 (2004).

\bibitem{BachasGaberdiel2004}
C.~Bachas and M.~R.~Gaberdiel,
JHEP \textbf{11}, 065 (2004).



\bibitem{Lindblad1976}
G.~Lindblad,
Commun. Math. Phys. \textbf{48}, 119 (1976).

\bibitem{BIPMSI2026}
Bureau International des Poids et Mesures,
\emph{The International System of Units (SI), 9th ed., version 3.01}
(2025), \url{https://www.bipm.org/en/publications/si-brochure}.

\bibitem{NISTPOWR}
National Institute of Standards and Technology,
``Primary Optical Watt Radiometer (POWR),''
\url{https://www.nist.gov/programs-projects/primary-optical-watt-radiometer-powr}.

\bibitem{Gentile1996}
T.~R.~Gentile, J.~M.~Houston, and C.~L.~Cromer,
Appl. Opt. \textbf{35}, 1056 (1996).

\bibitem{PantheonDataRelease}
Pantheon+SH0ES Collaboration,
``Pantheon+SH0ES Data Release,''
\url{https://github.com/PantheonPlusSH0ES/DataRelease}.

\bibitem{Schutz1986}
B.~F.~Schutz,
Nature \textbf{323}, 310 (1986).

\bibitem{Abbott2017}
B.~P.~Abbott et al. (LIGO Scientific and Virgo Collaborations),
Nature \textbf{551}, 85 (2017).

\bibitem{JimenezLoeb2002}
R.~Jimenez and A.~Loeb,
Astrophys. J. \textbf{573}, 37 (2002).

\bibitem{Moresco2012}
M.~Moresco et al.,
J. Cosmol. Astropart. Phys. \textbf{08} (2012) 006.

\bibitem{Birrer2024}
S.~Birrer et al.,
Space Sci. Rev. \textbf{220}, 48 (2024).

\bibitem{Herrnstein1999}
J.~R.~Herrnstein et al.,
Nature \textbf{400}, 539 (1999).

\bibitem{Pesce2020}
D.~W.~Pesce et al.,
Astrophys. J. Lett. \textbf{891}, L1 (2020).

\end{thebibliography}


\end{document}
